\documentclass[10pt,a4paper,oneside]{report}

\usepackage[T1]{fontenc}
\usepackage{lmodern}
\usepackage[a4paper,margin=2.25cm,headheight=15pt]{geometry}
\usepackage{microtype}
\usepackage{setspace}
\usepackage{parskip}
\usepackage{enumitem}
\usepackage{array}
\usepackage{booktabs}
\usepackage{longtable}
\usepackage{tabularx}
\usepackage{ragged2e}
\usepackage[table]{xcolor}
\usepackage{fvextra}
\usepackage{pdflscape}
\usepackage{titlesec}
\usepackage{fancyhdr}

\definecolor{ReportBlue}{HTML}{1F4E79}
\definecolor{ReportNeutral}{HTML}{475569}
\definecolor{ReportRule}{HTML}{CBD5E1}
\definecolor{ReportSoft}{HTML}{F6F8FA}
\definecolor{ReportTableStripe}{HTML}{F4F7FB}
\definecolor{ReportCode}{HTML}{F7F7F4}
\usepackage[colorlinks=true,linkcolor=ReportBlue,urlcolor=ReportBlue,citecolor=ReportBlue]{hyperref}
\arrayrulecolor{ReportRule}

\setcounter{secnumdepth}{-1}
\setcounter{tocdepth}{2}
\setstretch{1.12}
\setlist[itemize,1]{leftmargin=1.45em,itemsep=0.18\baselineskip,topsep=0.35\baselineskip,label=\textcolor{ReportBlue}{\small\textbullet}}
\setlist[itemize,2]{leftmargin=1.8em,itemsep=0.08\baselineskip,topsep=0.15\baselineskip,label=\textcolor{ReportNeutral}{\scriptsize$\triangleright$}}
\setlist[enumerate]{leftmargin=1.75em,itemsep=0.18\baselineskip,topsep=0.35\baselineskip}
\renewcommand{\arraystretch}{1.18}
\emergencystretch=3em
\sloppy

\title{CRA - Additive Manufacturing Threat Taxonomy}
\author{%
Sami Gawad\\
University of Southern Denmark\\
Denmark\\
\texttt{sagaw22@student.sdu.dk}
\and
Kenny Nguyen\\
University of Southern Denmark\\
Denmark\\
\texttt{kengu22@student.sdu.dk}
\and
Danieal Johnbaskar\\
University of Southern Denmark\\
Denmark\\
\texttt{dajoh22@student.sdu.dk}
\and
Marcel Omid Rasmussen Ghasemi\\
University of Southern Denmark\\
Denmark\\
\texttt{magha21@student.sdu.dk}
}
\date{\today}

\titleformat{\chapter}
  {\normalfont\huge\bfseries\color{ReportBlue}}
  {}
  {0pt}
  {}
  [\vspace{0.25\baselineskip}{\color{ReportRule}\titlerule[0.6pt]}]
\titlespacing*{\chapter}{0pt}{-0.5\baselineskip}{1.1\baselineskip}

\titleformat{\section}
  {\normalfont\Large\bfseries\color{ReportBlue}}
  {}
  {0pt}
  {}
\titlespacing*{\section}{0pt}{1.1\baselineskip}{0.35\baselineskip}

\titleformat{\subsection}
  {\normalfont\large\bfseries\color{ReportNeutral}}
  {}
  {0pt}
  {}
\titlespacing*{\subsection}{0pt}{0.9\baselineskip}{0.25\baselineskip}

\titleformat{\subsubsection}
  {\normalfont\normalsize\bfseries\color{ReportBlue}}
  {}
  {0pt}
  {}
\titlespacing*{\subsubsection}{0pt}{0.75\baselineskip}{0.15\baselineskip}

\pagestyle{fancy}
\fancyhf{}
\fancyhead[L]{\small CRA - Additive Manufacturing Threat Taxonomy}
\fancyhead[R]{\small\thepage}
\renewcommand{\headrulewidth}{0.35pt}
\fancypagestyle{plain}{%
  \fancyhf{}%
  \fancyfoot[C]{\small\thepage}%
  \renewcommand{\headrulewidth}{0pt}%
}

\fvset{
  fontsize=\small,
  breaklines=true,
  breakanywhere=true,
  frame=single,
  framesep=3mm,
  rulecolor=\color{ReportRule}
}

\begin{document}
\pagenumbering{roman}

\begin{titlepage}
\thispagestyle{empty}
\begin{center}
\vspace*{1.5cm}
\vspace{1.2cm}
{\Huge\bfseries\color{ReportBlue} CRA - Additive\\[0.25em]Manufacturing Threat Taxonomy\par}
\vspace{1.2cm}
\vfill
\begin{tabularx}{0.92\linewidth}{@{}>{\centering\arraybackslash}X>{\centering\arraybackslash}X@{}}
\textbf{Sami Gawad} & \textbf{Kenny Nguyen}\\
University of Southern Denmark & University of Southern Denmark\\
Denmark & Denmark\\
\texttt{sagaw22@student.sdu.dk} & \texttt{kengu22@student.sdu.dk}\\[1.4cm]
\textbf{Danieal Johnbaskar} & \textbf{Marcel Omid Rasmussen Ghasemi}\\
University of Southern Denmark & University of Southern Denmark\\
Denmark & Denmark\\
\texttt{dajoh22@student.sdu.dk} & \texttt{magha21@student.sdu.dk}
\end{tabularx}
\vfill
{\large \today\par}
\vspace*{1.2cm}
\end{center}
\end{titlepage}

\begin{abstract}
This report defines a merged threat taxonomy for additive manufacturing by combining two source taxonomies into a single goal-target-method structure. The taxonomy preserves additive-manufacturing-specific attack detail while extending the scope to supply-chain assets, production dependencies, equipment compromise, part compromise, and contextual disruption. Each taxonomy node is labelled using the CIA triad where relevant, and the method-level entries are connected to MITRE ATT\&CK mitigation families. The report also includes a Cyber Resilience Act crosswalk to show where taxonomy nodes correspond to product-security, vulnerability-handling, documentation, and incident-reporting requirements.
\end{abstract}

\tableofcontents
\clearpage
\pagenumbering{arabic}



\chapter{Purpose}


This document defines a \textbf{merged threat taxonomy for additive manufacturing (AM)} by combining the two source taxonomies below into one coherent structure:

\begin{enumerate}
\item \textbf{Yampolskiy et al. (2018)}, \emph{Security of additive manufacturing: Attack taxonomy and survey}
\item \textbf{Kumar et al. (2025)}, \emph{Security of cyber-physical Additive Manufacturing supply chain: Survey, attack taxonomy and solutions}
\end{enumerate}

The CRA/ATT\&CK mitigation-alignment rationale is additionally informed by:

\begin{enumerate}
\setcounter{enumi}{2}
\item \textbf{Ruohonen et al. (2025)}, \emph{An Alignment Between the CRA's Essential Requirements and the ATT\&CK's Mitigations}
\end{enumerate}

The merged taxonomy is designed to be:

\begin{itemize}
\item \textbf{supply-chain-wide} rather than limited to a single AM stage
\item \textbf{three-dimensional} rather than purely target/method based
\item \textbf{careful about duplicates}, merging only where the concepts are materially the same
\item \textbf{explicit about scope differences}, preserving useful distinctions from both papers
\end{itemize}


\chapter{Source-model summary}



\section{Yampolskiy et al. (2018)}


Yampolskiy structures AM threats into \textbf{two major threat categories}:

\begin{itemize}
\item \textbf{Theft of technical data}
\item \textbf{AM sabotage}
\end{itemize}

For each category, the taxonomy uses \textbf{two top-level branches}:

\begin{itemize}
\item \textbf{Attack targets}
\item \textbf{Attack methods}
\end{itemize}

This taxonomy is especially strong on:

\begin{itemize}
\item detailed AM-specific sabotage logic
\item distinction between \textbf{technical data theft} and \textbf{sabotage}
\item detailed treatment of \textbf{geometry}, \textbf{required properties}, \textbf{process parameters}, \textbf{post-processing}, \textbf{equipment compromise}, and \textbf{environmental impact}
\end{itemize}


\section{Kumar et al. (2025)}


Kumar structures threats across the \textbf{entire AM supply chain} using \textbf{three top-level dimensions}:

\begin{itemize}
\item \textbf{Attack goals}
\item \textbf{Attack targets}
\item \textbf{Attack methods}
\end{itemize}

This taxonomy is especially strong on:

\begin{itemize}
\item supply-chain-wide scope
\item broader target coverage beyond printing itself
\item explicit inclusion of \textbf{unauthorised production}, \textbf{restricted/export-controlled parts}, \textbf{overproduction}, \textbf{distribution/logistics}, and \textbf{broader business/supply-chain disruption}
\end{itemize}


\section{Ruohonen et al. (2025)}


Ruohonen et al. evaluate the alignment between the Cyber Resilience Act's essential cybersecurity requirements and MITRE ATT\&CK mitigations.

In this document, Ruohonen et al. is used as a \textbf{supporting CRA/ATT\&CK alignment source}. It is not used as an AM threat-taxonomy source, and it does not replace MITRE ATT\&CK as the source for mitigation names, IDs, or catalog descriptions.

This source is especially useful for:

\begin{itemize}
\item supporting the conceptual link between CRA essential requirements and ATT\&CK mitigation categories
\item identifying where ATT\&CK mitigations align well or poorly with CRA requirements
\item reinforcing that ATT\&CK mitigations can support CRA-oriented security analysis at a conceptual level, while not proving product-specific compliance
\end{itemize}


\chapter{Merge decision}


The merged taxonomy uses \textbf{Kumar's high-level shape} as the umbrella structure:

\begin{enumerate}
\item \textbf{Attack goals} = why the attacker is acting
\item \textbf{Attack targets} = what in the AM ecosystem is being attacked
\item \textbf{Attack methods} = how the attacker achieves the attack
\end{enumerate}

Within that umbrella, \textbf{Yampolskiy's detailed AM-specific categories} are used to deepen and refine the target and method branches.

Ruohonen et al. is not merged into the AM taxonomy structure. It is used only to support the CRA/ATT\&CK alignment rationale behind the mitigation and regulatory crosswalks.

This means:

\begin{itemize}
\item \textbf{Yampolskiy's threat families} are mapped into \textbf{Kumar's attack goals}
\item \textbf{Yampolskiy's targets and methods} are preserved where they add AM-specific detail
\item \textbf{duplicates are normalized}, but only when they are semantically equivalent
\item \textbf{non-equivalent concepts are kept separate}
\end{itemize}


\chapter{CIA Triad Labeling Convention}


Each attack goal, attack target, and attack method below is labeled with one or more CIA triad tags:

\begin{itemize}
\item \textbf{\texttt{[C]} Confidentiality:} unauthorized access to, disclosure of, or reconstruction of AM-related information, designs, process data, or protected knowledge
\item \textbf{\texttt{[I]} Integrity:} unauthorized modification, corruption, substitution, or falsification of AM data, processes, equipment, or produced parts
\item \textbf{\texttt{[A]} Availability:} disruption, delay, degradation, denial, or loss of access to AM systems, equipment, production stages, or supply-chain operations
\end{itemize}

A node may carry multiple labels (for example, \texttt{[C, I]}) where it directly affects more than one security property. Labels are applied based on the direct security role of the taxonomy entry, not on worst-case speculation.


\chapter{MITRE ATT\&CK Mitigation Mapping Convention}


The mitigation mappings below are a \textbf{derived analytical crosswalk} from the AM attack-method taxonomy to MITRE ATT\&CK mitigations. They are not presented as MITRE-authored AM-specific mappings, and they should not be read as a claim that a listed mitigation fully eliminates the attack.

The mappings aim to be exhaustive at the level of \textbf{materially relevant ATT\&CK mitigations} for each AM attack method. They intentionally do not repeat every ATT\&CK mitigation under every method where the connection would be merely generic or weak.

The mapping uses MITRE ATT\&CK \textbf{Enterprise} and \textbf{Industrial Control Systems (ICS)} mitigations as the primary catalogs because AM environments combine enterprise IT, engineering workstations, production networks, cyber-physical equipment, and industrial safety constraints. The MITRE pages checked for this version identify ATT\&CK v19 as current. MITRE ATT\&CK \textbf{Mobile} mitigations are not repeated under every method; they should be added only when the AM product or workflow includes a mobile app, mobile device management surface, or mobile-controlled production interface.

The source hierarchy for mitigation entries is: MITRE ATT\&CK provides the mitigation names, IDs, and catalog descriptions; Ruohonen et al. (2025) supports the broader rationale for aligning ATT\&CK mitigations with CRA essential requirements; this document derives the AM-specific method-level mapping.

The CRA connection is specific rather than generic. Article 13 makes the Annex I cybersecurity requirements operative for products with digital elements, while Annex I Part I defines product-security properties and Annex I Part II defines vulnerability-handling requirements. In this document, ATT\&CK mitigations are used as a traceability layer between concrete AM attack methods and CRA-relevant control areas:

\begin{itemize}
\item theft of CAD files, process data, recipes, or machine-resident know-how maps most directly to \texttt{Annex I.I(2)(d)} on unauthorised access, \texttt{Annex I.I(2)(e)} on confidentiality, and \texttt{Annex I.I(2)(l)} on monitoring, with ATT\&CK mitigations such as access management, authentication, encryption, data loss prevention, and audit logging
\item CAD, STL, AMF, 3MF, tool-path, G-code, recipe, sensor, or configuration tampering maps most directly to \texttt{Annex I.I(2)(f)} on integrity of data, commands, programs, and configuration, supported by \texttt{Annex I.I(2)(j)} on attack-surface limitation and \texttt{Annex I.I(2)(k)} on exploitation mitigation, with ATT\&CK mitigations such as code signing, communication authenticity, software process and device authentication, input validation, boot integrity, and execution prevention
\item compromise of AM equipment, firmware, slicers, cloud design platforms, remote services, update channels, or third-party dependencies maps to \texttt{Annex I.I(2)(a)-(c)}, \texttt{Annex I.I(2)(j)-(l)}, and \texttt{Annex I.II(1)-(8)}, with ATT\&CK mitigations such as supply chain management, vulnerability scanning, update software, secure update paths, boot integrity, audit logging, data backup, and network segmentation
\end{itemize}

Each method-level mitigation block follows this logic:

\begin{itemize}
\item \textbf{Primary mitigations} are the ATT\&CK mitigations most directly relevant to preventing or constraining the method.
\item \textbf{Supporting mitigations} are useful hardening, recovery, monitoring, or governance measures that reduce preconditions, blast radius, or recovery time.
\item \textbf{AM-specific application} explains how the generic ATT\&CK mitigation concept maps onto additive-manufacturing assets such as CAD files, build jobs, slicers, process recipes, feedstock, post-processing equipment, quality data, and machine-control networks.
\item \textbf{Residual coverage notes} are used where ATT\&CK does not provide a complete preventive answer, especially for physical tampering, metrology-based reverse engineering, process-material science issues, natural disruptions, and legal or commercial abuse.
\end{itemize}

Mitigation entries are written in a name-first format for readability:

\begin{itemize}
\item Read each mitigation entry as: \textbf{mitigation name}, followed by the relevant \textbf{MITRE ATT\&CK catalog} and \textbf{mitigation ID} in parentheses.
\item \texttt{Data Loss Prevention (Enterprise M1057; ICS M0803)} means the same mitigation concept appears in both the Enterprise and Industrial Control Systems (ICS) ATT\&CK catalogs under different IDs.
\item \texttt{Remote Data Storage (Enterprise M1029)} means the mitigation is drawn from only the Enterprise catalog in this mapping.
\item The sub-bullet below each mitigation entry gives a short, source-based explanation of what that MITRE mitigation means.
\end{itemize}


\chapter{CRA Mapping Convention}


The Cyber Resilience Act (CRA) mapping below is an \textbf{analytical crosswalk} from this AM threat taxonomy to Regulation (EU) 2024/2847. It is not legal advice and it is not a claim that a product is compliant or non-compliant.

The mapping uses \textbf{Annex I} as the primary reference because Annex I contains the essential cybersecurity requirements that can be meaningfully linked to threat categories. Articles are used as \textbf{legal anchors} that make those requirements operative:

\begin{itemize}
\item \textbf{Article 2} and \textbf{Article 3} define the scope and key terms, including products with digital elements.
\item \textbf{Article 6} states the market requirement that products with digital elements must meet Annex I Part I and that manufacturer processes must meet Annex I Part II.
\item \textbf{Article 13} sets the manufacturer's design, development, production, risk-assessment, component due-diligence, vulnerability-handling, support-period, documentation, user-information, and corrective-action obligations.
\item \textbf{Article 14} adds mandatory reporting for actively exploited vulnerabilities and severe incidents affecting the security of products with digital elements.
\item \textbf{Article 31} and \textbf{Annex VII} define technical-documentation expectations.
\item \textbf{Article 32} and \textbf{Annex VIII} define conformity-assessment procedures.
\item \textbf{Annex II} defines user information and instructions, including secure use, support, updates, foreseeable risks, and vulnerability-reporting contact information.
\end{itemize}

The CRA applicability gate is important: the Regulation applies to \textbf{products with digital elements} whose intended purpose or reasonably foreseeable use includes a direct or indirect logical or physical data connection to a device or network. Therefore, an AM threat maps strongly to the CRA only where it affects a connected AM product, software, firmware, component, remote data processing solution, vulnerability-handling process, security update, or related manufacturer process. Purely physical, contractual, IP-law, export-control, or logistics threats are marked as weak, conditional, or outside the direct CRA cybersecurity requirement set.

Mapping strength labels:

\begin{itemize}
\item \textbf{Direct:} the threat clearly corresponds to an Annex I cybersecurity requirement where the affected AM asset is a product with digital elements or part of one.
\item \textbf{Conditional:} the mapping depends on the AM implementation, for example whether the asset is digitally connected, software-controlled, firmware-based, remotely processed, or sold as a component.
\item \textbf{Partial/gap:} the CRA addresses only a supporting cybersecurity aspect of the threat, not the threat as a whole.
\item \textbf{Mostly outside CRA:} the threat is mainly physical, legal, commercial, export-control, or operational, and only becomes CRA-relevant if it is enabled through a product with digital elements.
\end{itemize}

Reference notation used in the CRA crosswalk:

\begin{itemize}
\item \texttt{Annex I.I(1)} means Annex I, Part I, point (1): risk-based appropriate cybersecurity.
\item \texttt{Annex I.I(2)(a-m)} means Annex I, Part I, point (2), subpoints (a) to (m), covering product security properties.
\item \texttt{Annex I.II(1-8)} means Annex I, Part II, points (1) to (8), covering vulnerability handling.
\item \texttt{Annex II} means information and instructions to the user.
\item \texttt{Annex VII} means technical documentation.
\end{itemize}


\chapter{Canonical merged taxonomy}



\chapter{1. Attack goals}


Attack goals describe the attacker's intended outcome.


\section{1.1 Intellectual property compromise / technical data theft [C]}


The attacker seeks to obtain, reconstruct, copy, or misuse protected AM-related knowledge.

Includes:

\begin{itemize}
\item theft of technical data
\begin{itemize}
\item Covers unauthorised acquisition of AM technical representations or process knowledge that can enable reproduction, modification, or competitive use.
\end{itemize}
\item intellectual property theft
\begin{itemize}
\item Targets AM IP such as unique designs, manufacturing techniques, product innovations, trade secrets, or other protected know-how.
\end{itemize}
\item CAD file theft
\begin{itemize}
\item Involves unauthorised access to proprietary CAD models or design files that define the product blueprint.
\end{itemize}
\item design theft
\begin{itemize}
\item Covers acquisition of protected design information that can be used to duplicate, modify, or understand an AM object.
\end{itemize}
\item process-knowledge theft
\begin{itemize}
\item Targets manufacturing procedures, printing settings, material specifications, or production techniques that reveal how the part is made.
\end{itemize}
\item trade secret theft
\begin{itemize}
\item Covers confidential AM knowledge such as material compositions, post-processing methods, or unique manufacturing processes.
\end{itemize}
\item counterfeiting and unlicensed reproduction
\begin{itemize}
\item Uses protected or copied AM knowledge to produce fake products or replicas without the required rights or permissions.
\end{itemize}
\item reverse-engineered reproduction
\begin{itemize}
\item Recreates AM designs or manufacturing information from physical objects, scans, side channels, or equipment analysis rather than direct file access.
\end{itemize}
\item industrial espionage
\begin{itemize}
\item Involves infiltrating or compromising AM organisations to obtain proprietary information for competitive or financial advantage.
\end{itemize}
\item unauthorised reproduction based on stolen or reconstructed knowledge
\begin{itemize}
\item Produces AM parts from stolen or reverse-engineered knowledge without authorisation from the legitimate rights holder.
\end{itemize}
\end{itemize}


\section{1.2 Unauthorised production / counterfeit / restricted-part production [I]}


The attacker seeks to produce parts that are not authorised by the legitimate owner, regulator, or internal control regime.

Includes:

\begin{itemize}
\item counterfeit production
\begin{itemize}
\item Produces parts that imitate genuine components and may deceive buyers or fail to meet expected quality and safety standards.
\end{itemize}
\item unlicensed reproduction
\begin{itemize}
\item Replicates proprietary AM parts or designs without the rights, licences, or permissions of the legitimate owner.
\end{itemize}
\item grey-market production
\begin{itemize}
\item Creates AM parts outside authorised production and distribution channels, often to supply unofficial markets.
\end{itemize}
\item grey-market operations
\begin{itemize}
\item Runs unauthorised AM production or sales activity without consent from original producers or authorised distributors.
\end{itemize}
\item unauthorised overproduction
\begin{itemize}
\item Produces more parts than the legitimate order, approval, or production authorisation allows.
\end{itemize}
\item production of export-controlled parts
\begin{itemize}
\item Uses AM to produce parts or technologies whose manufacture or transfer is restricted by export-control rules.
\end{itemize}
\item production of internally restricted parts
\begin{itemize}
\item Produces parts that an organisation or facility restricts because of safety, security, legal, regulatory, or policy concerns.
\end{itemize}
\item evasion of supply-chain controls
\begin{itemize}
\item Uses decentralised or localised AM production to avoid permitted manufacturing routes and make source tracing harder.
\end{itemize}
\item bypassing regulatory or contractual constraints
\begin{itemize}
\item Produces parts outside required licences, certifications, contractual limits, or other authorisation controls.
\end{itemize}
\item bypassing regulatory compliance
\begin{itemize}
\item Manufactures AM components without required compliance steps, especially in regulated sectors where safety and legal risks follow.
\end{itemize}
\end{itemize}


\section{1.3 Data breach and data manipulation [C, I]}


The attacker seeks unauthorised access to AM-related data or seeks to alter that data to influence production, inspection, quality, or downstream use.

Includes:

\begin{itemize}
\item data infiltration
\begin{itemize}
\item Breaks into AM systems to access valuable data such as CAD files, STL models, slicing settings, G-Code, or confidential design knowledge.
\end{itemize}
\item data exfiltration
\begin{itemize}
\item Removes AM data from compromised systems, potentially exposing sensitive information over time.
\end{itemize}
\item design and process information access
\begin{itemize}
\item Accesses design parameters, material specifications, or printing settings to learn about quality, performance, or distinctive product features.
\end{itemize}
\item manipulation of design files
\begin{itemize}
\item Changes AM design files to influence the geometry, integrity, safety, or functionality of the printed object.
\end{itemize}
\item manipulation of process data
\begin{itemize}
\item Alters AM process-related data such as printing settings, slicing parameters, material information, or production instructions.
\end{itemize}
\item manipulation of quality-control data
\begin{itemize}
\item Changes inspection, testing, or acceptance data so defective or manipulated parts may appear compliant.
\end{itemize}
\item manipulation of sensor readings
\begin{itemize}
\item Falsifies process or environmental readings used for monitoring, control, quality assurance, or safety decisions.
\end{itemize}
\item manipulation of machine instructions
\begin{itemize}
\item Alters G-Code, tool paths, or equivalent machine instructions that determine how the printer executes the build.
\end{itemize}
\item tampering with AM-related data or control representations
\begin{itemize}
\item Broadly covers unauthorised changes to AM digital artefacts or control data that influence production or downstream use.
\end{itemize}
\end{itemize}


\section{1.4 Sabotage / integrity degradation [I, A]}


The attacker seeks to reduce the integrity, reliability, safety, or functional fitness of the AM process or the resulting object.

Includes:

\begin{itemize}
\item reduction of mechanical strength
\begin{itemize}
\item Degrades the part below required mechanical tolerance so it may fail under normal operational conditions.
\end{itemize}
\item reduction of fatigue life
\begin{itemize}
\item Shortens the part's useful lifetime so failure can occur prematurely during repeated or long-term use.
\end{itemize}
\item introduction of hidden defects
\begin{itemize}
\item Inserts flaws such as internal weaknesses, voids, or defects that may remain undetected during quality control.
\end{itemize}
\item introducing weaknesses
\begin{itemize}
\item Creates vulnerabilities in the part or process that reduce durability, reliability, or expected performance.
\end{itemize}
\item material/property degradation
\begin{itemize}
\item Alters material characteristics or part properties so the final object no longer meets its required function.
\end{itemize}
\item tampering with material properties
\begin{itemize}
\item Manipulates material composition, process parameters, or related conditions to change the properties of the produced part.
\end{itemize}
\item contamination
\begin{itemize}
\item Introduces unwanted material, foreign objects, or NBC contamination into feedstock, equipment, or parts.
\end{itemize}
\item compromising part integrity
\begin{itemize}
\item Undermines the final part's fit, form, function, structure, or conformance to its authorised specification.
\end{itemize}
\item equipment damage
\begin{itemize}
\item Harms AM mechanical, electrical, cyber, or safety-critical equipment, causing wear, malfunction, or irreparable damage.
\end{itemize}
\item environmental damage
\begin{itemize}
\item Causes harm beyond the part itself, including effects on nearby equipment, facility systems, customers, or deployed environments.
\end{itemize}
\item disrupting production processes
\begin{itemize}
\item Interferes with printers, software systems, network connectivity, or production steps to cause delays, waste, or degraded output.
\end{itemize}
\item targeting critical components
\begin{itemize}
\item Focuses sabotage on safety-critical or mission-critical parts where failure can cause system failures or accidents.
\end{itemize}
\item concealing defects
\begin{itemize}
\item Makes compromised parts appear visually or procedurally acceptable despite hidden weaknesses or structural problems.
\end{itemize}
\item quality degradation that leads to unsafe or substandard output
\begin{itemize}
\item Produces parts whose integrity, functionality, or safety is reduced below expected quality or performance standards.
\end{itemize}
\end{itemize}


\section{1.5 Availability disruption / denial of service [A]}


The attacker seeks to interrupt, delay, disable, overload, or render unavailable AM operations or supply-chain flows.

Includes:

\begin{itemize}
\item denial of service
\begin{itemize}
\item Disrupts service availability so AM systems, equipment, production stages, or supply-chain operations cannot function as needed.
\end{itemize}
\item ransomware-driven lockout
\begin{itemize}
\item Locks AM control systems or related infrastructure, delaying production and creating operational or financial losses.
\end{itemize}
\item production stoppage
\begin{itemize}
\item Halts or interrupts AM production, causing downtime, delayed output, material waste, or recovery costs.
\end{itemize}
\item logistics disruption
\begin{itemize}
\item Disrupts transportation, routing, delivery, or distribution flows that AM supply-chain continuity depends on.
\end{itemize}
\item equipment unavailability
\begin{itemize}
\item Makes AM printers, post-processing equipment, control systems, or supporting devices unavailable for production.
\end{itemize}
\item power-related disruption
\begin{itemize}
\item Interrupts or destabilises power needed for AM temperature control, environmental stability, equipment operation, or safe completion of builds.
\end{itemize}
\item natural disasters
\begin{itemize}
\item Represents exploited disaster conditions such as floods, earthquakes, hurricanes, or wildfires that can damage AM facilities or infrastructure.
\end{itemize}
\item power outages
\begin{itemize}
\item Interrupt printing and supporting environmental conditions, potentially causing incomplete builds, material waste, or equipment damage.
\end{itemize}
\item material contamination
\begin{itemize}
\item Covers disruption conditions where AM materials absorb moisture, particulates, or other contaminants that affect material properties or part quality.
\end{itemize}
\item temperature and humidity fluctuations
\begin{itemize}
\item Affect AM material behaviour and process consistency, especially where powders or thermal conditions are sensitive.
\end{itemize}
\item disruption of supply-chain continuity
\begin{itemize}
\item Breaks the dependable flow of materials, data, equipment, parts, or services needed to sustain AM operations.
\end{itemize}
\end{itemize}


\chapter{2. Attack targets}


Attack targets describe the assets, stages, artefacts, systems, environments, or protected objects that can be attacked.

To make the merged taxonomy usable, the targets are grouped by \textbf{AM lifecycle zone}.


\section{2.1 Design and specification assets [C, I]}


Targets in this family are the digital and specification artefacts that define the intended 3D object, including geometry, material characteristics, required properties, design representations, and machine-ready instructions.


\subsection{2.1.1 Part specification [C, I]}


The part specification captures the geometry, properties, appearance, material choices, and interface details needed to define what the AM part should be and how it should perform.

Merged from Yampolskiy's detailed theft/sabotage treatment and Kumar's broader design/specification treatment.

Includes:

\begin{itemize}
\item \textbf{Part geometry}
\begin{itemize}
\item external shape
\item internal cavities
\item dimensions
\item tolerances
\item surface finish requirements
\item mass / center-of-mass relevant geometry
\end{itemize}
\item \textbf{Required properties}
\begin{itemize}
\item mechanical properties
\item thermal properties
\item chemical properties
\item manufacturing properties
\item electrical or other functional properties where relevant
\end{itemize}
\item \textbf{Colour / texture / appearance specifications}
\item \textbf{Material specification}
\item \textbf{Assembly/interface features}
\item \textbf{Feature addition/removal opportunities}
\end{itemize}


\subsection{2.1.2 CAD model integrity [C, I]}


The CAD model is the native design representation of the 3D object and may contain shape, dimensions, features, design intent, revisions, constraints, and proprietary product knowledge.

Includes:

\begin{itemize}
\item CAD files
\item model revisions
\item embedded design intent
\item metadata and constraints inside CAD tooling
\end{itemize}


\subsection{2.1.3 STL / AMF / 3MF / intermediate model integrity [C, I]}


The intermediate model is the converted downstream representation of the design, such as STL surface geometry or richer AMF/3MF object data, used before slicing and machine execution.

Includes:

\begin{itemize}
\item STL integrity
\item mesh integrity
\item converted geometry integrity
\item format-specific object representation
\end{itemize}


\subsection{2.1.4 Tool-path / slicing output / machine-instruction integrity [C, I]}


Tool paths, slicing outputs, and machine instructions are the executable manufacturing representations that translate the design into printer movements, layer behavior, process commands, and command ordering.

Includes:

\begin{itemize}
\item slicing parameters
\item generated tool-paths
\item G-code or equivalent control representations
\item proprietary machine instructions
\item timing/order of manufacturing commands
\end{itemize}


\section{2.2 Process knowledge and production recipe assets [C, I, A]}


Targets in this family are the process recipes, live production states, post-processing instructions, sensor data, and dependent production knowledge that determine how AM material is formed, treated, inspected, and accepted.


\subsection{2.2.1 Manufacturing process specification [C, I]}


The manufacturing process specification defines the AM process category and parameters, such as power, speed, hatch spacing, build direction, supports, atmosphere, feedstock settings, and machine-specific know-how that influence microstructure and properties.

A major Yampolskiy target preserved as a first-class node.

Includes:

\begin{itemize}
\item scan strategy
\item power / speed / hatch spacing
\item layer parameters
\item build direction
\item support strategy
\item atmosphere/environment settings
\item feedstock-related process settings
\item machine-specific process know-how
\item \textbf{Power Bed Fusion (PBF)}
\begin{itemize}
\item build parameters
\item support structures
\item build direction
\item feedstock
\item environment
\item build temperature
\end{itemize}
\item \textbf{Direct Energy Deposition (DED)}
\begin{itemize}
\item environment
\end{itemize}
\item additional AM process categories where the source taxonomy used open-ended ellipsis nodes
\end{itemize}


\subsection{2.2.2 Manufacturing process execution [I, A]}


Manufacturing process execution is the live build-stage operation where control commands, sequencing, parameter changes, machine state, material handling, and system access affect the printed output.

Includes:

\begin{itemize}
\item process control commands
\item execution sequencing
\item live parameter control
\item real-time machine behavior
\item machine state transitions
\end{itemize}


\subsection{2.2.3 Post-processing specification and execution [C, I, A]}


Post-processing covers the required post-build treatments and their execution, including heat treatment, HIPing, finish machining, surface treatment, inspection, and related settings that affect material properties, fit, surface quality, and acceptance.

Merged from Yampolskiy's \textbf{post-processing specification} and Kumar's \textbf{post-processing stage}.

Includes:

\begin{itemize}
\item heat-treatment specification
\item hot isostatic pressing specification
\item finish machining specification
\item coating/surface-treatment specification
\item inspection/test configuration used during post-processing
\item live execution of post-processing steps
\item \textbf{Heat treatment}
\begin{itemize}
\item HIPing
\item annealing
\end{itemize}
\item \textbf{Finish machining}
\begin{itemize}
\item EDM
\item surface finish
\end{itemize}
\item surface contamination
\item structural modifications
\item introduction of hidden components
\item substitution of post-processing materials
\item data alteration during post-processing
\item sabotage of post-processing equipment
\end{itemize}


\subsection{2.2.4 Sensor information and quality-control data [C, I, A]}


Sensor information and quality-control data are monitoring, inspection, and feedback artefacts used to estimate process state, guide control loops, verify conformance, or decide whether a part should be accepted.

Pulled from Yampolskiy's process-control discussion and retained as a practical target class.

Includes:

\begin{itemize}
\item in-situ monitoring data
\item IR thermography data
\item environmental sensor data
\item quality-control measurements
\item inspection/test results
\item closed-loop control feedback
\end{itemize}


\subsection{2.2.5 Indirect manufacturing dependencies [C, I, A]}


Indirect manufacturing dependencies are AM-produced or third-party assets, such as tooling, molds, supporting infrastructure, suppliers, software, or power dependencies, whose compromise can affect the final manufactured outcome.

Merged from both papers.

Includes:

\begin{itemize}
\item molds, tools, jigs, and intermediate AM-produced manufacturing aids
\item external dependencies whose compromise affects the final AM outcome
\item dependent processes needed to produce the eventual final part
\item tooling
\item molds
\item third-party supplier dependencies
\item cyber or power infrastructure dependencies
\item insider and social-engineering exposure in dependent processes
\end{itemize}


\section{2.3 Supply-chain assets and dependencies [C, I, A]}


Targets in this family are the materials, parts, equipment, software, services, logistics flows, and hand-off points that connect AM production to upstream suppliers and downstream distribution.


\subsection{2.3.1 Feedstock / source materials [I, A]}


Feedstock and source materials are the powders, filaments, resins, wires, recycled streams, and material characteristics whose chemistry, contamination, size distribution, moisture, or oxygen exposure can affect part quality and process safety.

A key Yampolskiy supply-chain target retained explicitly.

Includes:

\begin{itemize}
\item powders
\item filaments
\item resins
\item wire feedstock
\item recycled/reblended material streams
\item material chemistry, size distribution, moisture, oxygen exposure, contaminants
\item chemical composition
\item powder characteristics
\item recycling
\item sieving
\item blending properties
\item NBC contamination
\end{itemize}


\subsection{2.3.2 Replacement parts and upgrades [I, A]}


Replacement parts and upgrades are maintenance, repair, software, firmware, mechanical, electrical, and electronic components introduced after deployment that can affect machine control, process behavior, or equipment reliability.

Includes:

\begin{itemize}
\item replacement mechanical parts
\item electrical parts
\item firmware updates
\item software updates
\item maintenance components
\item third-party upgrades
\item software / firmware replacement parts
\item mechanical replacement parts
\item electrical replacement parts
\item electronic replacement parts
\end{itemize}


\subsection{2.3.3 AM equipment delivered through the supply chain [C, I, A]}


This target covers AM machines and related equipment before, during, or after delivery, where printers, post-processing devices, or quality-control devices may arrive from suppliers or intermediaries in a compromised state.

Includes:

\begin{itemize}
\item 3D printers
\item post-processing equipment
\item quality-control devices
\item compromised machines delivered from suppliers or intermediaries
\end{itemize}


\subsection{2.3.4 Third-party software, firmware, and service dependencies [C, I, A]}


Third-party software, firmware, and services are external digital dependencies in the AM tool chain, including slicers, machine-control software, process modeling, quality software, cloud collaboration, and remote maintenance links.

Expanded from Kumar's supply-chain scope and Yampolskiy's cyber infrastructure notion.

Includes:

\begin{itemize}
\item OEM firmware
\item slicers
\item machine-control software
\item process-modeling software
\item quality software
\item cloud platforms
\item collaboration systems
\item networked maintenance connections
\item remote support links
\end{itemize}


\subsection{2.3.5 Logistics and distribution chain [C, I, A]}


The logistics and distribution chain covers movement, storage, tracking, packaging, and hand-off of design data, feedstock, spare parts, equipment, and finished parts across AM supply-chain actors.

Explicitly retained from Kumar.

Includes:

\begin{itemize}
\item movement of design data
\item movement of feedstock
\item movement of spare parts
\item movement of finished parts
\item packaging, storage, and hand-off points
\item shipment tracking systems
\end{itemize}


\subsection{2.3.6 Manufactured objects in transit or return flows [I, A]}


Manufactured objects in transit or return flows are finished or returned parts that may be replaced, contaminated, tampered with, or otherwise compromised during shipping, warehousing, repair, recycling, or return handling.

Merged from Yampolskiy's "manufactured object" supply-chain target and Kumar's downstream delivery/return context.

Includes:

\begin{itemize}
\item parts in shipping
\item parts in warehousing
\item parts returned for repair/recycling
\item counterfeit replacement during delivery/return
\end{itemize}


\section{2.4 Physical production system assets [C, I, A]}


Targets in this family are the physical AM outputs, equipment, production surroundings, and downstream environments that can suffer confidentiality, integrity, availability, safety, or functional consequences.


\subsection{2.4.1 Final manufactured part [C, I]}


The final manufactured part is the central sabotage target because AM creates the material and the part together, so fit, form, properties, strength, fatigue life, microstructure, and contamination can all affect function and safety.

Includes:

\begin{itemize}
\item \textbf{Fit and form}
\begin{itemize}
\item shape and dimensions
\item mass
\item center of mass
\end{itemize}
\item \textbf{Function}
\begin{itemize}
\item physical properties
\item chemical properties
\item thermal properties
\item electrical properties
\item magnetic properties
\item optical properties
\item mechanical properties
\end{itemize}
\item \textbf{Mechanical strength}
\item \textbf{Fatigue life}
\item \textbf{Microstructure}
\item \textbf{Hidden flaws / voids / microfractures}
\item \textbf{Material composition and internal integrity}
\item \textbf{Contamination / embedded objects / hidden components}
\end{itemize}


\subsection{2.4.2 AM equipment [C, I, A]}


AM equipment includes printers, post-processing devices, quality-control devices, mechanical and electrical components, cyber infrastructure, networks, safety subsystems, and machine elements whose compromise can disrupt production or alter output.

Merged from both papers, keeping Yampolskiy's detailed split.

Includes:

\begin{itemize}
\item \textbf{Mechanical components}
\begin{itemize}
\item increased wear
\item irreparable damage
\end{itemize}
\item \textbf{Electrical components}
\item \textbf{Cyber infrastructure}
\begin{itemize}
\item software
\item firmware
\item hardware
\item networks
\item availability
\item integrity
\item timing
\end{itemize}
\item \textbf{Safety-critical subsystems}
\item \textbf{Inspection / post-processing machines}
\item explosion / fire
\end{itemize}


\subsection{2.4.3 AM production environment [I, A]}


The AM production environment is the physical surrounding in which printing occurs, including temperature, humidity, contamination, ventilation, vibration, power stability, physical access, and emergency conditions that affect process quality and availability.

Merged from Yampolskiy's environment focus and Kumar's production-environment target.

Includes:

\begin{itemize}
\item temperature
\item humidity
\item contaminants/dust
\item ventilation/fumes
\item vibration/noise
\item power-supply stability
\item physical access security
\item nearby equipment and surrounding facility conditions
\item emergency preparedness
\end{itemize}


\subsection{2.4.4 Downstream environment affected by deployed part [I, A]}


The downstream environment is the customer, host system, nearby equipment, facility, or mission context that can be harmed when a compromised AM machine or deployed part causes physical, cyber-physical, functional, or stealth-related effects.

A Yampolskiy-style environmental target retained because it captures downstream physical harm.

Includes:

\begin{itemize}
\item host system using the sabotaged part
\item nearby equipment harmed by machine malfunction
\item customer environment exposed to contaminated or compromised parts
\item mission/system environment impacted by stealth-signature or functional degradation
\item machine shop tools
\begin{itemize}
\item disruption
\item physical damage
\item interlocks
\item safety violation
\end{itemize}
\item electrical components in the surrounding environment
\begin{itemize}
\item overload
\item damage / burn
\end{itemize}
\item device properties
\begin{itemize}
\item efficiency
\item dependability
\item reliability
\item life time
\item damage / failure
\end{itemize}
\item CPS stealth properties
\begin{itemize}
\item emanations
\item reflectivity
\end{itemize}
\end{itemize}


\section{2.5 Rights, governance, and restricted-object targets [C, I]}


Targets in this family are protected rights, production limits, and restricted part categories where AM can enable infringement, unauthorised production, unauthorised transfer, or policy violations.


\subsection{2.5.1 Intellectual property rights and legally protected knowledge [C]}


This target covers copyrighted designs, patented products or processes, trade secrets, and proprietary AM design or manufacturing knowledge that can be infringed, stolen, pirated, counterfeited, or reproduced without authorisation.

Merged from Kumar's infringement/IP-right target and Yampolskiy's technical-data theft focus.

Includes:

\begin{itemize}
\item copyrighted designs
\item patented processes
\item trade secrets
\item proprietary process knowledge
\item confidential design/manufacturing know-how
\item intellectual property theft
\item counterfeiting
\item unauthorised reproduction
\item design piracy
\item patent infringement
\item brand dilution
\end{itemize}


\subsection{2.5.2 Unauthorised overproduction [I, A]}


Unauthorised overproduction targets the production quantity and authorisation boundary, where AM parts are produced beyond legitimate orders or approved runs, causing resource, inventory, legal, value, and reputational consequences.

A Kumar-specific target retained as distinct rather than merged away.

Includes:

\begin{itemize}
\item excess production beyond legitimate order quantity
\item hidden parallel production
\item off-books manufacturing runs
\item resource wastage
\item inventory surplus
\item loss of value
\item opportunity costs
\item legal implications
\item reputation damage
\end{itemize}


\subsection{2.5.3 Export-controlled parts [C, I]}


Export-controlled parts are AM-relevant components, technologies, or products whose production, transfer, or distribution is legally restricted, often because of defence, military, or dual-use significance.

Includes:

\begin{itemize}
\item unauthorised transfer
\item reverse engineering
\item intellectual property theft
\item distribution networks
\item dual-use technologies
\item supply-chain vulnerabilities
\end{itemize}


\subsection{2.5.4 Internally restricted parts [C, I]}


Internally restricted parts are components or objects restricted inside an organisation or facility because of safety, security, legal, contractual, regulatory, customer, or national-security handling concerns.

Includes:

\begin{itemize}
\item unauthorised production
\item sabotage
\item insider threats
\item regulatory non-compliance
\item intellectual property concerns
\item brand reputation
\end{itemize}


\chapter{3. Attack methods and mitigation mappings}


Attack methods describe the means by which attackers compromise the targets above.

The merged taxonomy groups methods into families, while preserving important source distinctions.


\section{3.1 Acquisition and exfiltration methods [C, I]}


Methods in this family obtain AM technical data either by direct theft or by reconstructing design, process, post-processing, or equipment knowledge from objects, communications, emissions, scans, or analysis.


\subsection{3.1.1 Theft / intellectual property theft [C]}


Direct acquisition of AM technical data, including model files, process files, post-processing know-how, and trade secrets, through hacking, malware, compromised equipment, physical access, insiders, or outsourcers.

Merged from Yampolskiy's \textbf{Theft} and Kumar's \textbf{Intellectual property theft}.

Includes:

\begin{itemize}
\item theft of design files
\item theft of process files
\item theft of post-processing know-how
\item theft of trade secrets
\item theft via malware, insider access, network intrusion, or compromised infrastructure
\item hacking
\item source file worm
\item compromised equipment
\item insider attack
\item malicious outsourcer
\item additional theft vectors where the source taxonomy used open-ended ellipsis nodes
\end{itemize}

\textbf{MITRE ATT\&CK mitigation mapping (derived):}

\begin{itemize}
\item \textbf{Primary mitigations:}
\begin{itemize}
\item Data Loss Prevention (\texttt{Enterprise M1057}; \texttt{ICS M0803})
\begin{itemize}
\item Identifies, monitors, and controls sensitive-data movement to reduce unauthorized access, transmission, or exfiltration.
\end{itemize}
\item Encrypt Sensitive Information (\texttt{Enterprise M1041}; \texttt{ICS M0941})
\begin{itemize}
\item Uses strong encryption to protect sensitive information against unauthorized access or tampering.
\end{itemize}
\item Encrypt Network Traffic (\texttt{ICS M0808})
\begin{itemize}
\item Uses strong cryptographic protocols to reduce eavesdropping on network communications.
\end{itemize}
\item Operational Information Confidentiality (\texttt{ICS M0809})
\begin{itemize}
\item Protects operational-process, facility, configuration, program, or database information that could reveal sensitive know-how or IP.
\end{itemize}
\item Access Management (\texttt{ICS M0801})
\begin{itemize}
\item Enforces access decisions, often through gateways or authentication services, before users can reach systems or devices.
\end{itemize}
\item Authorization Enforcement (\texttt{ICS M0800})
\begin{itemize}
\item Restricts read, manipulation, or execution privileges to authenticated users who require access under approved policies.
\end{itemize}
\item Human User Authentication (\texttt{ICS M0804})
\begin{itemize}
\item Requires users to authenticate before accessing data or sending commands to a device or system.
\end{itemize}
\item Multi-factor Authentication (\texttt{Enterprise M1032}; \texttt{ICS M0932})
\begin{itemize}
\item Requires two or more forms of identity evidence before granting access to a system.
\end{itemize}
\item Privileged Account Management (\texttt{Enterprise M1026}; \texttt{ICS M0926})
\begin{itemize}
\item Controls creation, use, permissions, and accountability for administrative, root, SYSTEM, or other privileged accounts.
\end{itemize}
\item User Account Management (\texttt{Enterprise M1018}; \texttt{ICS M0918})
\begin{itemize}
\item Manages account creation, modification, use, and permissions to reduce unauthorized access and privilege misuse.
\end{itemize}
\item Account Use Policies (\texttt{Enterprise M1036}; \texttt{ICS M0936})
\begin{itemize}
\item Configures account-use rules such as lockouts, login times, and inactivity timeouts to reduce unauthorized access.
\end{itemize}
\end{itemize}
\end{itemize}

\begin{itemize}
\item \textbf{Supporting mitigations:}
\begin{itemize}
\item Network Segmentation (\texttt{Enterprise M1030}; \texttt{ICS M0930})
\begin{itemize}
\item Divides networks into isolated segments or zones to restrict access, reduce lateral movement, and protect critical systems.
\end{itemize}
\item Limit Access to Resource Over Network (\texttt{Enterprise M1035}; \texttt{ICS M0935})
\begin{itemize}
\item Restricts network access to resources such as file shares, remote systems, and services to legitimate users or systems.
\end{itemize}
\item Filter Network Traffic (\texttt{Enterprise M1037}; \texttt{ICS M0937})
\begin{itemize}
\item Filters ingress, egress, or lateral traffic, including protocol-aware rules, to restrict malicious or unnecessary communication.
\end{itemize}
\item Audit (\texttt{Enterprise M1047}; \texttt{ICS M0947})
\begin{itemize}
\item Records and reviews activity, permissions, configurations, and integrity indicators to identify weaknesses or anomalies.
\end{itemize}
\item Remote Data Storage (\texttt{Enterprise M1029})
\begin{itemize}
\item Moves critical logs or sensitive files to secure off-host storage to reduce tampering, destruction, or unauthorized access.
\end{itemize}
\item User Training (\texttt{Enterprise M1017}; \texttt{ICS M0917})
\begin{itemize}
\item Trains users to recognize and respond to phishing, social engineering, and other user-interaction threats.
\end{itemize}
\item Threat Intelligence Program (\texttt{Enterprise M1019}; \texttt{ICS M0919})
\begin{itemize}
\item Collects and analyzes threat intelligence to track trends, guide defensive priorities, and reduce risk.
\end{itemize}
\item Vulnerability Scanning (\texttt{Enterprise M1016}; \texttt{ICS M0916})
\begin{itemize}
\item Assesses systems, applications, and networks for vulnerabilities, misconfigurations, or unpatched software to prioritize remediation.
\end{itemize}
\item Update Software (\texttt{Enterprise M1051}; \texttt{ICS M0951})
\begin{itemize}
\item Applies vendor patches and upgrades to reduce exploitation of known software or firmware vulnerabilities.
\end{itemize}
\item Exploit Protection (\texttt{Enterprise M1050}; \texttt{ICS M0950})
\begin{itemize}
\item Detects, blocks, or mitigates conditions that indicate or enable software exploitation.
\end{itemize}
\end{itemize}
\end{itemize}

\begin{itemize}
\item \textbf{AM-specific application:} protect CAD models, STL/3MF files, build jobs, process recipes, post-processing recipes, quality records, and machine-resident know-how with least-privilege access, encryption, egress controls, segmented engineering networks, audited access paths, and off-host logs.
\end{itemize}

\begin{itemize}
\item \textbf{Residual coverage note:} ATT\&CK mitigations reduce cyber-enabled theft paths, but insider collusion, lawful access misuse, and knowledge leakage through legitimate supplier collaboration still require contractual, organizational, and need-to-know controls outside ATT\&CK's mitigation catalog.
\end{itemize}


\subsection{3.1.2 Reverse engineering [C, I]}


Reconstruction of AM design or manufacturing information without direct file theft, including observation during manufacturing, side-channel collection, scanning of finished objects, and equipment analysis that exposes process or post-processing knowledge.

Merged from both papers.

Includes:

\begin{itemize}
\item reverse engineering of AM design
\item reverse engineering of finished parts
\item reproduction based on reconstructed digital or physical knowledge
\item eavesdropping
\item side-channel analysis
\item scanning
\end{itemize}

Submethods:

\begin{itemize}
\item \textbf{Eavesdropping and side-channel analysis}
\begin{itemize}
\item network eavesdropping
\item acoustic / sound side channels
\item thermal side channels
\item electro-magnetic side channels
\item power-consumption side channels
\item magnetic/IR side channels
\item observation of machine emissions
\end{itemize}
\item \textbf{Scanning and metrology-based reconstruction}
\begin{itemize}
\item 3D scanning / CMS / CMM-based measurement
\item CT / X-ray / radiographic reconstruction
\end{itemize}
\item \textbf{AM equipment analysis}
\begin{itemize}
\item cyber analysis of software
\item cyber analysis of firmware
\item cyber analysis of file formats
\item cyber analysis of network protocols
\item physical analysis of mechanical elements
\item physical analysis of electrical elements
\item physical analysis of heat sources
\item physical analysis of electronics
\item physical analysis of environmental characteristics
\item additional physical-analysis dimensions where the source taxonomy used open-ended ellipsis nodes
\end{itemize}
\end{itemize}

\textbf{MITRE ATT\&CK mitigation mapping (derived):}

\begin{itemize}
\item \textbf{Primary mitigations:}
\begin{itemize}
\item Operational Information Confidentiality (\texttt{ICS M0809})
\begin{itemize}
\item Protects operational-process, facility, configuration, program, or database information that could reveal sensitive know-how or IP.
\end{itemize}
\item Encrypt Sensitive Information (\texttt{Enterprise M1041}; \texttt{ICS M0941})
\begin{itemize}
\item Uses strong encryption to protect sensitive information against unauthorized access or tampering.
\end{itemize}
\item Encrypt Network Traffic (\texttt{ICS M0808})
\begin{itemize}
\item Uses strong cryptographic protocols to reduce eavesdropping on network communications.
\end{itemize}
\item Data Loss Prevention (\texttt{Enterprise M1057}; \texttt{ICS M0803})
\begin{itemize}
\item Identifies, monitors, and controls sensitive-data movement to reduce unauthorized access, transmission, or exfiltration.
\end{itemize}
\item Network Segmentation (\texttt{Enterprise M1030}; \texttt{ICS M0930})
\begin{itemize}
\item Divides networks into isolated segments or zones to restrict access, reduce lateral movement, and protect critical systems.
\end{itemize}
\item Limit Access to Resource Over Network (\texttt{Enterprise M1035}; \texttt{ICS M0935})
\begin{itemize}
\item Restricts network access to resources such as file shares, remote systems, and services to legitimate users or systems.
\end{itemize}
\item Filter Network Traffic (\texttt{Enterprise M1037}; \texttt{ICS M0937})
\begin{itemize}
\item Filters ingress, egress, or lateral traffic, including protocol-aware rules, to restrict malicious or unnecessary communication.
\end{itemize}
\item Minimize Wireless Signal Propagation (\texttt{ICS M0806})
\begin{itemize}
\item Detects, understands, and reduces unnecessary RF propagation to limit wireless monitoring or unauthorized access opportunities.
\end{itemize}
\end{itemize}
\end{itemize}

\begin{itemize}
\item \textbf{Supporting mitigations:}
\begin{itemize}
\item Access Management (\texttt{ICS M0801})
\begin{itemize}
\item Enforces access decisions, often through gateways or authentication services, before users can reach systems or devices.
\end{itemize}
\item Authorization Enforcement (\texttt{ICS M0800})
\begin{itemize}
\item Restricts read, manipulation, or execution privileges to authenticated users who require access under approved policies.
\end{itemize}
\item Human User Authentication (\texttt{ICS M0804})
\begin{itemize}
\item Requires users to authenticate before accessing data or sending commands to a device or system.
\end{itemize}
\item Multi-factor Authentication (\texttt{Enterprise M1032}; \texttt{ICS M0932})
\begin{itemize}
\item Requires two or more forms of identity evidence before granting access to a system.
\end{itemize}
\item Software Process and Device Authentication (\texttt{ICS M0813})
\begin{itemize}
\item Requires devices and software processes to authenticate, especially for remote connections and API access.
\end{itemize}
\item Audit (\texttt{Enterprise M1047}; \texttt{ICS M0947})
\begin{itemize}
\item Records and reviews activity, permissions, configurations, and integrity indicators to identify weaknesses or anomalies.
\end{itemize}
\item Application Developer Guidance (\texttt{Enterprise M1013}; \texttt{ICS M0913})
\begin{itemize}
\item Provides secure-development guidance to reduce weaknesses in applications, systems, and APIs.
\end{itemize}
\item Limit Hardware Installation (\texttt{Enterprise M1034}; \texttt{ICS M0934})
\begin{itemize}
\item Blocks users or groups from installing or using unapproved hardware, including removable devices.
\end{itemize}
\end{itemize}
\end{itemize}

\begin{itemize}
\item \textbf{AM-specific application:} reduce exposure of design files, machine telemetry, process recipes, network traffic, acoustic/RF emissions where feasible, and firmware or protocol details that make reconstruction of part geometry or process knowledge easier.
\end{itemize}

\begin{itemize}
\item \textbf{Residual coverage note:} ATT\&CK only partially covers physical reverse engineering of a legitimately obtained part through 3D scanning, CT, CMM, destructive analysis, or metallurgical analysis. Those cases require IP strategy, contractual restrictions, watermarking/provenance, export controls, or physical product-design choices beyond ATT\&CK.
\end{itemize}


\section{3.2 Protection-circumvention and unauthorised reproduction methods [C, I]}


Methods in this family defeat controls intended to restrict AM access, use, reproduction, attribution, or production authorisation, enabling protected designs or parts to be copied, altered, or produced without permission.


\subsection{3.2.1 Bypassing digital and physical protection / protection removal [C, I]}


Removal or bypass of safeguards around AM files, equipment, firmware, access controls, watermarks, signatures, device binding, encryption, or copy protection so protected IP or production functions can be accessed or altered.

Merged from Kumar's \textbf{Bypassing digital and physical protection} and Yampolskiy's \textbf{Protection removal}.

Includes:

\begin{itemize}
\item access-control bypass
\item digital-signature tampering/removal
\item copy-protection defeat
\item encryption-key extraction
\item DRM-like protection bypass
\item firmware modification to bypass safeguards
\item watermark removal
\item device-binding removal
\item additional protection-removal techniques where the source taxonomy used open-ended ellipsis nodes
\end{itemize}

\textbf{MITRE ATT\&CK mitigation mapping (derived):}

\begin{itemize}
\item \textbf{Primary mitigations:}
\begin{itemize}
\item Code Signing (\texttt{Enterprise M1045}; \texttt{ICS M0945})
\begin{itemize}
\item Verifies digital signatures so only trusted, untampered code is allowed to execute.
\end{itemize}
\item Boot Integrity (\texttt{Enterprise M1046}; \texttt{ICS M0946})
\begin{itemize}
\item Verifies the boot process, operating system, and loading mechanisms to prevent tampered components from starting trusted systems.
\end{itemize}
\item Execution Prevention (\texttt{Enterprise M1038}; \texttt{ICS M0938})
\begin{itemize}
\item Blocks unauthorized or malicious code execution through application control, script blocking, or related controls.
\end{itemize}
\item Application Isolation and Sandboxing (\texttt{Enterprise M1048}; \texttt{ICS M0948})
\begin{itemize}
\item Runs code in a controlled, isolated environment to limit access to sensitive resources and reduce impact.
\end{itemize}
\item Privileged Process Integrity (\texttt{Enterprise M1025})
\begin{itemize}
\item Protects highly privileged processes from tampering, injection, or compromise.
\end{itemize}
\item Software Process and Device Authentication (\texttt{ICS M0813})
\begin{itemize}
\item Requires devices and software processes to authenticate, especially for remote connections and API access.
\end{itemize}
\item Communication Authenticity (\texttt{ICS M0802})
\begin{itemize}
\item Authenticates message senders and verifies message integrity to detect spoofed messages or unauthorized connections.
\end{itemize}
\item Authorization Enforcement (\texttt{ICS M0800})
\begin{itemize}
\item Restricts read, manipulation, or execution privileges to authenticated users who require access under approved policies.
\end{itemize}
\item Access Management (\texttt{ICS M0801})
\begin{itemize}
\item Enforces access decisions, often through gateways or authentication services, before users can reach systems or devices.
\end{itemize}
\end{itemize}
\end{itemize}

\begin{itemize}
\item \textbf{Supporting mitigations:}
\begin{itemize}
\item Restrict File and Directory Permissions (\texttt{Enterprise M1022}; \texttt{ICS M0922})
\begin{itemize}
\item Limits file and directory access so only authorized users, groups, or processes can read, write, or execute them.
\end{itemize}
\item Privileged Account Management (\texttt{Enterprise M1026}; \texttt{ICS M0926})
\begin{itemize}
\item Controls creation, use, permissions, and accountability for administrative, root, SYSTEM, or other privileged accounts.
\end{itemize}
\item Credential Access Protection (\texttt{Enterprise M1043})
\begin{itemize}
\item Protects credentials by hardening credential storage, limiting access, and detecting suspicious credential-related activity.
\end{itemize}
\item Multi-factor Authentication (\texttt{Enterprise M1032}; \texttt{ICS M0932})
\begin{itemize}
\item Requires two or more forms of identity evidence before granting access to a system.
\end{itemize}
\item Software Configuration (\texttt{Enterprise M1054}; \texttt{ICS M0954})
\begin{itemize}
\item Applies security-focused software setting changes to reduce risky behavior and attack surface.
\end{itemize}
\item Operating System Configuration (\texttt{Enterprise M1028}; \texttt{ICS M0928})
\begin{itemize}
\item Hardens operating-system settings to reduce attack surface and opportunities for adversary abuse.
\end{itemize}
\item Update Software (\texttt{Enterprise M1051}; \texttt{ICS M0951})
\begin{itemize}
\item Applies vendor patches and upgrades to reduce exploitation of known software or firmware vulnerabilities.
\end{itemize}
\item Vulnerability Scanning (\texttt{Enterprise M1016}; \texttt{ICS M0916})
\begin{itemize}
\item Assesses systems, applications, and networks for vulnerabilities, misconfigurations, or unpatched software to prioritize remediation.
\end{itemize}
\item Mechanical Protection Layers (\texttt{ICS M0805})
\begin{itemize}
\item Uses physical or mechanical protection layers, such as interlocks or safety devices, to prevent damage or harm.
\end{itemize}
\item Safety Instrumented Systems (\texttt{ICS M0812})
\begin{itemize}
\item Uses independent safety systems to detect hazardous conditions and automatically drive equipment to a safer state.
\end{itemize}
\item Limit Hardware Installation (\texttt{Enterprise M1034}; \texttt{ICS M0934})
\begin{itemize}
\item Blocks users or groups from installing or using unapproved hardware, including removable devices.
\end{itemize}
\end{itemize}
\end{itemize}

\begin{itemize}
\item \textbf{AM-specific application:} enforce signed firmware, signed machine-control software, authenticated build-job submission, protected license or entitlement checks, hardened boot chains, and physical interlocks around machine safety or production-limit bypass paths.
\end{itemize}

\begin{itemize}
\item \textbf{Residual coverage note:} ATT\&CK covers many cyber bypass paths, but legal-rights enforcement, DRM business logic, watermark robustness, and physical tamper resistance need separate product-security and governance treatment.
\end{itemize}


\subsection{3.2.2 Counterfeiting / unlicensed reproduction [I]}


Unauthorised production or sale of copied AM parts or designs, including counterfeit products and unlicensed replicas that may use inferior materials, lack safety compliance, or damage brand and economic value.

Normalized umbrella category capturing the reproduction dimension in Kumar's goals/targets and the practical consequence of stolen or reconstructed IP.

Includes:

\begin{itemize}
\item counterfeit production
\item "Counterfeiting Part Protection" attack logic from Kumar's method list
\item illicit replication
\item unauthorized resale reproduction
\item grey-market part cloning
\end{itemize}

\textbf{MITRE ATT\&CK mitigation mapping (derived):}

\begin{itemize}
\item \textbf{Primary mitigations:}
\begin{itemize}
\item Supply Chain Management (\texttt{ICS M0817})
\begin{itemize}
\item Uses supplier policies, procedures, and integrity testing to ensure devices and components come from trusted sources.
\end{itemize}
\item Operational Information Confidentiality (\texttt{ICS M0809})
\begin{itemize}
\item Protects operational-process, facility, configuration, program, or database information that could reveal sensitive know-how or IP.
\end{itemize}
\item Data Loss Prevention (\texttt{Enterprise M1057}; \texttt{ICS M0803})
\begin{itemize}
\item Identifies, monitors, and controls sensitive-data movement to reduce unauthorized access, transmission, or exfiltration.
\end{itemize}
\item Encrypt Sensitive Information (\texttt{Enterprise M1041}; \texttt{ICS M0941})
\begin{itemize}
\item Uses strong encryption to protect sensitive information against unauthorized access or tampering.
\end{itemize}
\item Code Signing (\texttt{Enterprise M1045}; \texttt{ICS M0945})
\begin{itemize}
\item Verifies digital signatures so only trusted, untampered code is allowed to execute.
\end{itemize}
\item Communication Authenticity (\texttt{ICS M0802})
\begin{itemize}
\item Authenticates message senders and verifies message integrity to detect spoofed messages or unauthorized connections.
\end{itemize}
\item Software Process and Device Authentication (\texttt{ICS M0813})
\begin{itemize}
\item Requires devices and software processes to authenticate, especially for remote connections and API access.
\end{itemize}
\item Authorization Enforcement (\texttt{ICS M0800})
\begin{itemize}
\item Restricts read, manipulation, or execution privileges to authenticated users who require access under approved policies.
\end{itemize}
\item Access Management (\texttt{ICS M0801})
\begin{itemize}
\item Enforces access decisions, often through gateways or authentication services, before users can reach systems or devices.
\end{itemize}
\item Multi-factor Authentication (\texttt{Enterprise M1032}; \texttt{ICS M0932})
\begin{itemize}
\item Requires two or more forms of identity evidence before granting access to a system.
\end{itemize}
\end{itemize}
\end{itemize}

\begin{itemize}
\item \textbf{Supporting mitigations:}
\begin{itemize}
\item Privileged Account Management (\texttt{Enterprise M1026}; \texttt{ICS M0926})
\begin{itemize}
\item Controls creation, use, permissions, and accountability for administrative, root, SYSTEM, or other privileged accounts.
\end{itemize}
\item User Account Management (\texttt{Enterprise M1018}; \texttt{ICS M0918})
\begin{itemize}
\item Manages account creation, modification, use, and permissions to reduce unauthorized access and privilege misuse.
\end{itemize}
\item Audit (\texttt{Enterprise M1047}; \texttt{ICS M0947})
\begin{itemize}
\item Records and reviews activity, permissions, configurations, and integrity indicators to identify weaknesses or anomalies.
\end{itemize}
\item Remote Data Storage (\texttt{Enterprise M1029})
\begin{itemize}
\item Moves critical logs or sensitive files to secure off-host storage to reduce tampering, destruction, or unauthorized access.
\end{itemize}
\item Network Segmentation (\texttt{Enterprise M1030}; \texttt{ICS M0930})
\begin{itemize}
\item Divides networks into isolated segments or zones to restrict access, reduce lateral movement, and protect critical systems.
\end{itemize}
\item Limit Access to Resource Over Network (\texttt{Enterprise M1035}; \texttt{ICS M0935})
\begin{itemize}
\item Restricts network access to resources such as file shares, remote systems, and services to legitimate users or systems.
\end{itemize}
\item Software Configuration (\texttt{Enterprise M1054}; \texttt{ICS M0954})
\begin{itemize}
\item Applies security-focused software setting changes to reduce risky behavior and attack surface.
\end{itemize}
\end{itemize}
\end{itemize}

\begin{itemize}
\item \textbf{AM-specific application:} gate production authorization through authenticated users, signed build packages, approved machine identities, supplier controls, auditable production records, and controls that prevent unauthorized reuse of protected designs or process recipes.
\end{itemize}

\begin{itemize}
\item \textbf{Residual coverage note:} ATT\&CK mitigations do not by themselves solve commercial authorization, licensing, export-control, or contractual overproduction problems. Those require business controls and traceability in addition to technical enforcement.
\end{itemize}


\section{3.3 Design and specification tampering methods [I]}


Methods in this family alter the information that defines the part before execution, including geometry, properties, materials, tolerances, interfaces, digital model representations, and executable build instructions.


\subsection{3.3.1 Part specification tampering [I]}


Changing the specification of the 3D-printed component so its dimensions, geometry, material, tolerances, surface finish, features, assembly interfaces, or required properties no longer match the authorised intent.

Merged from Kumar's \textbf{Tempering with part specification} and Yampolskiy's \textbf{3D part specification}.

Includes:

\begin{itemize}
\item geometry modification
\item external-shape manipulation
\item internal-cavity manipulation
\item position changes
\item tolerance modification
\item support-feature changes
\item assembly-interface manipulation
\item source-material substitution in multi-material contexts
\item addition/removal of features
\item surface-finish specification changes
\item mass changes
\item center-of-mass changes
\item required-property changes
\begin{itemize}
\item mechanical
\item chemical
\item thermal
\end{itemize}
\item scaling inaccuracies
\end{itemize}

\textbf{MITRE ATT\&CK mitigation mapping (derived):}

\begin{itemize}
\item \textbf{Primary mitigations:}
\begin{itemize}
\item Authorization Enforcement (\texttt{ICS M0800})
\begin{itemize}
\item Restricts read, manipulation, or execution privileges to authenticated users who require access under approved policies.
\end{itemize}
\item Access Management (\texttt{ICS M0801})
\begin{itemize}
\item Enforces access decisions, often through gateways or authentication services, before users can reach systems or devices.
\end{itemize}
\item Human User Authentication (\texttt{ICS M0804})
\begin{itemize}
\item Requires users to authenticate before accessing data or sending commands to a device or system.
\end{itemize}
\item Multi-factor Authentication (\texttt{Enterprise M1032}; \texttt{ICS M0932})
\begin{itemize}
\item Requires two or more forms of identity evidence before granting access to a system.
\end{itemize}
\item Privileged Account Management (\texttt{Enterprise M1026}; \texttt{ICS M0926})
\begin{itemize}
\item Controls creation, use, permissions, and accountability for administrative, root, SYSTEM, or other privileged accounts.
\end{itemize}
\item Code Signing (\texttt{Enterprise M1045}; \texttt{ICS M0945})
\begin{itemize}
\item Verifies digital signatures so only trusted, untampered code is allowed to execute.
\end{itemize}
\item Communication Authenticity (\texttt{ICS M0802})
\begin{itemize}
\item Authenticates message senders and verifies message integrity to detect spoofed messages or unauthorized connections.
\end{itemize}
\item Software Process and Device Authentication (\texttt{ICS M0813})
\begin{itemize}
\item Requires devices and software processes to authenticate, especially for remote connections and API access.
\end{itemize}
\item Restrict File and Directory Permissions (\texttt{Enterprise M1022}; \texttt{ICS M0922})
\begin{itemize}
\item Limits file and directory access so only authorized users, groups, or processes can read, write, or execute them.
\end{itemize}
\end{itemize}
\end{itemize}

\begin{itemize}
\item \textbf{Supporting mitigations:}
\begin{itemize}
\item Audit (\texttt{Enterprise M1047}; \texttt{ICS M0947})
\begin{itemize}
\item Records and reviews activity, permissions, configurations, and integrity indicators to identify weaknesses or anomalies.
\end{itemize}
\item Remote Data Storage (\texttt{Enterprise M1029})
\begin{itemize}
\item Moves critical logs or sensitive files to secure off-host storage to reduce tampering, destruction, or unauthorized access.
\end{itemize}
\item Data Backup (\texttt{Enterprise M1053}; \texttt{ICS M0953})
\begin{itemize}
\item Creates and protects backups of critical data and configurations so systems can recover after compromise or disruption.
\end{itemize}
\item Application Developer Guidance (\texttt{Enterprise M1013}; \texttt{ICS M0913})
\begin{itemize}
\item Provides secure-development guidance to reduce weaknesses in applications, systems, and APIs.
\end{itemize}
\item Validate Program Inputs (\texttt{ICS M0818})
\begin{itemize}
\item Checks input and action values against expected format, content, and operational ranges before accepting them.
\end{itemize}
\item Software Configuration (\texttt{Enterprise M1054}; \texttt{ICS M0954})
\begin{itemize}
\item Applies security-focused software setting changes to reduce risky behavior and attack surface.
\end{itemize}
\item Network Segmentation (\texttt{Enterprise M1030}; \texttt{ICS M0930})
\begin{itemize}
\item Divides networks into isolated segments or zones to restrict access, reduce lateral movement, and protect critical systems.
\end{itemize}
\item Limit Access to Resource Over Network (\texttt{Enterprise M1035}; \texttt{ICS M0935})
\begin{itemize}
\item Restricts network access to resources such as file shares, remote systems, and services to legitimate users or systems.
\end{itemize}
\end{itemize}
\end{itemize}

\begin{itemize}
\item \textbf{AM-specific application:} preserve approved part specifications through controlled design repositories, signed design baselines, peer-reviewed release workflows, auditable change histories, protected conversion pipelines, and validation of geometry or tolerance values before manufacturing release.
\end{itemize}

\begin{itemize}
\item \textbf{Residual coverage note:} ATT\&CK can support integrity protection of the digital specification, but it does not define AM engineering acceptance criteria for detecting malicious geometry changes; those criteria must come from design assurance, simulation, inspection, and domain-specific quality controls.
\end{itemize}


\subsection{3.3.2 CAD/STL/intermediate-representation tampering [I]}


Tampering with CAD, STL, mesh, or intermediate representations in the AM tool chain to introduce design changes, hidden cavities, file corruption, or conversion-stage defects before slicing or execution.

A practical subfamily separated out for implementation use.

Includes:

\begin{itemize}
\item CAD model alteration
\item STL mesh alteration
\item hidden-cavity insertion
\item intermediate-file corruption
\item conversion-stage manipulation
\end{itemize}

\textbf{MITRE ATT\&CK mitigation mapping (derived):}

\begin{itemize}
\item \textbf{Primary mitigations:}
\begin{itemize}
\item Code Signing (\texttt{Enterprise M1045}; \texttt{ICS M0945})
\begin{itemize}
\item Verifies digital signatures so only trusted, untampered code is allowed to execute.
\end{itemize}
\item Communication Authenticity (\texttt{ICS M0802})
\begin{itemize}
\item Authenticates message senders and verifies message integrity to detect spoofed messages or unauthorized connections.
\end{itemize}
\item Software Process and Device Authentication (\texttt{ICS M0813})
\begin{itemize}
\item Requires devices and software processes to authenticate, especially for remote connections and API access.
\end{itemize}
\item Restrict File and Directory Permissions (\texttt{Enterprise M1022}; \texttt{ICS M0922})
\begin{itemize}
\item Limits file and directory access so only authorized users, groups, or processes can read, write, or execute them.
\end{itemize}
\item Authorization Enforcement (\texttt{ICS M0800})
\begin{itemize}
\item Restricts read, manipulation, or execution privileges to authenticated users who require access under approved policies.
\end{itemize}
\item Access Management (\texttt{ICS M0801})
\begin{itemize}
\item Enforces access decisions, often through gateways or authentication services, before users can reach systems or devices.
\end{itemize}
\item Multi-factor Authentication (\texttt{Enterprise M1032}; \texttt{ICS M0932})
\begin{itemize}
\item Requires two or more forms of identity evidence before granting access to a system.
\end{itemize}
\end{itemize}
\end{itemize}

\begin{itemize}
\item \textbf{Supporting mitigations:}
\begin{itemize}
\item Application Developer Guidance (\texttt{Enterprise M1013}; \texttt{ICS M0913})
\begin{itemize}
\item Provides secure-development guidance to reduce weaknesses in applications, systems, and APIs.
\end{itemize}
\item Validate Program Inputs (\texttt{ICS M0818})
\begin{itemize}
\item Checks input and action values against expected format, content, and operational ranges before accepting them.
\end{itemize}
\item Application Isolation and Sandboxing (\texttt{Enterprise M1048}; \texttt{ICS M0948})
\begin{itemize}
\item Runs code in a controlled, isolated environment to limit access to sensitive resources and reduce impact.
\end{itemize}
\item Execution Prevention (\texttt{Enterprise M1038}; \texttt{ICS M0938})
\begin{itemize}
\item Blocks unauthorized or malicious code execution through application control, script blocking, or related controls.
\end{itemize}
\item Antivirus/Antimalware (\texttt{Enterprise M1049}; \texttt{ICS M0949})
\begin{itemize}
\item Uses signatures, heuristics, or behavior analysis to detect, block, and remediate malicious software.
\end{itemize}
\item Exploit Protection (\texttt{Enterprise M1050}; \texttt{ICS M0950})
\begin{itemize}
\item Detects, blocks, or mitigates conditions that indicate or enable software exploitation.
\end{itemize}
\item Audit (\texttt{Enterprise M1047}; \texttt{ICS M0947})
\begin{itemize}
\item Records and reviews activity, permissions, configurations, and integrity indicators to identify weaknesses or anomalies.
\end{itemize}
\item Data Backup (\texttt{Enterprise M1053}; \texttt{ICS M0953})
\begin{itemize}
\item Creates and protects backups of critical data and configurations so systems can recover after compromise or disruption.
\end{itemize}
\item Remote Data Storage (\texttt{Enterprise M1029})
\begin{itemize}
\item Moves critical logs or sensitive files to secure off-host storage to reduce tampering, destruction, or unauthorized access.
\end{itemize}
\end{itemize}
\end{itemize}

\begin{itemize}
\item \textbf{AM-specific application:} require integrity checks and approval states across CAD, STL, AMF, 3MF, mesh repair, nesting, and conversion tools; preserve known-good originals and log every conversion-stage change that can alter geometry or manufacturability.
\end{itemize}

\begin{itemize}
\item \textbf{Residual coverage note:} ATT\&CK mitigations can protect files and tools, but hidden-cavity or subtle mesh-manipulation detection still needs AM-aware geometry diffing, simulation, inspection, and release-review practices.
\end{itemize}


\subsection{3.3.3 Tool-path and execution-file tampering [I]}


Changing slicing outputs, tool paths, G-code or equivalent machine instructions, or execution sequencing so the machine receives build instructions that differ from the approved model or process.

Includes:

\begin{itemize}
\item slicing manipulation
\item tool-path corruption
\item machine-instruction modification
\item command sequencing manipulation
\item timing/order changes in executable build instructions
\end{itemize}

\textbf{MITRE ATT\&CK mitigation mapping (derived):}

\begin{itemize}
\item \textbf{Primary mitigations:}
\begin{itemize}
\item Code Signing (\texttt{Enterprise M1045}; \texttt{ICS M0945})
\begin{itemize}
\item Verifies digital signatures so only trusted, untampered code is allowed to execute.
\end{itemize}
\item Communication Authenticity (\texttt{ICS M0802})
\begin{itemize}
\item Authenticates message senders and verifies message integrity to detect spoofed messages or unauthorized connections.
\end{itemize}
\item Software Process and Device Authentication (\texttt{ICS M0813})
\begin{itemize}
\item Requires devices and software processes to authenticate, especially for remote connections and API access.
\end{itemize}
\item Validate Program Inputs (\texttt{ICS M0818})
\begin{itemize}
\item Checks input and action values against expected format, content, and operational ranges before accepting them.
\end{itemize}
\item Authorization Enforcement (\texttt{ICS M0800})
\begin{itemize}
\item Restricts read, manipulation, or execution privileges to authenticated users who require access under approved policies.
\end{itemize}
\item Access Management (\texttt{ICS M0801})
\begin{itemize}
\item Enforces access decisions, often through gateways or authentication services, before users can reach systems or devices.
\end{itemize}
\item Multi-factor Authentication (\texttt{Enterprise M1032}; \texttt{ICS M0932})
\begin{itemize}
\item Requires two or more forms of identity evidence before granting access to a system.
\end{itemize}
\item Filter Network Traffic (\texttt{ICS M0937}; \texttt{Enterprise M1037})
\begin{itemize}
\item Filters ingress, egress, or lateral traffic, including protocol-aware rules, to restrict malicious or unnecessary communication.
\end{itemize}
\item Network Allowlists (\texttt{ICS M0807})
\begin{itemize}
\item Allows only approved device connections based on defined addresses, ports, protocols, or related network properties.
\end{itemize}
\end{itemize}
\end{itemize}

\begin{itemize}
\item \textbf{Supporting mitigations:}
\begin{itemize}
\item Network Segmentation (\texttt{Enterprise M1030}; \texttt{ICS M0930})
\begin{itemize}
\item Divides networks into isolated segments or zones to restrict access, reduce lateral movement, and protect critical systems.
\end{itemize}
\item Static Network Configuration (\texttt{ICS M0814})
\begin{itemize}
\item Uses static network settings where possible to reduce abuse of dynamic discovery or addressing protocols.
\end{itemize}
\item Software Configuration (\texttt{Enterprise M1054}; \texttt{ICS M0954})
\begin{itemize}
\item Applies security-focused software setting changes to reduce risky behavior and attack surface.
\end{itemize}
\item Execution Prevention (\texttt{Enterprise M1038}; \texttt{ICS M0938})
\begin{itemize}
\item Blocks unauthorized or malicious code execution through application control, script blocking, or related controls.
\end{itemize}
\item Application Isolation and Sandboxing (\texttt{Enterprise M1048}; \texttt{ICS M0948})
\begin{itemize}
\item Runs code in a controlled, isolated environment to limit access to sensitive resources and reduce impact.
\end{itemize}
\item Audit (\texttt{Enterprise M1047}; \texttt{ICS M0947})
\begin{itemize}
\item Records and reviews activity, permissions, configurations, and integrity indicators to identify weaknesses or anomalies.
\end{itemize}
\item Data Backup (\texttt{Enterprise M1053}; \texttt{ICS M0953})
\begin{itemize}
\item Creates and protects backups of critical data and configurations so systems can recover after compromise or disruption.
\end{itemize}
\item Remote Data Storage (\texttt{Enterprise M1029})
\begin{itemize}
\item Moves critical logs or sensitive files to secure off-host storage to reduce tampering, destruction, or unauthorized access.
\end{itemize}
\end{itemize}
\end{itemize}

\begin{itemize}
\item \textbf{AM-specific application:} sign and authenticate sliced build jobs, G-code or equivalent machine instructions, machine profiles, tool-path files, queue submissions, and job-transfer channels before a machine accepts execution.
\end{itemize}

\begin{itemize}
\item \textbf{Residual coverage note:} ATT\&CK does not specify what a physically safe or mechanically valid AM tool path is; bounded process windows, machine-specific validation, and independent build simulation remain necessary.
\end{itemize}


\section{3.4 Process tampering methods [I, A]}


Methods in this family manipulate the live manufacturing process, control logic, sensing, timing, or power conditions that determine how material is formed during AM.


\subsection{3.4.1 Compromising AM process parameters [I, A]}


Unauthorised changes to AM process parameters such as laser or beam power, scan strategy, speed, hatch spacing, build direction, feedstock handling, or chamber/environment conditions, affecting part quality or properties.

Merged from Kumar's method and Yampolskiy's AM process manipulation logic.

Includes:

\begin{itemize}
\item laser/electron-beam/power setting changes
\item scan strategy changes
\item speed/path/hatch changes
\item build-direction changes
\item powder-spread or deposition changes
\item environmental process setting changes
\item process-control scanning-strategy manipulation
\item process-control build-direction manipulation
\end{itemize}

\textbf{MITRE ATT\&CK mitigation mapping (derived):}

\begin{itemize}
\item \textbf{Primary mitigations:}
\begin{itemize}
\item Validate Program Inputs (\texttt{ICS M0818})
\begin{itemize}
\item Checks input and action values against expected format, content, and operational ranges before accepting them.
\end{itemize}
\item Authorization Enforcement (\texttt{ICS M0800})
\begin{itemize}
\item Restricts read, manipulation, or execution privileges to authenticated users who require access under approved policies.
\end{itemize}
\item Access Management (\texttt{ICS M0801})
\begin{itemize}
\item Enforces access decisions, often through gateways or authentication services, before users can reach systems or devices.
\end{itemize}
\item Human User Authentication (\texttt{ICS M0804})
\begin{itemize}
\item Requires users to authenticate before accessing data or sending commands to a device or system.
\end{itemize}
\item Multi-factor Authentication (\texttt{Enterprise M1032}; \texttt{ICS M0932})
\begin{itemize}
\item Requires two or more forms of identity evidence before granting access to a system.
\end{itemize}
\item Communication Authenticity (\texttt{ICS M0802})
\begin{itemize}
\item Authenticates message senders and verifies message integrity to detect spoofed messages or unauthorized connections.
\end{itemize}
\item Software Process and Device Authentication (\texttt{ICS M0813})
\begin{itemize}
\item Requires devices and software processes to authenticate, especially for remote connections and API access.
\end{itemize}
\item Network Allowlists (\texttt{ICS M0807})
\begin{itemize}
\item Allows only approved device connections based on defined addresses, ports, protocols, or related network properties.
\end{itemize}
\item Filter Network Traffic (\texttt{ICS M0937}; \texttt{Enterprise M1037})
\begin{itemize}
\item Filters ingress, egress, or lateral traffic, including protocol-aware rules, to restrict malicious or unnecessary communication.
\end{itemize}
\end{itemize}
\end{itemize}

\begin{itemize}
\item \textbf{Supporting mitigations:}
\begin{itemize}
\item Network Segmentation (\texttt{Enterprise M1030}; \texttt{ICS M0930})
\begin{itemize}
\item Divides networks into isolated segments or zones to restrict access, reduce lateral movement, and protect critical systems.
\end{itemize}
\item Static Network Configuration (\texttt{ICS M0814})
\begin{itemize}
\item Uses static network settings where possible to reduce abuse of dynamic discovery or addressing protocols.
\end{itemize}
\item Limit Access to Resource Over Network (\texttt{Enterprise M1035}; \texttt{ICS M0935})
\begin{itemize}
\item Restricts network access to resources such as file shares, remote systems, and services to legitimate users or systems.
\end{itemize}
\item Software Configuration (\texttt{Enterprise M1054}; \texttt{ICS M0954})
\begin{itemize}
\item Applies security-focused software setting changes to reduce risky behavior and attack surface.
\end{itemize}
\item Audit (\texttt{Enterprise M1047}; \texttt{ICS M0947})
\begin{itemize}
\item Records and reviews activity, permissions, configurations, and integrity indicators to identify weaknesses or anomalies.
\end{itemize}
\item Mechanical Protection Layers (\texttt{ICS M0805})
\begin{itemize}
\item Uses physical or mechanical protection layers, such as interlocks or safety devices, to prevent damage or harm.
\end{itemize}
\item Safety Instrumented Systems (\texttt{ICS M0812})
\begin{itemize}
\item Uses independent safety systems to detect hazardous conditions and automatically drive equipment to a safer state.
\end{itemize}
\item Watchdog Timers (\texttt{ICS M0815})
\begin{itemize}
\item Uses timers to quickly detect when a device or system becomes unresponsive.
\end{itemize}
\item Redundancy of Service (\texttt{ICS M0811})
\begin{itemize}
\item Provides backup or standby devices and services so critical functions can continue during failures or disruption.
\end{itemize}
\end{itemize}
\end{itemize}

\begin{itemize}
\item \textbf{AM-specific application:} constrain laser/electron-beam power, scan strategy, hatch spacing, speed, build orientation, feedstock handling, chamber atmosphere, and machine profiles to authenticated, approved, range-checked process windows.
\end{itemize}

\begin{itemize}
\item \textbf{Residual coverage note:} ATT\&CK supports access, authentication, validation, and network control, but it does not define material- or machine-specific process limits. Those must come from qualified AM process specifications and production validation.
\end{itemize}


\subsection{3.4.2 Process-control manipulation [I, A]}


Manipulation of machine-control commands, control-loop behavior, or actuator instructions so AM equipment executes attacker-influenced actions during the build.

Includes:

\begin{itemize}
\item direct manipulation of machine-control commands
\item modification of control-loop behavior
\item intervention in actuator instructions
\end{itemize}

\textbf{MITRE ATT\&CK mitigation mapping (derived):}

\begin{itemize}
\item \textbf{Primary mitigations:}
\begin{itemize}
\item Authorization Enforcement (\texttt{ICS M0800})
\begin{itemize}
\item Restricts read, manipulation, or execution privileges to authenticated users who require access under approved policies.
\end{itemize}
\item Access Management (\texttt{ICS M0801})
\begin{itemize}
\item Enforces access decisions, often through gateways or authentication services, before users can reach systems or devices.
\end{itemize}
\item Human User Authentication (\texttt{ICS M0804})
\begin{itemize}
\item Requires users to authenticate before accessing data or sending commands to a device or system.
\end{itemize}
\item Multi-factor Authentication (\texttt{Enterprise M1032}; \texttt{ICS M0932})
\begin{itemize}
\item Requires two or more forms of identity evidence before granting access to a system.
\end{itemize}
\item Privileged Account Management (\texttt{Enterprise M1026}; \texttt{ICS M0926})
\begin{itemize}
\item Controls creation, use, permissions, and accountability for administrative, root, SYSTEM, or other privileged accounts.
\end{itemize}
\item Communication Authenticity (\texttt{ICS M0802})
\begin{itemize}
\item Authenticates message senders and verifies message integrity to detect spoofed messages or unauthorized connections.
\end{itemize}
\item Software Process and Device Authentication (\texttt{ICS M0813})
\begin{itemize}
\item Requires devices and software processes to authenticate, especially for remote connections and API access.
\end{itemize}
\item Validate Program Inputs (\texttt{ICS M0818})
\begin{itemize}
\item Checks input and action values against expected format, content, and operational ranges before accepting them.
\end{itemize}
\item Network Allowlists (\texttt{ICS M0807})
\begin{itemize}
\item Allows only approved device connections based on defined addresses, ports, protocols, or related network properties.
\end{itemize}
\end{itemize}
\end{itemize}

\begin{itemize}
\item \textbf{Supporting mitigations:}
\begin{itemize}
\item Filter Network Traffic (\texttt{ICS M0937}; \texttt{Enterprise M1037})
\begin{itemize}
\item Filters ingress, egress, or lateral traffic, including protocol-aware rules, to restrict malicious or unnecessary communication.
\end{itemize}
\item Network Segmentation (\texttt{ICS M0930}; \texttt{Enterprise M1030})
\begin{itemize}
\item Divides networks into isolated segments or zones to restrict access, reduce lateral movement, and protect critical systems.
\end{itemize}
\item Static Network Configuration (\texttt{ICS M0814})
\begin{itemize}
\item Uses static network settings where possible to reduce abuse of dynamic discovery or addressing protocols.
\end{itemize}
\item Network Intrusion Prevention (\texttt{ICS M0931}; \texttt{Enterprise M1031})
\begin{itemize}
\item Uses intrusion-detection signatures to block traffic at network boundaries, with care not to disrupt control or safety functions.
\end{itemize}
\item Code Signing (\texttt{Enterprise M1045}; \texttt{ICS M0945})
\begin{itemize}
\item Verifies digital signatures so only trusted, untampered code is allowed to execute.
\end{itemize}
\item Boot Integrity (\texttt{Enterprise M1046}; \texttt{ICS M0946})
\begin{itemize}
\item Verifies the boot process, operating system, and loading mechanisms to prevent tampered components from starting trusted systems.
\end{itemize}
\item Execution Prevention (\texttt{Enterprise M1038}; \texttt{ICS M0938})
\begin{itemize}
\item Blocks unauthorized or malicious code execution through application control, script blocking, or related controls.
\end{itemize}
\item Audit (\texttt{Enterprise M1047}; \texttt{ICS M0947})
\begin{itemize}
\item Records and reviews activity, permissions, configurations, and integrity indicators to identify weaknesses or anomalies.
\end{itemize}
\item Mechanical Protection Layers (\texttt{ICS M0805})
\begin{itemize}
\item Uses physical or mechanical protection layers, such as interlocks or safety devices, to prevent damage or harm.
\end{itemize}
\item Safety Instrumented Systems (\texttt{ICS M0812})
\begin{itemize}
\item Uses independent safety systems to detect hazardous conditions and automatically drive equipment to a safer state.
\end{itemize}
\item Watchdog Timers (\texttt{ICS M0815})
\begin{itemize}
\item Uses timers to quickly detect when a device or system becomes unresponsive.
\end{itemize}
\end{itemize}
\end{itemize}

\begin{itemize}
\item \textbf{AM-specific application:} isolate and authenticate machine-control interfaces, controller APIs, HMI commands, actuator paths, build queues, maintenance sessions, and closed-loop control channels before they can alter live machine behavior.
\end{itemize}

\begin{itemize}
\item \textbf{Residual coverage note:} real-time production constraints may limit heavy inline inspection or endpoint controls; representative test validation is needed before deploying mitigations to production machines.
\end{itemize}


\subsection{3.4.3 Sensor-information falsification [I]}


Falsification of sensor data used by AM monitoring or control loops, such as oxygen, pressure, chamber, infrared, or quality-control readings, causing incorrect state estimation or unsafe process adjustments.

A distinct Yampolskiy sabotage method retained explicitly.

Includes:

\begin{itemize}
\item false state-estimation attacks
\item manipulated IR thermography data
\item false chamber-condition readings
\item compromised quality-control sensor streams
\item false oxygen readings
\item false infrared readings
\item additional sensor-information falsification where the source taxonomy used open-ended ellipsis nodes
\end{itemize}

\textbf{MITRE ATT\&CK mitigation mapping (derived):}

\begin{itemize}
\item \textbf{Primary mitigations:}
\begin{itemize}
\item Communication Authenticity (\texttt{ICS M0802})
\begin{itemize}
\item Authenticates message senders and verifies message integrity to detect spoofed messages or unauthorized connections.
\end{itemize}
\item Software Process and Device Authentication (\texttt{ICS M0813})
\begin{itemize}
\item Requires devices and software processes to authenticate, especially for remote connections and API access.
\end{itemize}
\item Validate Program Inputs (\texttt{ICS M0818})
\begin{itemize}
\item Checks input and action values against expected format, content, and operational ranges before accepting them.
\end{itemize}
\item Network Allowlists (\texttt{ICS M0807})
\begin{itemize}
\item Allows only approved device connections based on defined addresses, ports, protocols, or related network properties.
\end{itemize}
\item Filter Network Traffic (\texttt{ICS M0937}; \texttt{Enterprise M1037})
\begin{itemize}
\item Filters ingress, egress, or lateral traffic, including protocol-aware rules, to restrict malicious or unnecessary communication.
\end{itemize}
\item Network Segmentation (\texttt{ICS M0930}; \texttt{Enterprise M1030})
\begin{itemize}
\item Divides networks into isolated segments or zones to restrict access, reduce lateral movement, and protect critical systems.
\end{itemize}
\item Static Network Configuration (\texttt{ICS M0814})
\begin{itemize}
\item Uses static network settings where possible to reduce abuse of dynamic discovery or addressing protocols.
\end{itemize}
\end{itemize}
\end{itemize}

\begin{itemize}
\item \textbf{Supporting mitigations:}
\begin{itemize}
\item Encrypt Network Traffic (\texttt{ICS M0808})
\begin{itemize}
\item Uses strong cryptographic protocols to reduce eavesdropping on network communications.
\end{itemize}
\item Audit (\texttt{Enterprise M1047}; \texttt{ICS M0947})
\begin{itemize}
\item Records and reviews activity, permissions, configurations, and integrity indicators to identify weaknesses or anomalies.
\end{itemize}
\item Watchdog Timers (\texttt{ICS M0815})
\begin{itemize}
\item Uses timers to quickly detect when a device or system becomes unresponsive.
\end{itemize}
\item Redundancy of Service (\texttt{ICS M0811})
\begin{itemize}
\item Provides backup or standby devices and services so critical functions can continue during failures or disruption.
\end{itemize}
\item Safety Instrumented Systems (\texttt{ICS M0812})
\begin{itemize}
\item Uses independent safety systems to detect hazardous conditions and automatically drive equipment to a safer state.
\end{itemize}
\item Mechanical Protection Layers (\texttt{ICS M0805})
\begin{itemize}
\item Uses physical or mechanical protection layers, such as interlocks or safety devices, to prevent damage or harm.
\end{itemize}
\item Software Configuration (\texttt{ICS M0954}; \texttt{Enterprise M1054})
\begin{itemize}
\item Applies security-focused software setting changes to reduce risky behavior and attack surface.
\end{itemize}
\item Update Software (\texttt{ICS M0951}; \texttt{Enterprise M1051})
\begin{itemize}
\item Applies vendor patches and upgrades to reduce exploitation of known software or firmware vulnerabilities.
\end{itemize}
\end{itemize}
\end{itemize}

\begin{itemize}
\item \textbf{AM-specific application:} authenticate and sanity-check oxygen, infrared, chamber, melt-pool, environmental, and inspection data streams; compare independent sensor channels where possible; retain trusted monitoring logs outside the machine being controlled.
\end{itemize}

\begin{itemize}
\item \textbf{Residual coverage note:} ATT\&CK mitigations help protect sensor communications and software paths, but compromised physical sensors, calibration drift, or spoofed physical phenomena require engineering calibration, redundancy, and process-quality controls beyond ATT\&CK.
\end{itemize}


\subsection{3.4.4 Timing disruption [I, A]}


Disruption of the timing or ordering of machine-control, sensor, or network messages, where early, late, delayed, or out-of-order data can change cyber-physical behavior.

Explicitly retained because Yampolskiy treats timing as a meaningful CPS-specific manipulation.

Includes:

\begin{itemize}
\item commands too early
\item commands too late
\item out-of-order packets
\item delayed or reordered sensor readings
\end{itemize}

\textbf{MITRE ATT\&CK mitigation mapping (derived):}

\begin{itemize}
\item \textbf{Primary mitigations:}
\begin{itemize}
\item Static Network Configuration (\texttt{ICS M0814})
\begin{itemize}
\item Uses static network settings where possible to reduce abuse of dynamic discovery or addressing protocols.
\end{itemize}
\item Network Allowlists (\texttt{ICS M0807})
\begin{itemize}
\item Allows only approved device connections based on defined addresses, ports, protocols, or related network properties.
\end{itemize}
\item Filter Network Traffic (\texttt{ICS M0937}; \texttt{Enterprise M1037})
\begin{itemize}
\item Filters ingress, egress, or lateral traffic, including protocol-aware rules, to restrict malicious or unnecessary communication.
\end{itemize}
\item Network Segmentation (\texttt{ICS M0930}; \texttt{Enterprise M1030})
\begin{itemize}
\item Divides networks into isolated segments or zones to restrict access, reduce lateral movement, and protect critical systems.
\end{itemize}
\item Communication Authenticity (\texttt{ICS M0802})
\begin{itemize}
\item Authenticates message senders and verifies message integrity to detect spoofed messages or unauthorized connections.
\end{itemize}
\item Software Process and Device Authentication (\texttt{ICS M0813})
\begin{itemize}
\item Requires devices and software processes to authenticate, especially for remote connections and API access.
\end{itemize}
\item Watchdog Timers (\texttt{ICS M0815})
\begin{itemize}
\item Uses timers to quickly detect when a device or system becomes unresponsive.
\end{itemize}
\end{itemize}
\end{itemize}

\begin{itemize}
\item \textbf{Supporting mitigations:}
\begin{itemize}
\item Redundancy of Service (\texttt{ICS M0811})
\begin{itemize}
\item Provides backup or standby devices and services so critical functions can continue during failures or disruption.
\end{itemize}
\item Out-of-Band Communications Channel (\texttt{ICS M0810}; \texttt{Enterprise M1060})
\begin{itemize}
\item Provides alternative communication paths when primary network communication fails, is disrupted, or cannot be trusted.
\end{itemize}
\item Network Intrusion Prevention (\texttt{ICS M0931}; \texttt{Enterprise M1031})
\begin{itemize}
\item Uses intrusion-detection signatures to block traffic at network boundaries, with care not to disrupt control or safety functions.
\end{itemize}
\item Audit (\texttt{Enterprise M1047}; \texttt{ICS M0947})
\begin{itemize}
\item Records and reviews activity, permissions, configurations, and integrity indicators to identify weaknesses or anomalies.
\end{itemize}
\item Mechanical Protection Layers (\texttt{ICS M0805})
\begin{itemize}
\item Uses physical or mechanical protection layers, such as interlocks or safety devices, to prevent damage or harm.
\end{itemize}
\item Safety Instrumented Systems (\texttt{ICS M0812})
\begin{itemize}
\item Uses independent safety systems to detect hazardous conditions and automatically drive equipment to a safer state.
\end{itemize}
\item Mitigation Limited or Not Effective (\texttt{ICS M0816})
\begin{itemize}
\item Marks techniques where preventive controls may be limited because the attack abuses legitimate system features.
\end{itemize}
\end{itemize}
\end{itemize}

\begin{itemize}
\item \textbf{AM-specific application:} protect deterministic production networks, job sequencing, sensor feedback timing, HMI-to-controller traffic, and machine-control messages from unauthorized reordering, delay, replay, or command-rate manipulation.
\end{itemize}

\begin{itemize}
\item \textbf{Residual coverage note:} where timing disruption arises from legitimate network jitter, equipment limitations, or abuse of allowed control features, preventive controls may be limited; monitoring, safe-state logic, watchdogs, and response planning become central.
\end{itemize}


\subsection{3.4.5 Power manipulation / power surge [I, A]}


Use of power instability, surge behavior, or coordinated machine power draw to interrupt production, damage facility power systems, or harm AM equipment.

Retained from Yampolskiy's sabotage logic.

Includes:

\begin{itemize}
\item disruption of power quality
\item coordinated surge behavior
\item power instability used to damage or interrupt equipment
\end{itemize}

\textbf{MITRE ATT\&CK mitigation mapping (derived):}

\begin{itemize}
\item \textbf{Primary mitigations:}
\begin{itemize}
\item Mechanical Protection Layers (\texttt{ICS M0805})
\begin{itemize}
\item Uses physical or mechanical protection layers, such as interlocks or safety devices, to prevent damage or harm.
\end{itemize}
\item Safety Instrumented Systems (\texttt{ICS M0812})
\begin{itemize}
\item Uses independent safety systems to detect hazardous conditions and automatically drive equipment to a safer state.
\end{itemize}
\item Redundancy of Service (\texttt{ICS M0811})
\begin{itemize}
\item Provides backup or standby devices and services so critical functions can continue during failures or disruption.
\end{itemize}
\item Watchdog Timers (\texttt{ICS M0815})
\begin{itemize}
\item Uses timers to quickly detect when a device or system becomes unresponsive.
\end{itemize}
\item Out-of-Band Communications Channel (\texttt{ICS M0810}; \texttt{Enterprise M1060})
\begin{itemize}
\item Provides alternative communication paths when primary network communication fails, is disrupted, or cannot be trusted.
\end{itemize}
\item Mitigation Limited or Not Effective (\texttt{ICS M0816})
\begin{itemize}
\item Marks techniques where preventive controls may be limited because the attack abuses legitimate system features.
\end{itemize}
\end{itemize}
\end{itemize}

\begin{itemize}
\item \textbf{Supporting mitigations:}
\begin{itemize}
\item Network Segmentation (\texttt{ICS M0930}; \texttt{Enterprise M1030})
\begin{itemize}
\item Divides networks into isolated segments or zones to restrict access, reduce lateral movement, and protect critical systems.
\end{itemize}
\item Filter Network Traffic (\texttt{ICS M0937}; \texttt{Enterprise M1037})
\begin{itemize}
\item Filters ingress, egress, or lateral traffic, including protocol-aware rules, to restrict malicious or unnecessary communication.
\end{itemize}
\item Network Allowlists (\texttt{ICS M0807})
\begin{itemize}
\item Allows only approved device connections based on defined addresses, ports, protocols, or related network properties.
\end{itemize}
\item Communication Authenticity (\texttt{ICS M0802})
\begin{itemize}
\item Authenticates message senders and verifies message integrity to detect spoofed messages or unauthorized connections.
\end{itemize}
\item Software Process and Device Authentication (\texttt{ICS M0813})
\begin{itemize}
\item Requires devices and software processes to authenticate, especially for remote connections and API access.
\end{itemize}
\item Data Backup (\texttt{ICS M0953}; \texttt{Enterprise M1053})
\begin{itemize}
\item Creates and protects backups of critical data and configurations so systems can recover after compromise or disruption.
\end{itemize}
\item Audit (\texttt{Enterprise M1047}; \texttt{ICS M0947})
\begin{itemize}
\item Records and reviews activity, permissions, configurations, and integrity indicators to identify weaknesses or anomalies.
\end{itemize}
\end{itemize}
\end{itemize}

\begin{itemize}
\item \textbf{AM-specific application:} protect cyber-accessible power-control paths, preserve safe shutdown behavior, maintain recoverable machine configurations, and use redundant or protective industrial layers around critical AM equipment.
\end{itemize}

\begin{itemize}
\item \textbf{Residual coverage note:} ATT\&CK is weak for electrical power-quality engineering. Surge protection, UPS design, facility electrical controls, grounding, and equipment certification are necessary controls but are not expressed as detailed ATT\&CK mitigations.
\end{itemize}


\section{3.5 Post-processing tampering methods [I]}


Methods in this family tamper with required post-build operations such as heat treatment, finishing, coating, cleaning, inspection, assembly, or material substitution, affecting final part properties or acceptance.


\subsection{3.5.1 Post-processing procedure tampering [I]}


Manipulation of post-processing recipes, operations, equipment, materials, or inspections after printing to introduce defects, alter properties, conceal flaws, or compromise part integrity.

Merged from both papers.

Includes:

\begin{itemize}
\item heat-treatment manipulation
\item HIP-cycle manipulation
\item finish-machining manipulation
\item coating/surface-treatment manipulation
\item concealment of defects during inspection
\item substitution of post-processing materials
\item sabotage of post-processing equipment
\end{itemize}

\textbf{MITRE ATT\&CK mitigation mapping (derived):}

\begin{itemize}
\item \textbf{Primary mitigations:}
\begin{itemize}
\item Authorization Enforcement (\texttt{ICS M0800})
\begin{itemize}
\item Restricts read, manipulation, or execution privileges to authenticated users who require access under approved policies.
\end{itemize}
\item Access Management (\texttt{ICS M0801})
\begin{itemize}
\item Enforces access decisions, often through gateways or authentication services, before users can reach systems or devices.
\end{itemize}
\item Human User Authentication (\texttt{ICS M0804})
\begin{itemize}
\item Requires users to authenticate before accessing data or sending commands to a device or system.
\end{itemize}
\item Multi-factor Authentication (\texttt{Enterprise M1032}; \texttt{ICS M0932})
\begin{itemize}
\item Requires two or more forms of identity evidence before granting access to a system.
\end{itemize}
\item Validate Program Inputs (\texttt{ICS M0818})
\begin{itemize}
\item Checks input and action values against expected format, content, and operational ranges before accepting them.
\end{itemize}
\item Code Signing (\texttt{Enterprise M1045}; \texttt{ICS M0945})
\begin{itemize}
\item Verifies digital signatures so only trusted, untampered code is allowed to execute.
\end{itemize}
\item Communication Authenticity (\texttt{ICS M0802})
\begin{itemize}
\item Authenticates message senders and verifies message integrity to detect spoofed messages or unauthorized connections.
\end{itemize}
\item Software Process and Device Authentication (\texttt{ICS M0813})
\begin{itemize}
\item Requires devices and software processes to authenticate, especially for remote connections and API access.
\end{itemize}
\item Audit (\texttt{Enterprise M1047}; \texttt{ICS M0947})
\begin{itemize}
\item Records and reviews activity, permissions, configurations, and integrity indicators to identify weaknesses or anomalies.
\end{itemize}
\end{itemize}
\end{itemize}

\begin{itemize}
\item \textbf{Supporting mitigations:}
\begin{itemize}
\item Network Segmentation (\texttt{ICS M0930}; \texttt{Enterprise M1030})
\begin{itemize}
\item Divides networks into isolated segments or zones to restrict access, reduce lateral movement, and protect critical systems.
\end{itemize}
\item Network Allowlists (\texttt{ICS M0807})
\begin{itemize}
\item Allows only approved device connections based on defined addresses, ports, protocols, or related network properties.
\end{itemize}
\item Filter Network Traffic (\texttt{ICS M0937}; \texttt{Enterprise M1037})
\begin{itemize}
\item Filters ingress, egress, or lateral traffic, including protocol-aware rules, to restrict malicious or unnecessary communication.
\end{itemize}
\item Software Configuration (\texttt{ICS M0954}; \texttt{Enterprise M1054})
\begin{itemize}
\item Applies security-focused software setting changes to reduce risky behavior and attack surface.
\end{itemize}
\item Data Backup (\texttt{ICS M0953}; \texttt{Enterprise M1053})
\begin{itemize}
\item Creates and protects backups of critical data and configurations so systems can recover after compromise or disruption.
\end{itemize}
\item Remote Data Storage (\texttt{Enterprise M1029})
\begin{itemize}
\item Moves critical logs or sensitive files to secure off-host storage to reduce tampering, destruction, or unauthorized access.
\end{itemize}
\item Mechanical Protection Layers (\texttt{ICS M0805})
\begin{itemize}
\item Uses physical or mechanical protection layers, such as interlocks or safety devices, to prevent damage or harm.
\end{itemize}
\item Safety Instrumented Systems (\texttt{ICS M0812})
\begin{itemize}
\item Uses independent safety systems to detect hazardous conditions and automatically drive equipment to a safer state.
\end{itemize}
\end{itemize}
\end{itemize}

\begin{itemize}
\item \textbf{AM-specific application:} require authenticated and approved post-processing recipes, inspection configurations, heat-treatment cycles, machining programs, coating parameters, and material substitutions; preserve tamper-evident records of post-processing results.
\end{itemize}

\begin{itemize}
\item \textbf{Residual coverage note:} ATT\&CK does not define metallurgical or surface-finish acceptance controls. Independent inspection, destructive or non-destructive testing, and qualified post-processing procedures remain necessary.
\end{itemize}


\subsection{3.5.2 Heat-treatment manipulation [I]}


Alteration of heat-treatment steps such as HIPing or annealing, including temperature, pressure, duration, or gas, in ways that can affect residual stress relief and mechanical properties.

Retained as a named submethod due to its importance in metal AM.

Includes:

\begin{itemize}
\item HIPing
\begin{itemize}
\item temperature
\item pressure
\item duration
\item gas
\end{itemize}
\item annealing
\end{itemize}

\textbf{MITRE ATT\&CK mitigation mapping (derived):}

\begin{itemize}
\item \textbf{Primary mitigations:}
\begin{itemize}
\item Validate Program Inputs (\texttt{ICS M0818})
\begin{itemize}
\item Checks input and action values against expected format, content, and operational ranges before accepting them.
\end{itemize}
\item Authorization Enforcement (\texttt{ICS M0800})
\begin{itemize}
\item Restricts read, manipulation, or execution privileges to authenticated users who require access under approved policies.
\end{itemize}
\item Access Management (\texttt{ICS M0801})
\begin{itemize}
\item Enforces access decisions, often through gateways or authentication services, before users can reach systems or devices.
\end{itemize}
\item Human User Authentication (\texttt{ICS M0804})
\begin{itemize}
\item Requires users to authenticate before accessing data or sending commands to a device or system.
\end{itemize}
\item Multi-factor Authentication (\texttt{Enterprise M1032}; \texttt{ICS M0932})
\begin{itemize}
\item Requires two or more forms of identity evidence before granting access to a system.
\end{itemize}
\item Communication Authenticity (\texttt{ICS M0802})
\begin{itemize}
\item Authenticates message senders and verifies message integrity to detect spoofed messages or unauthorized connections.
\end{itemize}
\item Software Process and Device Authentication (\texttt{ICS M0813})
\begin{itemize}
\item Requires devices and software processes to authenticate, especially for remote connections and API access.
\end{itemize}
\item Audit (\texttt{Enterprise M1047}; \texttt{ICS M0947})
\begin{itemize}
\item Records and reviews activity, permissions, configurations, and integrity indicators to identify weaknesses or anomalies.
\end{itemize}
\end{itemize}
\end{itemize}

\begin{itemize}
\item \textbf{Supporting mitigations:}
\begin{itemize}
\item Software Configuration (\texttt{ICS M0954}; \texttt{Enterprise M1054})
\begin{itemize}
\item Applies security-focused software setting changes to reduce risky behavior and attack surface.
\end{itemize}
\item Network Segmentation (\texttt{ICS M0930}; \texttt{Enterprise M1030})
\begin{itemize}
\item Divides networks into isolated segments or zones to restrict access, reduce lateral movement, and protect critical systems.
\end{itemize}
\item Network Allowlists (\texttt{ICS M0807})
\begin{itemize}
\item Allows only approved device connections based on defined addresses, ports, protocols, or related network properties.
\end{itemize}
\item Mechanical Protection Layers (\texttt{ICS M0805})
\begin{itemize}
\item Uses physical or mechanical protection layers, such as interlocks or safety devices, to prevent damage or harm.
\end{itemize}
\item Safety Instrumented Systems (\texttt{ICS M0812})
\begin{itemize}
\item Uses independent safety systems to detect hazardous conditions and automatically drive equipment to a safer state.
\end{itemize}
\item Watchdog Timers (\texttt{ICS M0815})
\begin{itemize}
\item Uses timers to quickly detect when a device or system becomes unresponsive.
\end{itemize}
\item Redundancy of Service (\texttt{ICS M0811})
\begin{itemize}
\item Provides backup or standby devices and services so critical functions can continue during failures or disruption.
\end{itemize}
\item Data Backup (\texttt{ICS M0953}; \texttt{Enterprise M1053})
\begin{itemize}
\item Creates and protects backups of critical data and configurations so systems can recover after compromise or disruption.
\end{itemize}
\end{itemize}
\end{itemize}

\begin{itemize}
\item \textbf{AM-specific application:} constrain HIP and annealing temperature, pressure, duration, atmosphere, ramp rates, and recipe selection to approved ranges and authenticated operator or system actions.
\end{itemize}

\begin{itemize}
\item \textbf{Residual coverage note:} the security controls protect recipe integrity, but proving that heat treatment produced the required microstructure and mechanical properties requires AM/materials engineering evidence outside ATT\&CK.
\end{itemize}


\subsection{3.5.3 Finish-machining manipulation [I]}


Manipulation of finish machining, EDM, surface-finish operations, offsets, or programs needed to achieve required fit, function, surface quality, and fatigue-related performance.

Retained as a named submethod due to its impact on fit, finish, and fatigue-related performance.

Includes:

\begin{itemize}
\item EDM
\item surface finish
\end{itemize}

\textbf{MITRE ATT\&CK mitigation mapping (derived):}

\begin{itemize}
\item \textbf{Primary mitigations:}
\begin{itemize}
\item Validate Program Inputs (\texttt{ICS M0818})
\begin{itemize}
\item Checks input and action values against expected format, content, and operational ranges before accepting them.
\end{itemize}
\item Authorization Enforcement (\texttt{ICS M0800})
\begin{itemize}
\item Restricts read, manipulation, or execution privileges to authenticated users who require access under approved policies.
\end{itemize}
\item Access Management (\texttt{ICS M0801})
\begin{itemize}
\item Enforces access decisions, often through gateways or authentication services, before users can reach systems or devices.
\end{itemize}
\item Code Signing (\texttt{Enterprise M1045}; \texttt{ICS M0945})
\begin{itemize}
\item Verifies digital signatures so only trusted, untampered code is allowed to execute.
\end{itemize}
\item Communication Authenticity (\texttt{ICS M0802})
\begin{itemize}
\item Authenticates message senders and verifies message integrity to detect spoofed messages or unauthorized connections.
\end{itemize}
\item Software Process and Device Authentication (\texttt{ICS M0813})
\begin{itemize}
\item Requires devices and software processes to authenticate, especially for remote connections and API access.
\end{itemize}
\item Audit (\texttt{Enterprise M1047}; \texttt{ICS M0947})
\begin{itemize}
\item Records and reviews activity, permissions, configurations, and integrity indicators to identify weaknesses or anomalies.
\end{itemize}
\end{itemize}
\end{itemize}

\begin{itemize}
\item \textbf{Supporting mitigations:}
\begin{itemize}
\item Network Segmentation (\texttt{ICS M0930}; \texttt{Enterprise M1030})
\begin{itemize}
\item Divides networks into isolated segments or zones to restrict access, reduce lateral movement, and protect critical systems.
\end{itemize}
\item Network Allowlists (\texttt{ICS M0807})
\begin{itemize}
\item Allows only approved device connections based on defined addresses, ports, protocols, or related network properties.
\end{itemize}
\item Filter Network Traffic (\texttt{ICS M0937}; \texttt{Enterprise M1037})
\begin{itemize}
\item Filters ingress, egress, or lateral traffic, including protocol-aware rules, to restrict malicious or unnecessary communication.
\end{itemize}
\item Software Configuration (\texttt{ICS M0954}; \texttt{Enterprise M1054})
\begin{itemize}
\item Applies security-focused software setting changes to reduce risky behavior and attack surface.
\end{itemize}
\item Execution Prevention (\texttt{Enterprise M1038}; \texttt{ICS M0938})
\begin{itemize}
\item Blocks unauthorized or malicious code execution through application control, script blocking, or related controls.
\end{itemize}
\item Mechanical Protection Layers (\texttt{ICS M0805})
\begin{itemize}
\item Uses physical or mechanical protection layers, such as interlocks or safety devices, to prevent damage or harm.
\end{itemize}
\item Safety Instrumented Systems (\texttt{ICS M0812})
\begin{itemize}
\item Uses independent safety systems to detect hazardous conditions and automatically drive equipment to a safer state.
\end{itemize}
\end{itemize}
\end{itemize}

\begin{itemize}
\item \textbf{AM-specific application:} protect CNC/EDM programs, fixture definitions, machining offsets, surface-finish instructions, inspection settings, and handoff data between AM build, post-processing, and quality-control systems.
\end{itemize}

\begin{itemize}
\item \textbf{Residual coverage note:} ATT\&CK does not cover all physical machining-quality assurance; independent dimensional inspection and surface/fatigue validation are still required.
\end{itemize}


\section{3.6 Supply-chain compromise methods [C, I, A]}


Methods in this family exploit upstream, downstream, or third-party AM supply-chain paths to introduce compromised materials, components, equipment, objects, software, or logistics disruption.


\subsection{3.6.1 Exploitation of supply-chain weaknesses [C, I, A]}


Compromise of weak suppliers, third-party systems, maintenance or update paths, material channels, or logistics actors to gain access, modify files, insert harmful components, or disrupt production.

Merged umbrella category covering upstream and downstream compromise opportunities.

Includes:

\begin{itemize}
\item supplier compromise
\item third-party software/firmware compromise
\item malicious maintenance/update pathways
\item counterfeit feedstock insertion
\item counterfeit component insertion
\item insider manipulation within supplier or logistics chain
\item supply-chain attack
\end{itemize}

\textbf{MITRE ATT\&CK mitigation mapping (derived):}

\begin{itemize}
\item \textbf{Primary mitigations:}
\begin{itemize}
\item Supply Chain Management (\texttt{ICS M0817})
\begin{itemize}
\item Uses supplier policies, procedures, and integrity testing to ensure devices and components come from trusted sources.
\end{itemize}
\item Code Signing (\texttt{Enterprise M1045}; \texttt{ICS M0945})
\begin{itemize}
\item Verifies digital signatures so only trusted, untampered code is allowed to execute.
\end{itemize}
\item Boot Integrity (\texttt{Enterprise M1046}; \texttt{ICS M0946})
\begin{itemize}
\item Verifies the boot process, operating system, and loading mechanisms to prevent tampered components from starting trusted systems.
\end{itemize}
\item Update Software (\texttt{Enterprise M1051}; \texttt{ICS M0951})
\begin{itemize}
\item Applies vendor patches and upgrades to reduce exploitation of known software or firmware vulnerabilities.
\end{itemize}
\item Vulnerability Scanning (\texttt{Enterprise M1016}; \texttt{ICS M0916})
\begin{itemize}
\item Assesses systems, applications, and networks for vulnerabilities, misconfigurations, or unpatched software to prioritize remediation.
\end{itemize}
\item Application Developer Guidance (\texttt{Enterprise M1013}; \texttt{ICS M0913})
\begin{itemize}
\item Provides secure-development guidance to reduce weaknesses in applications, systems, and APIs.
\end{itemize}
\item Threat Intelligence Program (\texttt{Enterprise M1019}; \texttt{ICS M0919})
\begin{itemize}
\item Collects and analyzes threat intelligence to track trends, guide defensive priorities, and reduce risk.
\end{itemize}
\end{itemize}
\end{itemize}

\begin{itemize}
\item \textbf{Supporting mitigations:}
\begin{itemize}
\item Limit Software Installation (\texttt{Enterprise M1033})
\begin{itemize}
\item Prevents installation of unauthorized or unapproved software to reduce malicious or vulnerable applications.
\end{itemize}
\item Limit Hardware Installation (\texttt{Enterprise M1034}; \texttt{ICS M0934})
\begin{itemize}
\item Blocks users or groups from installing or using unapproved hardware, including removable devices.
\end{itemize}
\item Execution Prevention (\texttt{Enterprise M1038}; \texttt{ICS M0938})
\begin{itemize}
\item Blocks unauthorized or malicious code execution through application control, script blocking, or related controls.
\end{itemize}
\item Antivirus/Antimalware (\texttt{Enterprise M1049}; \texttt{ICS M0949})
\begin{itemize}
\item Uses signatures, heuristics, or behavior analysis to detect, block, and remediate malicious software.
\end{itemize}
\item Data Loss Prevention (\texttt{Enterprise M1057}; \texttt{ICS M0803})
\begin{itemize}
\item Identifies, monitors, and controls sensitive-data movement to reduce unauthorized access, transmission, or exfiltration.
\end{itemize}
\item Operational Information Confidentiality (\texttt{ICS M0809})
\begin{itemize}
\item Protects operational-process, facility, configuration, program, or database information that could reveal sensitive know-how or IP.
\end{itemize}
\item Audit (\texttt{Enterprise M1047}; \texttt{ICS M0947})
\begin{itemize}
\item Records and reviews activity, permissions, configurations, and integrity indicators to identify weaknesses or anomalies.
\end{itemize}
\item Network Segmentation (\texttt{Enterprise M1030}; \texttt{ICS M0930})
\begin{itemize}
\item Divides networks into isolated segments or zones to restrict access, reduce lateral movement, and protect critical systems.
\end{itemize}
\item User Training (\texttt{Enterprise M1017}; \texttt{ICS M0917})
\begin{itemize}
\item Trains users to recognize and respond to phishing, social engineering, and other user-interaction threats.
\end{itemize}
\end{itemize}
\end{itemize}

\begin{itemize}
\item \textbf{AM-specific application:} qualify suppliers, verify firmware/software integrity, control maintenance and update channels, inspect incoming feedstock/components/equipment, and protect supplier-shared designs, recipes, and telemetry.
\end{itemize}

\begin{itemize}
\item \textbf{Residual coverage note:} ATT\&CK's supply-chain mitigation is intentionally broad. It does not replace supplier audits, material acceptance testing, contractual controls, provenance records, or regulatory/export-control due diligence.
\end{itemize}


\subsection{3.6.2 Feedstock compromise [I, A]}


Manipulation or substitution of AM feedstock, including chemical composition, powder characteristics, contamination, moisture or oxygen exposure, recycling, sieving, blending, or counterfeit material insertion.

Retained from Yampolskiy as a concrete supply-chain submethod.

Includes:

\begin{itemize}
\item contamination
\item chemical-composition manipulation
\item powder-characteristic manipulation
\item moisture/oxygen exposure manipulation
\item particle-size-distribution manipulation
\item poisoning during blending/sieving/recycling
\item counterfeit material substitution
\item NBC contamination
\end{itemize}

\textbf{MITRE ATT\&CK mitigation mapping (derived):}

\begin{itemize}
\item \textbf{Primary mitigations:}
\begin{itemize}
\item Supply Chain Management (\texttt{ICS M0817})
\begin{itemize}
\item Uses supplier policies, procedures, and integrity testing to ensure devices and components come from trusted sources.
\end{itemize}
\item Mechanical Protection Layers (\texttt{ICS M0805})
\begin{itemize}
\item Uses physical or mechanical protection layers, such as interlocks or safety devices, to prevent damage or harm.
\end{itemize}
\item Authorization Enforcement (\texttt{ICS M0800})
\begin{itemize}
\item Restricts read, manipulation, or execution privileges to authenticated users who require access under approved policies.
\end{itemize}
\item Access Management (\texttt{ICS M0801})
\begin{itemize}
\item Enforces access decisions, often through gateways or authentication services, before users can reach systems or devices.
\end{itemize}
\item Audit (\texttt{Enterprise M1047}; \texttt{ICS M0947})
\begin{itemize}
\item Records and reviews activity, permissions, configurations, and integrity indicators to identify weaknesses or anomalies.
\end{itemize}
\item Limit Hardware Installation (\texttt{Enterprise M1034}; \texttt{ICS M0934})
\begin{itemize}
\item Blocks users or groups from installing or using unapproved hardware, including removable devices.
\end{itemize}
\item Mitigation Limited or Not Effective (\texttt{ICS M0816})
\begin{itemize}
\item Marks techniques where preventive controls may be limited because the attack abuses legitimate system features.
\end{itemize}
\end{itemize}
\end{itemize}

\begin{itemize}
\item \textbf{Supporting mitigations:}
\begin{itemize}
\item Operational Information Confidentiality (\texttt{ICS M0809})
\begin{itemize}
\item Protects operational-process, facility, configuration, program, or database information that could reveal sensitive know-how or IP.
\end{itemize}
\item Data Loss Prevention (\texttt{Enterprise M1057}; \texttt{ICS M0803})
\begin{itemize}
\item Identifies, monitors, and controls sensitive-data movement to reduce unauthorized access, transmission, or exfiltration.
\end{itemize}
\item Network Segmentation (\texttt{Enterprise M1030}; \texttt{ICS M0930})
\begin{itemize}
\item Divides networks into isolated segments or zones to restrict access, reduce lateral movement, and protect critical systems.
\end{itemize}
\item Limit Access to Resource Over Network (\texttt{Enterprise M1035}; \texttt{ICS M0935})
\begin{itemize}
\item Restricts network access to resources such as file shares, remote systems, and services to legitimate users or systems.
\end{itemize}
\item Software Configuration (\texttt{ICS M0954}; \texttt{Enterprise M1054})
\begin{itemize}
\item Applies security-focused software setting changes to reduce risky behavior and attack surface.
\end{itemize}
\item Safety Instrumented Systems (\texttt{ICS M0812})
\begin{itemize}
\item Uses independent safety systems to detect hazardous conditions and automatically drive equipment to a safer state.
\end{itemize}
\end{itemize}
\end{itemize}

\begin{itemize}
\item \textbf{AM-specific application:} control supplier approval, material lots, powder recycling, sieving/blending records, storage conditions, authenticated material-handling systems, and access to feedstock containers or material-preparation equipment.
\end{itemize}

\begin{itemize}
\item \textbf{Residual coverage note:} chemical composition, moisture, oxygen exposure, particle-size distribution, and contamination verification depend on materials testing and chain-of-custody controls that ATT\&CK only covers indirectly.
\end{itemize}


\subsection{3.6.3 Replacement-parts compromise [I, A]}


Insertion of compromised spare parts, software, firmware, mechanical parts, electrical parts, or electronics through maintenance, repair, replacement, or upgrade channels.

Includes:

\begin{itemize}
\item malicious spare-part substitution
\item compromised upgrades
\item malicious firmware/software updates
\item software / firmware replacement-part compromise
\item mechanical replacement-part compromise
\item electrical replacement-part compromise
\item electronics replacement-part compromise
\end{itemize}

\textbf{MITRE ATT\&CK mitigation mapping (derived):}

\begin{itemize}
\item \textbf{Primary mitigations:}
\begin{itemize}
\item Supply Chain Management (\texttt{ICS M0817})
\begin{itemize}
\item Uses supplier policies, procedures, and integrity testing to ensure devices and components come from trusted sources.
\end{itemize}
\item Code Signing (\texttt{Enterprise M1045}; \texttt{ICS M0945})
\begin{itemize}
\item Verifies digital signatures so only trusted, untampered code is allowed to execute.
\end{itemize}
\item Boot Integrity (\texttt{Enterprise M1046}; \texttt{ICS M0946})
\begin{itemize}
\item Verifies the boot process, operating system, and loading mechanisms to prevent tampered components from starting trusted systems.
\end{itemize}
\item Update Software (\texttt{Enterprise M1051}; \texttt{ICS M0951})
\begin{itemize}
\item Applies vendor patches and upgrades to reduce exploitation of known software or firmware vulnerabilities.
\end{itemize}
\item Vulnerability Scanning (\texttt{Enterprise M1016}; \texttt{ICS M0916})
\begin{itemize}
\item Assesses systems, applications, and networks for vulnerabilities, misconfigurations, or unpatched software to prioritize remediation.
\end{itemize}
\item Limit Hardware Installation (\texttt{Enterprise M1034}; \texttt{ICS M0934})
\begin{itemize}
\item Blocks users or groups from installing or using unapproved hardware, including removable devices.
\end{itemize}
\item Limit Software Installation (\texttt{Enterprise M1033})
\begin{itemize}
\item Prevents installation of unauthorized or unapproved software to reduce malicious or vulnerable applications.
\end{itemize}
\end{itemize}
\end{itemize}

\begin{itemize}
\item \textbf{Supporting mitigations:}
\begin{itemize}
\item Application Developer Guidance (\texttt{Enterprise M1013}; \texttt{ICS M0913})
\begin{itemize}
\item Provides secure-development guidance to reduce weaknesses in applications, systems, and APIs.
\end{itemize}
\item Audit (\texttt{Enterprise M1047}; \texttt{ICS M0947})
\begin{itemize}
\item Records and reviews activity, permissions, configurations, and integrity indicators to identify weaknesses or anomalies.
\end{itemize}
\item Access Management (\texttt{ICS M0801})
\begin{itemize}
\item Enforces access decisions, often through gateways or authentication services, before users can reach systems or devices.
\end{itemize}
\item Authorization Enforcement (\texttt{ICS M0800})
\begin{itemize}
\item Restricts read, manipulation, or execution privileges to authenticated users who require access under approved policies.
\end{itemize}
\item Mechanical Protection Layers (\texttt{ICS M0805})
\begin{itemize}
\item Uses physical or mechanical protection layers, such as interlocks or safety devices, to prevent damage or harm.
\end{itemize}
\item Safety Instrumented Systems (\texttt{ICS M0812})
\begin{itemize}
\item Uses independent safety systems to detect hazardous conditions and automatically drive equipment to a safer state.
\end{itemize}
\item Software Configuration (\texttt{ICS M0954}; \texttt{Enterprise M1054})
\begin{itemize}
\item Applies security-focused software setting changes to reduce risky behavior and attack surface.
\end{itemize}
\end{itemize}
\end{itemize}

\begin{itemize}
\item \textbf{AM-specific application:} accept only trusted spare parts, firmware, software updates, electronics, sensors, actuators, and maintenance kits; verify digital signatures, supplier provenance, compatibility, and installation authorization.
\end{itemize}

\begin{itemize}
\item \textbf{Residual coverage note:} ATT\&CK covers software, firmware, and broad supply-chain integrity better than it covers mechanical-part quality. Physical acceptance inspection and equipment qualification remain necessary.
\end{itemize}


\subsection{3.6.4 AM equipment compromise before deployment [C, I, A]}


Delivery of AM machines, controllers, firmware, software, or hardware in a compromised state before the equipment is commissioned into production.

Includes:

\begin{itemize}
\item compromised machine delivery
\item compromised hardware/firmware/software from supplier channel
\item supply-chain compromise of AM equipment prior to deployment
\end{itemize}

\textbf{MITRE ATT\&CK mitigation mapping (derived):}

\begin{itemize}
\item \textbf{Primary mitigations:}
\begin{itemize}
\item Supply Chain Management (\texttt{ICS M0817})
\begin{itemize}
\item Uses supplier policies, procedures, and integrity testing to ensure devices and components come from trusted sources.
\end{itemize}
\item Boot Integrity (\texttt{Enterprise M1046}; \texttt{ICS M0946})
\begin{itemize}
\item Verifies the boot process, operating system, and loading mechanisms to prevent tampered components from starting trusted systems.
\end{itemize}
\item Code Signing (\texttt{Enterprise M1045}; \texttt{ICS M0945})
\begin{itemize}
\item Verifies digital signatures so only trusted, untampered code is allowed to execute.
\end{itemize}
\item Update Software (\texttt{Enterprise M1051}; \texttt{ICS M0951})
\begin{itemize}
\item Applies vendor patches and upgrades to reduce exploitation of known software or firmware vulnerabilities.
\end{itemize}
\item Vulnerability Scanning (\texttt{Enterprise M1016}; \texttt{ICS M0916})
\begin{itemize}
\item Assesses systems, applications, and networks for vulnerabilities, misconfigurations, or unpatched software to prioritize remediation.
\end{itemize}
\item Audit (\texttt{Enterprise M1047}; \texttt{ICS M0947})
\begin{itemize}
\item Records and reviews activity, permissions, configurations, and integrity indicators to identify weaknesses or anomalies.
\end{itemize}
\item Limit Hardware Installation (\texttt{Enterprise M1034}; \texttt{ICS M0934})
\begin{itemize}
\item Blocks users or groups from installing or using unapproved hardware, including removable devices.
\end{itemize}
\end{itemize}
\end{itemize}

\begin{itemize}
\item \textbf{Supporting mitigations:}
\begin{itemize}
\item Limit Software Installation (\texttt{Enterprise M1033})
\begin{itemize}
\item Prevents installation of unauthorized or unapproved software to reduce malicious or vulnerable applications.
\end{itemize}
\item Access Management (\texttt{ICS M0801})
\begin{itemize}
\item Enforces access decisions, often through gateways or authentication services, before users can reach systems or devices.
\end{itemize}
\item Authorization Enforcement (\texttt{ICS M0800})
\begin{itemize}
\item Restricts read, manipulation, or execution privileges to authenticated users who require access under approved policies.
\end{itemize}
\item Operational Information Confidentiality (\texttt{ICS M0809})
\begin{itemize}
\item Protects operational-process, facility, configuration, program, or database information that could reveal sensitive know-how or IP.
\end{itemize}
\item Software Configuration (\texttt{ICS M0954}; \texttt{Enterprise M1054})
\begin{itemize}
\item Applies security-focused software setting changes to reduce risky behavior and attack surface.
\end{itemize}
\item Operating System Configuration (\texttt{ICS M0928}; \texttt{Enterprise M1028})
\begin{itemize}
\item Hardens operating-system settings to reduce attack surface and opportunities for adversary abuse.
\end{itemize}
\item Mechanical Protection Layers (\texttt{ICS M0805})
\begin{itemize}
\item Uses physical or mechanical protection layers, such as interlocks or safety devices, to prevent damage or harm.
\end{itemize}
\item Safety Instrumented Systems (\texttt{ICS M0812})
\begin{itemize}
\item Uses independent safety systems to detect hazardous conditions and automatically drive equipment to a safer state.
\end{itemize}
\end{itemize}
\end{itemize}

\begin{itemize}
\item \textbf{AM-specific application:} verify delivered AM machines, controllers, HMIs, embedded firmware, calibration files, default accounts, installed software, service links, and safety components before connecting equipment to production networks.
\end{itemize}

\begin{itemize}
\item \textbf{Residual coverage note:} ATT\&CK cannot by itself prove that a delivered machine is mechanically or materially trustworthy; commissioning tests, factory acceptance testing, site acceptance testing, and supplier assurance are still required.
\end{itemize}


\subsection{3.6.5 Manufactured-object compromise in logistics [I, A]}


Replacement, contamination, tampering, or cyber infection of manufactured objects during shipment, warehousing, delivery, return, or other logistics handoffs.

Includes:

\begin{itemize}
\item replacement in transit
\item contamination in transit
\item insertion of cyber-infected embedded electronics
\item tampering during warehousing/delivery/return
\item manufactured-object replacement
\item manufactured-object cyber infection
\end{itemize}

\textbf{MITRE ATT\&CK mitigation mapping (derived):}

\begin{itemize}
\item \textbf{Primary mitigations:}
\begin{itemize}
\item Supply Chain Management (\texttt{ICS M0817})
\begin{itemize}
\item Uses supplier policies, procedures, and integrity testing to ensure devices and components come from trusted sources.
\end{itemize}
\item Audit (\texttt{Enterprise M1047}; \texttt{ICS M0947})
\begin{itemize}
\item Records and reviews activity, permissions, configurations, and integrity indicators to identify weaknesses or anomalies.
\end{itemize}
\item Access Management (\texttt{ICS M0801})
\begin{itemize}
\item Enforces access decisions, often through gateways or authentication services, before users can reach systems or devices.
\end{itemize}
\item Authorization Enforcement (\texttt{ICS M0800})
\begin{itemize}
\item Restricts read, manipulation, or execution privileges to authenticated users who require access under approved policies.
\end{itemize}
\item Mechanical Protection Layers (\texttt{ICS M0805})
\begin{itemize}
\item Uses physical or mechanical protection layers, such as interlocks or safety devices, to prevent damage or harm.
\end{itemize}
\item Mitigation Limited or Not Effective (\texttt{ICS M0816})
\begin{itemize}
\item Marks techniques where preventive controls may be limited because the attack abuses legitimate system features.
\end{itemize}
\end{itemize}
\end{itemize}

\begin{itemize}
\item \textbf{Supporting mitigations:}
\begin{itemize}
\item Operational Information Confidentiality (\texttt{ICS M0809})
\begin{itemize}
\item Protects operational-process, facility, configuration, program, or database information that could reveal sensitive know-how or IP.
\end{itemize}
\item Data Loss Prevention (\texttt{Enterprise M1057}; \texttt{ICS M0803})
\begin{itemize}
\item Identifies, monitors, and controls sensitive-data movement to reduce unauthorized access, transmission, or exfiltration.
\end{itemize}
\item Code Signing (\texttt{Enterprise M1045}; \texttt{ICS M0945})
\begin{itemize}
\item Verifies digital signatures so only trusted, untampered code is allowed to execute.
\end{itemize}
\item Boot Integrity (\texttt{Enterprise M1046}; \texttt{ICS M0946})
\begin{itemize}
\item Verifies the boot process, operating system, and loading mechanisms to prevent tampered components from starting trusted systems.
\end{itemize}
\item Update Software (\texttt{Enterprise M1051}; \texttt{ICS M0951})
\begin{itemize}
\item Applies vendor patches and upgrades to reduce exploitation of known software or firmware vulnerabilities.
\end{itemize}
\item Network Segmentation (\texttt{Enterprise M1030}; \texttt{ICS M0930})
\begin{itemize}
\item Divides networks into isolated segments or zones to restrict access, reduce lateral movement, and protect critical systems.
\end{itemize}
\item Data Backup (\texttt{ICS M0953}; \texttt{Enterprise M1053})
\begin{itemize}
\item Creates and protects backups of critical data and configurations so systems can recover after compromise or disruption.
\end{itemize}
\end{itemize}
\end{itemize}

\begin{itemize}
\item \textbf{AM-specific application:} preserve chain of custody for shipped parts, returned parts, embedded electronics, serialized components, and warehoused inventory; verify any cyber-capable embedded element before use.
\end{itemize}

\begin{itemize}
\item \textbf{Residual coverage note:} physical replacement, contamination, or tampering in transit is only partially represented in ATT\&CK. Tamper-evident packaging, physical inspection, logistics controls, and acceptance testing are outside the detailed ATT\&CK mitigation catalog.
\end{itemize}


\subsection{3.6.6 Disrupting supply-chain logistics [C, I, A]}


Operational disruption of AM supply-chain logistics by manipulating shipments, routing, tracking, distribution timing, or part custody to delay production or swap or modify items.

Retained from Kumar as a distinct operational method.

Includes:

\begin{itemize}
\item shipment disruption
\item tracking-system manipulation
\item delay induction
\item routing manipulation
\item part swapping during distribution
\end{itemize}

\textbf{MITRE ATT\&CK mitigation mapping (derived):}

\begin{itemize}
\item \textbf{Primary mitigations:}
\begin{itemize}
\item Supply Chain Management (\texttt{ICS M0817})
\begin{itemize}
\item Uses supplier policies, procedures, and integrity testing to ensure devices and components come from trusted sources.
\end{itemize}
\item Redundancy of Service (\texttt{ICS M0811})
\begin{itemize}
\item Provides backup or standby devices and services so critical functions can continue during failures or disruption.
\end{itemize}
\item Out-of-Band Communications Channel (\texttt{ICS M0810}; \texttt{Enterprise M1060})
\begin{itemize}
\item Provides alternative communication paths when primary network communication fails, is disrupted, or cannot be trusted.
\end{itemize}
\item Data Backup (\texttt{ICS M0953}; \texttt{Enterprise M1053})
\begin{itemize}
\item Creates and protects backups of critical data and configurations so systems can recover after compromise or disruption.
\end{itemize}
\item Network Segmentation (\texttt{Enterprise M1030}; \texttt{ICS M0930})
\begin{itemize}
\item Divides networks into isolated segments or zones to restrict access, reduce lateral movement, and protect critical systems.
\end{itemize}
\item Limit Access to Resource Over Network (\texttt{Enterprise M1035}; \texttt{ICS M0935})
\begin{itemize}
\item Restricts network access to resources such as file shares, remote systems, and services to legitimate users or systems.
\end{itemize}
\end{itemize}
\end{itemize}

\begin{itemize}
\item \textbf{Supporting mitigations:}
\begin{itemize}
\item Filter Network Traffic (\texttt{ICS M0937}; \texttt{Enterprise M1037})
\begin{itemize}
\item Filters ingress, egress, or lateral traffic, including protocol-aware rules, to restrict malicious or unnecessary communication.
\end{itemize}
\item Network Intrusion Prevention (\texttt{ICS M0931}; \texttt{Enterprise M1031})
\begin{itemize}
\item Uses intrusion-detection signatures to block traffic at network boundaries, with care not to disrupt control or safety functions.
\end{itemize}
\item Access Management (\texttt{ICS M0801})
\begin{itemize}
\item Enforces access decisions, often through gateways or authentication services, before users can reach systems or devices.
\end{itemize}
\item Authorization Enforcement (\texttt{ICS M0800})
\begin{itemize}
\item Restricts read, manipulation, or execution privileges to authenticated users who require access under approved policies.
\end{itemize}
\item Audit (\texttt{Enterprise M1047}; \texttt{ICS M0947})
\begin{itemize}
\item Records and reviews activity, permissions, configurations, and integrity indicators to identify weaknesses or anomalies.
\end{itemize}
\item Threat Intelligence Program (\texttt{Enterprise M1019}; \texttt{ICS M0919})
\begin{itemize}
\item Collects and analyzes threat intelligence to track trends, guide defensive priorities, and reduce risk.
\end{itemize}
\item User Training (\texttt{Enterprise M1017}; \texttt{ICS M0917})
\begin{itemize}
\item Trains users to recognize and respond to phishing, social engineering, and other user-interaction threats.
\end{itemize}
\item Update Software (\texttt{Enterprise M1051}; \texttt{ICS M0951})
\begin{itemize}
\item Applies vendor patches and upgrades to reduce exploitation of known software or firmware vulnerabilities.
\end{itemize}
\end{itemize}
\end{itemize}

\begin{itemize}
\item \textbf{AM-specific application:} protect shipment tracking, supplier portals, logistics APIs, production scheduling, inventory systems, emergency supplier communications, and fallback delivery or production-routing plans.
\end{itemize}

\begin{itemize}
\item \textbf{Residual coverage note:} ATT\&CK helps for cyber-enabled logistics disruption, but carrier failure, geopolitical disruption, strikes, and physical routing constraints require business-continuity and supplier-contingency planning beyond ATT\&CK.
\end{itemize}


\section{3.7 Equipment-compromise methods [C, I, A]}


Methods in this family compromise AM printers, controllers, firmware, software, sensors, actuators, safety components, or machine-control infrastructure, using the equipment itself as a path to theft, tampering, or disruption.


\subsection{3.7.1 Cyber intrusion on AM equipment [C, I, A]}


Cyber compromise of AM equipment or its control environment through unauthorised access, malware, ransomware, firmware or software compromise, or network intrusion.

Merged with Yampolskiy's compromised-equipment logic.

Includes:

\begin{itemize}
\item unauthorised remote access
\item malware injection
\item ransomware
\item firmware compromise
\item software compromise
\item network compromise of machine-control environment
\end{itemize}

\textbf{MITRE ATT\&CK mitigation mapping (derived):}

\begin{itemize}
\item \textbf{Primary mitigations:}
\begin{itemize}
\item Network Segmentation (\texttt{Enterprise M1030}; \texttt{ICS M0930})
\begin{itemize}
\item Divides networks into isolated segments or zones to restrict access, reduce lateral movement, and protect critical systems.
\end{itemize}
\item Limit Access to Resource Over Network (\texttt{Enterprise M1035}; \texttt{ICS M0935})
\begin{itemize}
\item Restricts network access to resources such as file shares, remote systems, and services to legitimate users or systems.
\end{itemize}
\item Filter Network Traffic (\texttt{Enterprise M1037}; \texttt{ICS M0937})
\begin{itemize}
\item Filters ingress, egress, or lateral traffic, including protocol-aware rules, to restrict malicious or unnecessary communication.
\end{itemize}
\item Network Allowlists (\texttt{ICS M0807})
\begin{itemize}
\item Allows only approved device connections based on defined addresses, ports, protocols, or related network properties.
\end{itemize}
\item Access Management (\texttt{ICS M0801})
\begin{itemize}
\item Enforces access decisions, often through gateways or authentication services, before users can reach systems or devices.
\end{itemize}
\item Authorization Enforcement (\texttt{ICS M0800})
\begin{itemize}
\item Restricts read, manipulation, or execution privileges to authenticated users who require access under approved policies.
\end{itemize}
\item Human User Authentication (\texttt{ICS M0804})
\begin{itemize}
\item Requires users to authenticate before accessing data or sending commands to a device or system.
\end{itemize}
\item Multi-factor Authentication (\texttt{Enterprise M1032}; \texttt{ICS M0932})
\begin{itemize}
\item Requires two or more forms of identity evidence before granting access to a system.
\end{itemize}
\item Privileged Account Management (\texttt{Enterprise M1026}; \texttt{ICS M0926})
\begin{itemize}
\item Controls creation, use, permissions, and accountability for administrative, root, SYSTEM, or other privileged accounts.
\end{itemize}
\item User Account Management (\texttt{Enterprise M1018}; \texttt{ICS M0918})
\begin{itemize}
\item Manages account creation, modification, use, and permissions to reduce unauthorized access and privilege misuse.
\end{itemize}
\item Account Use Policies (\texttt{Enterprise M1036}; \texttt{ICS M0936})
\begin{itemize}
\item Configures account-use rules such as lockouts, login times, and inactivity timeouts to reduce unauthorized access.
\end{itemize}
\end{itemize}
\end{itemize}

\begin{itemize}
\item \textbf{Supporting mitigations:}
\begin{itemize}
\item Antivirus/Antimalware (\texttt{Enterprise M1049}; \texttt{ICS M0949})
\begin{itemize}
\item Uses signatures, heuristics, or behavior analysis to detect, block, and remediate malicious software.
\end{itemize}
\item Behavior Prevention on Endpoint (\texttt{Enterprise M1040})
\begin{itemize}
\item Detects and blocks suspicious endpoint behavior based on process, file, API, or other activity patterns.
\end{itemize}
\item Execution Prevention (\texttt{Enterprise M1038}; \texttt{ICS M0938})
\begin{itemize}
\item Blocks unauthorized or malicious code execution through application control, script blocking, or related controls.
\end{itemize}
\item Application Isolation and Sandboxing (\texttt{Enterprise M1048}; \texttt{ICS M0948})
\begin{itemize}
\item Runs code in a controlled, isolated environment to limit access to sensitive resources and reduce impact.
\end{itemize}
\item Code Signing (\texttt{Enterprise M1045}; \texttt{ICS M0945})
\begin{itemize}
\item Verifies digital signatures so only trusted, untampered code is allowed to execute.
\end{itemize}
\item Boot Integrity (\texttt{Enterprise M1046}; \texttt{ICS M0946})
\begin{itemize}
\item Verifies the boot process, operating system, and loading mechanisms to prevent tampered components from starting trusted systems.
\end{itemize}
\item Exploit Protection (\texttt{Enterprise M1050}; \texttt{ICS M0950})
\begin{itemize}
\item Detects, blocks, or mitigates conditions that indicate or enable software exploitation.
\end{itemize}
\item Update Software (\texttt{Enterprise M1051}; \texttt{ICS M0951})
\begin{itemize}
\item Applies vendor patches and upgrades to reduce exploitation of known software or firmware vulnerabilities.
\end{itemize}
\item Vulnerability Scanning (\texttt{Enterprise M1016}; \texttt{ICS M0916})
\begin{itemize}
\item Assesses systems, applications, and networks for vulnerabilities, misconfigurations, or unpatched software to prioritize remediation.
\end{itemize}
\item Operating System Configuration (\texttt{Enterprise M1028}; \texttt{ICS M0928})
\begin{itemize}
\item Hardens operating-system settings to reduce attack surface and opportunities for adversary abuse.
\end{itemize}
\item Software Configuration (\texttt{Enterprise M1054}; \texttt{ICS M0954})
\begin{itemize}
\item Applies security-focused software setting changes to reduce risky behavior and attack surface.
\end{itemize}
\item Restrict File and Directory Permissions (\texttt{Enterprise M1022}; \texttt{ICS M0922})
\begin{itemize}
\item Limits file and directory access so only authorized users, groups, or processes can read, write, or execute them.
\end{itemize}
\item Restrict Library Loading (\texttt{Enterprise M1044}; \texttt{ICS M0944})
\begin{itemize}
\item Ensures only trusted libraries are loaded into processes to prevent abuse of library-loading mechanisms.
\end{itemize}
\item Restrict Registry Permissions (\texttt{Enterprise M1024}; \texttt{ICS M0924})
\begin{itemize}
\item Limits who can modify sensitive Windows registry keys or hives to reduce persistence, privilege escalation, or evasion.
\end{itemize}
\item Audit (\texttt{Enterprise M1047}; \texttt{ICS M0947})
\begin{itemize}
\item Records and reviews activity, permissions, configurations, and integrity indicators to identify weaknesses or anomalies.
\end{itemize}
\item Data Backup (\texttt{Enterprise M1053}; \texttt{ICS M0953})
\begin{itemize}
\item Creates and protects backups of critical data and configurations so systems can recover after compromise or disruption.
\end{itemize}
\item Remote Data Storage (\texttt{Enterprise M1029})
\begin{itemize}
\item Moves critical logs or sensitive files to secure off-host storage to reduce tampering, destruction, or unauthorized access.
\end{itemize}
\item Threat Intelligence Program (\texttt{Enterprise M1019}; \texttt{ICS M0919})
\begin{itemize}
\item Collects and analyzes threat intelligence to track trends, guide defensive priorities, and reduce risk.
\end{itemize}
\end{itemize}
\end{itemize}

\begin{itemize}
\item \textbf{AM-specific application:} harden printers, controllers, HMIs, slicer hosts, build servers, remote maintenance links, firmware-update paths, engineering workstations, quality systems, and production networks against unauthorized remote access, malware, ransomware, and machine-control compromise.
\end{itemize}

\begin{itemize}
\item \textbf{Residual coverage note:} some endpoint controls can disrupt real-time industrial operation if deployed carelessly. Validate them in representative AM test environments before production rollout.
\end{itemize}


\subsection{3.7.2 Physical sabotage on AM equipment [I, A]}


Physical interference with AM hardware, components, operating conditions, or safety-critical elements to obstruct operation, cause malfunction, introduce defects, or damage equipment.

Includes:

\begin{itemize}
\item tampering with sensors/motors/actuators
\item insertion of contaminants/debris
\item mechanical overload
\item physical damage/vandalism
\item unsafe modification of safety-critical components
\end{itemize}

\textbf{MITRE ATT\&CK mitigation mapping (derived):}

\begin{itemize}
\item \textbf{Primary mitigations:}
\begin{itemize}
\item Mechanical Protection Layers (\texttt{ICS M0805})
\begin{itemize}
\item Uses physical or mechanical protection layers, such as interlocks or safety devices, to prevent damage or harm.
\end{itemize}
\item Safety Instrumented Systems (\texttt{ICS M0812})
\begin{itemize}
\item Uses independent safety systems to detect hazardous conditions and automatically drive equipment to a safer state.
\end{itemize}
\item Access Management (\texttt{ICS M0801})
\begin{itemize}
\item Enforces access decisions, often through gateways or authentication services, before users can reach systems or devices.
\end{itemize}
\item Authorization Enforcement (\texttt{ICS M0800})
\begin{itemize}
\item Restricts read, manipulation, or execution privileges to authenticated users who require access under approved policies.
\end{itemize}
\item Limit Hardware Installation (\texttt{Enterprise M1034}; \texttt{ICS M0934})
\begin{itemize}
\item Blocks users or groups from installing or using unapproved hardware, including removable devices.
\end{itemize}
\item Audit (\texttt{Enterprise M1047}; \texttt{ICS M0947})
\begin{itemize}
\item Records and reviews activity, permissions, configurations, and integrity indicators to identify weaknesses or anomalies.
\end{itemize}
\item Mitigation Limited or Not Effective (\texttt{ICS M0816})
\begin{itemize}
\item Marks techniques where preventive controls may be limited because the attack abuses legitimate system features.
\end{itemize}
\end{itemize}
\end{itemize}

\begin{itemize}
\item \textbf{Supporting mitigations:}
\begin{itemize}
\item Watchdog Timers (\texttt{ICS M0815})
\begin{itemize}
\item Uses timers to quickly detect when a device or system becomes unresponsive.
\end{itemize}
\item Redundancy of Service (\texttt{ICS M0811})
\begin{itemize}
\item Provides backup or standby devices and services so critical functions can continue during failures or disruption.
\end{itemize}
\item Out-of-Band Communications Channel (\texttt{ICS M0810}; \texttt{Enterprise M1060})
\begin{itemize}
\item Provides alternative communication paths when primary network communication fails, is disrupted, or cannot be trusted.
\end{itemize}
\item User Training (\texttt{Enterprise M1017}; \texttt{ICS M0917})
\begin{itemize}
\item Trains users to recognize and respond to phishing, social engineering, and other user-interaction threats.
\end{itemize}
\item Data Backup (\texttt{ICS M0953}; \texttt{Enterprise M1053})
\begin{itemize}
\item Creates and protects backups of critical data and configurations so systems can recover after compromise or disruption.
\end{itemize}
\item Network Segmentation (\texttt{Enterprise M1030}; \texttt{ICS M0930})
\begin{itemize}
\item Divides networks into isolated segments or zones to restrict access, reduce lateral movement, and protect critical systems.
\end{itemize}
\item Software Configuration (\texttt{ICS M0954}; \texttt{Enterprise M1054})
\begin{itemize}
\item Applies security-focused software setting changes to reduce risky behavior and attack surface.
\end{itemize}
\end{itemize}
\end{itemize}

\begin{itemize}
\item \textbf{AM-specific application:} protect physical access to printers, material bays, sensors, actuators, safety interlocks, electrical cabinets, build chambers, post-processing equipment, and maintenance ports.
\end{itemize}

\begin{itemize}
\item \textbf{Residual coverage note:} ATT\&CK is primarily cyber-focused. Badge controls, surveillance, guards, locked rooms, tamper-evident seals, and facility security are necessary but not fully specified as ATT\&CK mitigations.
\end{itemize}


\subsection{3.7.3 Manipulating 3D equipment [C, I, A]}


Unauthorised control or reconfiguration of AM machinery so attackers can change parameters, bypass safeguards, manipulate the print process, introduce defects, steal data, or disrupt production.

Retained from Kumar as a practical umbrella label.

Includes:

\begin{itemize}
\item parameter manipulation via compromised machine
\item bypassing safety checks
\item malicious machine reconfiguration
\item use of compromised equipment to affect part quality or output quantity
\end{itemize}

\textbf{MITRE ATT\&CK mitigation mapping (derived):}

\begin{itemize}
\item \textbf{Primary mitigations:}
\begin{itemize}
\item Authorization Enforcement (\texttt{ICS M0800})
\begin{itemize}
\item Restricts read, manipulation, or execution privileges to authenticated users who require access under approved policies.
\end{itemize}
\item Access Management (\texttt{ICS M0801})
\begin{itemize}
\item Enforces access decisions, often through gateways or authentication services, before users can reach systems or devices.
\end{itemize}
\item Human User Authentication (\texttt{ICS M0804})
\begin{itemize}
\item Requires users to authenticate before accessing data or sending commands to a device or system.
\end{itemize}
\item Multi-factor Authentication (\texttt{Enterprise M1032}; \texttt{ICS M0932})
\begin{itemize}
\item Requires two or more forms of identity evidence before granting access to a system.
\end{itemize}
\item Privileged Account Management (\texttt{Enterprise M1026}; \texttt{ICS M0926})
\begin{itemize}
\item Controls creation, use, permissions, and accountability for administrative, root, SYSTEM, or other privileged accounts.
\end{itemize}
\item Validate Program Inputs (\texttt{ICS M0818})
\begin{itemize}
\item Checks input and action values against expected format, content, and operational ranges before accepting them.
\end{itemize}
\item Communication Authenticity (\texttt{ICS M0802})
\begin{itemize}
\item Authenticates message senders and verifies message integrity to detect spoofed messages or unauthorized connections.
\end{itemize}
\item Software Process and Device Authentication (\texttt{ICS M0813})
\begin{itemize}
\item Requires devices and software processes to authenticate, especially for remote connections and API access.
\end{itemize}
\item Code Signing (\texttt{Enterprise M1045}; \texttt{ICS M0945})
\begin{itemize}
\item Verifies digital signatures so only trusted, untampered code is allowed to execute.
\end{itemize}
\item Boot Integrity (\texttt{Enterprise M1046}; \texttt{ICS M0946})
\begin{itemize}
\item Verifies the boot process, operating system, and loading mechanisms to prevent tampered components from starting trusted systems.
\end{itemize}
\end{itemize}
\end{itemize}

\begin{itemize}
\item \textbf{Supporting mitigations:}
\begin{itemize}
\item Network Segmentation (\texttt{Enterprise M1030}; \texttt{ICS M0930})
\begin{itemize}
\item Divides networks into isolated segments or zones to restrict access, reduce lateral movement, and protect critical systems.
\end{itemize}
\item Network Allowlists (\texttt{ICS M0807})
\begin{itemize}
\item Allows only approved device connections based on defined addresses, ports, protocols, or related network properties.
\end{itemize}
\item Filter Network Traffic (\texttt{ICS M0937}; \texttt{Enterprise M1037})
\begin{itemize}
\item Filters ingress, egress, or lateral traffic, including protocol-aware rules, to restrict malicious or unnecessary communication.
\end{itemize}
\item Static Network Configuration (\texttt{ICS M0814})
\begin{itemize}
\item Uses static network settings where possible to reduce abuse of dynamic discovery or addressing protocols.
\end{itemize}
\item Software Configuration (\texttt{ICS M0954}; \texttt{Enterprise M1054})
\begin{itemize}
\item Applies security-focused software setting changes to reduce risky behavior and attack surface.
\end{itemize}
\item Audit (\texttt{Enterprise M1047}; \texttt{ICS M0947})
\begin{itemize}
\item Records and reviews activity, permissions, configurations, and integrity indicators to identify weaknesses or anomalies.
\end{itemize}
\item Mechanical Protection Layers (\texttt{ICS M0805})
\begin{itemize}
\item Uses physical or mechanical protection layers, such as interlocks or safety devices, to prevent damage or harm.
\end{itemize}
\item Safety Instrumented Systems (\texttt{ICS M0812})
\begin{itemize}
\item Uses independent safety systems to detect hazardous conditions and automatically drive equipment to a safer state.
\end{itemize}
\item Watchdog Timers (\texttt{ICS M0815})
\begin{itemize}
\item Uses timers to quickly detect when a device or system becomes unresponsive.
\end{itemize}
\item Redundancy of Service (\texttt{ICS M0811})
\begin{itemize}
\item Provides backup or standby devices and services so critical functions can continue during failures or disruption.
\end{itemize}
\end{itemize}
\end{itemize}

\begin{itemize}
\item \textbf{AM-specific application:} control and log machine reconfiguration, safety-check bypass attempts, production-count changes, parameter edits, calibration changes, maintenance-mode access, and machine-profile updates.
\end{itemize}

\begin{itemize}
\item \textbf{Residual coverage note:} security controls can prevent unauthorized manipulation, but detecting all malicious-but-plausible equipment settings requires AM process expertise and quality assurance.
\end{itemize}


\subsection{3.7.4 Compromised AM equipment as a method-enabler [C, I, A]}


Use of already compromised AM equipment as a trusted attack platform that can alter models, process parameters, post-processing, quality control, data access, or power behavior.

Retained from Yampolskiy as a special note: compromised equipment is not just a target; it also enables other methods.

It can enable:

\begin{itemize}
\item part-specification tampering
\item process manipulation
\item post-processing manipulation
\item data theft
\item quality-control bypass
\item environmental harm
\item cyber attacks
\item power surge
\end{itemize}

\textbf{MITRE ATT\&CK mitigation mapping (derived):}

\begin{itemize}
\item \textbf{Primary mitigations:}
\begin{itemize}
\item Boot Integrity (\texttt{Enterprise M1046}; \texttt{ICS M0946})
\begin{itemize}
\item Verifies the boot process, operating system, and loading mechanisms to prevent tampered components from starting trusted systems.
\end{itemize}
\item Code Signing (\texttt{Enterprise M1045}; \texttt{ICS M0945})
\begin{itemize}
\item Verifies digital signatures so only trusted, untampered code is allowed to execute.
\end{itemize}
\item Update Software (\texttt{Enterprise M1051}; \texttt{ICS M0951})
\begin{itemize}
\item Applies vendor patches and upgrades to reduce exploitation of known software or firmware vulnerabilities.
\end{itemize}
\item Vulnerability Scanning (\texttt{Enterprise M1016}; \texttt{ICS M0916})
\begin{itemize}
\item Assesses systems, applications, and networks for vulnerabilities, misconfigurations, or unpatched software to prioritize remediation.
\end{itemize}
\item Execution Prevention (\texttt{Enterprise M1038}; \texttt{ICS M0938})
\begin{itemize}
\item Blocks unauthorized or malicious code execution through application control, script blocking, or related controls.
\end{itemize}
\item Antivirus/Antimalware (\texttt{Enterprise M1049}; \texttt{ICS M0949})
\begin{itemize}
\item Uses signatures, heuristics, or behavior analysis to detect, block, and remediate malicious software.
\end{itemize}
\item Application Isolation and Sandboxing (\texttt{Enterprise M1048}; \texttt{ICS M0948})
\begin{itemize}
\item Runs code in a controlled, isolated environment to limit access to sensitive resources and reduce impact.
\end{itemize}
\item Access Management (\texttt{ICS M0801})
\begin{itemize}
\item Enforces access decisions, often through gateways or authentication services, before users can reach systems or devices.
\end{itemize}
\item Authorization Enforcement (\texttt{ICS M0800})
\begin{itemize}
\item Restricts read, manipulation, or execution privileges to authenticated users who require access under approved policies.
\end{itemize}
\item Network Segmentation (\texttt{Enterprise M1030}; \texttt{ICS M0930})
\begin{itemize}
\item Divides networks into isolated segments or zones to restrict access, reduce lateral movement, and protect critical systems.
\end{itemize}
\end{itemize}
\end{itemize}

\begin{itemize}
\item \textbf{Supporting mitigations:}
\begin{itemize}
\item Limit Access to Resource Over Network (\texttt{Enterprise M1035}; \texttt{ICS M0935})
\begin{itemize}
\item Restricts network access to resources such as file shares, remote systems, and services to legitimate users or systems.
\end{itemize}
\item Filter Network Traffic (\texttt{Enterprise M1037}; \texttt{ICS M0937})
\begin{itemize}
\item Filters ingress, egress, or lateral traffic, including protocol-aware rules, to restrict malicious or unnecessary communication.
\end{itemize}
\item Network Allowlists (\texttt{ICS M0807})
\begin{itemize}
\item Allows only approved device connections based on defined addresses, ports, protocols, or related network properties.
\end{itemize}
\item Software Process and Device Authentication (\texttt{ICS M0813})
\begin{itemize}
\item Requires devices and software processes to authenticate, especially for remote connections and API access.
\end{itemize}
\item Audit (\texttt{Enterprise M1047}; \texttt{ICS M0947})
\begin{itemize}
\item Records and reviews activity, permissions, configurations, and integrity indicators to identify weaknesses or anomalies.
\end{itemize}
\item Remote Data Storage (\texttt{Enterprise M1029})
\begin{itemize}
\item Moves critical logs or sensitive files to secure off-host storage to reduce tampering, destruction, or unauthorized access.
\end{itemize}
\item Data Backup (\texttt{ICS M0953}; \texttt{Enterprise M1053})
\begin{itemize}
\item Creates and protects backups of critical data and configurations so systems can recover after compromise or disruption.
\end{itemize}
\item Mechanical Protection Layers (\texttt{ICS M0805})
\begin{itemize}
\item Uses physical or mechanical protection layers, such as interlocks or safety devices, to prevent damage or harm.
\end{itemize}
\item Safety Instrumented Systems (\texttt{ICS M0812})
\begin{itemize}
\item Uses independent safety systems to detect hazardous conditions and automatically drive equipment to a safer state.
\end{itemize}
\item Watchdog Timers (\texttt{ICS M0815})
\begin{itemize}
\item Uses timers to quickly detect when a device or system becomes unresponsive.
\end{itemize}
\item Out-of-Band Communications Channel (\texttt{ICS M0810}; \texttt{Enterprise M1060})
\begin{itemize}
\item Provides alternative communication paths when primary network communication fails, is disrupted, or cannot be trusted.
\end{itemize}
\end{itemize}
\end{itemize}

\begin{itemize}
\item \textbf{AM-specific application:} treat compromised equipment as both a target and a platform for downstream AM attacks; harden the machine enough that it cannot silently become the trusted path for theft, tampering, unsafe process changes, quality-record falsification, or sabotage.
\end{itemize}

\begin{itemize}
\item \textbf{Residual coverage note:} this node is an enabler rather than one discrete technique. In practice, combine these baseline equipment mitigations with the mitigation block for the specific enabled method.
\end{itemize}


\section{3.8 Part-compromise methods [I, A]}


Methods in this family directly compromise the produced part or make the part itself a carrier of hidden weakness, contamination, embedded components, or downstream attack capability.


\subsection{3.8.1 Compromising 3D-printed parts [I]}


Production or introduction of intentionally altered, weakened, or defective 3D-printed parts through design manipulation, material substitution, hidden flaws, parameter tampering, scaling errors, or defect concealment.

Merged from Kumar and Yampolskiy's sabotage logic.

Includes:

\begin{itemize}
\item introducing hidden flaws
\item introducing voids or microfractures
\item inducing weakness through parameter changes
\item incorrect scaling or tolerances
\item weakening via material substitution
\item concealing defects
\end{itemize}

\textbf{MITRE ATT\&CK mitigation mapping (derived):}

\begin{itemize}
\item \textbf{Primary mitigations:}
\begin{itemize}
\item Validate Program Inputs (\texttt{ICS M0818})
\begin{itemize}
\item Checks input and action values against expected format, content, and operational ranges before accepting them.
\end{itemize}
\item Authorization Enforcement (\texttt{ICS M0800})
\begin{itemize}
\item Restricts read, manipulation, or execution privileges to authenticated users who require access under approved policies.
\end{itemize}
\item Access Management (\texttt{ICS M0801})
\begin{itemize}
\item Enforces access decisions, often through gateways or authentication services, before users can reach systems or devices.
\end{itemize}
\item Code Signing (\texttt{Enterprise M1045}; \texttt{ICS M0945})
\begin{itemize}
\item Verifies digital signatures so only trusted, untampered code is allowed to execute.
\end{itemize}
\item Communication Authenticity (\texttt{ICS M0802})
\begin{itemize}
\item Authenticates message senders and verifies message integrity to detect spoofed messages or unauthorized connections.
\end{itemize}
\item Software Process and Device Authentication (\texttt{ICS M0813})
\begin{itemize}
\item Requires devices and software processes to authenticate, especially for remote connections and API access.
\end{itemize}
\item Audit (\texttt{Enterprise M1047}; \texttt{ICS M0947})
\begin{itemize}
\item Records and reviews activity, permissions, configurations, and integrity indicators to identify weaknesses or anomalies.
\end{itemize}
\item Supply Chain Management (\texttt{ICS M0817})
\begin{itemize}
\item Uses supplier policies, procedures, and integrity testing to ensure devices and components come from trusted sources.
\end{itemize}
\end{itemize}
\end{itemize}

\begin{itemize}
\item \textbf{Supporting mitigations:}
\begin{itemize}
\item Network Segmentation (\texttt{Enterprise M1030}; \texttt{ICS M0930})
\begin{itemize}
\item Divides networks into isolated segments or zones to restrict access, reduce lateral movement, and protect critical systems.
\end{itemize}
\item Network Allowlists (\texttt{ICS M0807})
\begin{itemize}
\item Allows only approved device connections based on defined addresses, ports, protocols, or related network properties.
\end{itemize}
\item Filter Network Traffic (\texttt{ICS M0937}; \texttt{Enterprise M1037})
\begin{itemize}
\item Filters ingress, egress, or lateral traffic, including protocol-aware rules, to restrict malicious or unnecessary communication.
\end{itemize}
\item Software Configuration (\texttt{ICS M0954}; \texttt{Enterprise M1054})
\begin{itemize}
\item Applies security-focused software setting changes to reduce risky behavior and attack surface.
\end{itemize}
\item Data Backup (\texttt{ICS M0953}; \texttt{Enterprise M1053})
\begin{itemize}
\item Creates and protects backups of critical data and configurations so systems can recover after compromise or disruption.
\end{itemize}
\item Remote Data Storage (\texttt{Enterprise M1029})
\begin{itemize}
\item Moves critical logs or sensitive files to secure off-host storage to reduce tampering, destruction, or unauthorized access.
\end{itemize}
\item Mechanical Protection Layers (\texttt{ICS M0805})
\begin{itemize}
\item Uses physical or mechanical protection layers, such as interlocks or safety devices, to prevent damage or harm.
\end{itemize}
\item Safety Instrumented Systems (\texttt{ICS M0812})
\begin{itemize}
\item Uses independent safety systems to detect hazardous conditions and automatically drive equipment to a safer state.
\end{itemize}
\end{itemize}
\end{itemize}

\begin{itemize}
\item \textbf{AM-specific application:} protect every upstream input that determines final-part integrity: design baseline, material lot, machine profile, tool path, build parameters, post-processing recipe, inspection data, and release decision.
\end{itemize}

\begin{itemize}
\item \textbf{Residual coverage note:} ATT\&CK reduces manipulation opportunities, but it does not replace NDE, destructive testing, dimensional inspection, fatigue testing, material certification, or safety-case evidence for critical AM parts.
\end{itemize}


\subsection{3.8.2 Contamination and embedded-component insertion [I, A]}


Insertion of contaminants, foreign materials, NBC substances, or active or passive embedded components into parts so the manufactured object can harm users, systems, equipment, or its environment.

Retained because it appears explicitly in Yampolskiy's sabotage treatment and fits modern downstream risk scenarios.

Includes:

\begin{itemize}
\item NBC contamination
\item conductive/active/passive embedded components
\item hidden foreign material insertion
\item malicious carriers embedded into finished parts
\item compromised part / compromised 3D object behavior
\end{itemize}

\textbf{MITRE ATT\&CK mitigation mapping (derived):}

\begin{itemize}
\item \textbf{Primary mitigations:}
\begin{itemize}
\item Supply Chain Management (\texttt{ICS M0817})
\begin{itemize}
\item Uses supplier policies, procedures, and integrity testing to ensure devices and components come from trusted sources.
\end{itemize}
\item Mechanical Protection Layers (\texttt{ICS M0805})
\begin{itemize}
\item Uses physical or mechanical protection layers, such as interlocks or safety devices, to prevent damage or harm.
\end{itemize}
\item Authorization Enforcement (\texttt{ICS M0800})
\begin{itemize}
\item Restricts read, manipulation, or execution privileges to authenticated users who require access under approved policies.
\end{itemize}
\item Access Management (\texttt{ICS M0801})
\begin{itemize}
\item Enforces access decisions, often through gateways or authentication services, before users can reach systems or devices.
\end{itemize}
\item Limit Hardware Installation (\texttt{Enterprise M1034}; \texttt{ICS M0934})
\begin{itemize}
\item Blocks users or groups from installing or using unapproved hardware, including removable devices.
\end{itemize}
\item Audit (\texttt{Enterprise M1047}; \texttt{ICS M0947})
\begin{itemize}
\item Records and reviews activity, permissions, configurations, and integrity indicators to identify weaknesses or anomalies.
\end{itemize}
\item Safety Instrumented Systems (\texttt{ICS M0812})
\begin{itemize}
\item Uses independent safety systems to detect hazardous conditions and automatically drive equipment to a safer state.
\end{itemize}
\item Mitigation Limited or Not Effective (\texttt{ICS M0816})
\begin{itemize}
\item Marks techniques where preventive controls may be limited because the attack abuses legitimate system features.
\end{itemize}
\end{itemize}
\end{itemize}

\begin{itemize}
\item \textbf{Supporting mitigations:}
\begin{itemize}
\item Code Signing (\texttt{Enterprise M1045}; \texttt{ICS M0945})
\begin{itemize}
\item Verifies digital signatures so only trusted, untampered code is allowed to execute.
\end{itemize}
\item Boot Integrity (\texttt{Enterprise M1046}; \texttt{ICS M0946})
\begin{itemize}
\item Verifies the boot process, operating system, and loading mechanisms to prevent tampered components from starting trusted systems.
\end{itemize}
\item Update Software (\texttt{Enterprise M1051}; \texttt{ICS M0951})
\begin{itemize}
\item Applies vendor patches and upgrades to reduce exploitation of known software or firmware vulnerabilities.
\end{itemize}
\item Software Process and Device Authentication (\texttt{ICS M0813})
\begin{itemize}
\item Requires devices and software processes to authenticate, especially for remote connections and API access.
\end{itemize}
\item Operational Information Confidentiality (\texttt{ICS M0809})
\begin{itemize}
\item Protects operational-process, facility, configuration, program, or database information that could reveal sensitive know-how or IP.
\end{itemize}
\item Data Loss Prevention (\texttt{Enterprise M1057}; \texttt{ICS M0803})
\begin{itemize}
\item Identifies, monitors, and controls sensitive-data movement to reduce unauthorized access, transmission, or exfiltration.
\end{itemize}
\item Data Backup (\texttt{ICS M0953}; \texttt{Enterprise M1053})
\begin{itemize}
\item Creates and protects backups of critical data and configurations so systems can recover after compromise or disruption.
\end{itemize}
\end{itemize}
\end{itemize}

\begin{itemize}
\item \textbf{AM-specific application:} control access to feedstock, build chambers, part-handling areas, embedded electronics, inserts, returned parts, and post-processing stations; verify cyber-capable embedded components before deployment.
\end{itemize}

\begin{itemize}
\item \textbf{Residual coverage note:} NBC contamination, hidden foreign material, and physical embedded-object detection require physical inspection, environmental health and safety controls, and material analysis beyond ATT\&CK.
\end{itemize}


\section{3.9 Context-exploitation and destabilization methods [I, A]}


Methods in this family do not primarily manipulate AM technology itself, but exploit external stress conditions that make AM supply chains easier to compromise.


\subsection{3.9.1 Exploitation of economic and market instability [I, A]}


Exploitation of inflation, shortages, market uncertainty, or misinformation to push inferior or counterfeit materials, destabilise supply decisions, or target strained AM supply-chain components.

Retained from Kumar, but treated as a contextual exploitation method rather than a core AM-native engineering manipulation.

Includes:

\begin{itemize}
\item exploiting inflation-driven shortages
\item exploiting cost pressure to insert inferior materials
\item destabilising supply decisions through market manipulation or misinformation
\end{itemize}

\textbf{MITRE ATT\&CK mitigation mapping (derived):}

\begin{itemize}
\item \textbf{Primary mitigations:}
\begin{itemize}
\item Supply Chain Management (\texttt{ICS M0817})
\begin{itemize}
\item Uses supplier policies, procedures, and integrity testing to ensure devices and components come from trusted sources.
\end{itemize}
\item Threat Intelligence Program (\texttt{Enterprise M1019}; \texttt{ICS M0919})
\begin{itemize}
\item Collects and analyzes threat intelligence to track trends, guide defensive priorities, and reduce risk.
\end{itemize}
\item Out-of-Band Communications Channel (\texttt{ICS M0810}; \texttt{Enterprise M1060})
\begin{itemize}
\item Provides alternative communication paths when primary network communication fails, is disrupted, or cannot be trusted.
\end{itemize}
\item Redundancy of Service (\texttt{ICS M0811})
\begin{itemize}
\item Provides backup or standby devices and services so critical functions can continue during failures or disruption.
\end{itemize}
\item User Training (\texttt{Enterprise M1017}; \texttt{ICS M0917})
\begin{itemize}
\item Trains users to recognize and respond to phishing, social engineering, and other user-interaction threats.
\end{itemize}
\item Mitigation Limited or Not Effective (\texttt{ICS M0816})
\begin{itemize}
\item Marks techniques where preventive controls may be limited because the attack abuses legitimate system features.
\end{itemize}
\end{itemize}
\end{itemize}

\begin{itemize}
\item \textbf{Supporting mitigations:}
\begin{itemize}
\item Audit (\texttt{Enterprise M1047}; \texttt{ICS M0947})
\begin{itemize}
\item Records and reviews activity, permissions, configurations, and integrity indicators to identify weaknesses or anomalies.
\end{itemize}
\item Data Backup (\texttt{ICS M0953}; \texttt{Enterprise M1053})
\begin{itemize}
\item Creates and protects backups of critical data and configurations so systems can recover after compromise or disruption.
\end{itemize}
\item Operational Information Confidentiality (\texttt{ICS M0809})
\begin{itemize}
\item Protects operational-process, facility, configuration, program, or database information that could reveal sensitive know-how or IP.
\end{itemize}
\item Data Loss Prevention (\texttt{Enterprise M1057}; \texttt{ICS M0803})
\begin{itemize}
\item Identifies, monitors, and controls sensitive-data movement to reduce unauthorized access, transmission, or exfiltration.
\end{itemize}
\item Network Segmentation (\texttt{Enterprise M1030}; \texttt{ICS M0930})
\begin{itemize}
\item Divides networks into isolated segments or zones to restrict access, reduce lateral movement, and protect critical systems.
\end{itemize}
\item Limit Access to Resource Over Network (\texttt{Enterprise M1035}; \texttt{ICS M0935})
\begin{itemize}
\item Restricts network access to resources such as file shares, remote systems, and services to legitimate users or systems.
\end{itemize}
\end{itemize}
\end{itemize}

\begin{itemize}
\item \textbf{AM-specific application:} preserve trusted supplier lists, emergency procurement rules, alternative qualified material sources, production decision records, and secure communication paths during market stress or misinformation campaigns.
\end{itemize}

\begin{itemize}
\item \textbf{Residual coverage note:} this is mostly a context-exploitation node, not a direct ATT\&CK-style technique. ATT\&CK can support resilience and intelligence, but it cannot prevent inflation, shortages, or market pressure.
\end{itemize}


\subsection{3.9.2 Exploitation of natural or man-made disruption windows [I, A]}


Exploitation of disasters, recovery periods, labour disruption, geopolitical instability, or other emergency windows when AM facilities, suppliers, or controls are degraded.

Derived from Kumar's attack-goal framing around exploited opportunities.

Includes:

\begin{itemize}
\item exploiting disaster recovery windows
\item exploiting labour strikes or geopolitical instability
\item exploiting emergency conditions when controls are weakened
\end{itemize}

\textbf{MITRE ATT\&CK mitigation mapping (derived):}

\begin{itemize}
\item \textbf{Primary mitigations:}
\begin{itemize}
\item Redundancy of Service (\texttt{ICS M0811})
\begin{itemize}
\item Provides backup or standby devices and services so critical functions can continue during failures or disruption.
\end{itemize}
\item Data Backup (\texttt{ICS M0953}; \texttt{Enterprise M1053})
\begin{itemize}
\item Creates and protects backups of critical data and configurations so systems can recover after compromise or disruption.
\end{itemize}
\item Out-of-Band Communications Channel (\texttt{ICS M0810}; \texttt{Enterprise M1060})
\begin{itemize}
\item Provides alternative communication paths when primary network communication fails, is disrupted, or cannot be trusted.
\end{itemize}
\item Supply Chain Management (\texttt{ICS M0817})
\begin{itemize}
\item Uses supplier policies, procedures, and integrity testing to ensure devices and components come from trusted sources.
\end{itemize}
\item Watchdog Timers (\texttt{ICS M0815})
\begin{itemize}
\item Uses timers to quickly detect when a device or system becomes unresponsive.
\end{itemize}
\item Mechanical Protection Layers (\texttt{ICS M0805})
\begin{itemize}
\item Uses physical or mechanical protection layers, such as interlocks or safety devices, to prevent damage or harm.
\end{itemize}
\item Safety Instrumented Systems (\texttt{ICS M0812})
\begin{itemize}
\item Uses independent safety systems to detect hazardous conditions and automatically drive equipment to a safer state.
\end{itemize}
\item Mitigation Limited or Not Effective (\texttt{ICS M0816})
\begin{itemize}
\item Marks techniques where preventive controls may be limited because the attack abuses legitimate system features.
\end{itemize}
\end{itemize}
\end{itemize}

\begin{itemize}
\item \textbf{Supporting mitigations:}
\begin{itemize}
\item Threat Intelligence Program (\texttt{Enterprise M1019}; \texttt{ICS M0919})
\begin{itemize}
\item Collects and analyzes threat intelligence to track trends, guide defensive priorities, and reduce risk.
\end{itemize}
\item Audit (\texttt{Enterprise M1047}; \texttt{ICS M0947})
\begin{itemize}
\item Records and reviews activity, permissions, configurations, and integrity indicators to identify weaknesses or anomalies.
\end{itemize}
\item Network Segmentation (\texttt{Enterprise M1030}; \texttt{ICS M0930})
\begin{itemize}
\item Divides networks into isolated segments or zones to restrict access, reduce lateral movement, and protect critical systems.
\end{itemize}
\item Filter Network Traffic (\texttt{Enterprise M1037}; \texttt{ICS M0937})
\begin{itemize}
\item Filters ingress, egress, or lateral traffic, including protocol-aware rules, to restrict malicious or unnecessary communication.
\end{itemize}
\item Access Management (\texttt{ICS M0801})
\begin{itemize}
\item Enforces access decisions, often through gateways or authentication services, before users can reach systems or devices.
\end{itemize}
\item Authorization Enforcement (\texttt{ICS M0800})
\begin{itemize}
\item Restricts read, manipulation, or execution privileges to authenticated users who require access under approved policies.
\end{itemize}
\item User Training (\texttt{Enterprise M1017}; \texttt{ICS M0917})
\begin{itemize}
\item Trains users to recognize and respond to phishing, social engineering, and other user-interaction threats.
\end{itemize}
\item Update Software (\texttt{ICS M0951}; \texttt{Enterprise M1051})
\begin{itemize}
\item Applies vendor patches and upgrades to reduce exploitation of known software or firmware vulnerabilities.
\end{itemize}
\end{itemize}
\end{itemize}

\begin{itemize}
\item \textbf{AM-specific application:} keep emergency production, recovery builds, supplier substitution, remote support, backup restoration, and degraded-mode operation from bypassing normal security, approval, and quality controls.
\end{itemize}

\begin{itemize}
\item \textbf{Residual coverage note:} natural disasters, strikes, war, and emergency conditions cannot be mitigated away by ATT\&CK controls. The relevant contribution is resilience, fallback communication, recovery integrity, and disciplined exception handling.
\end{itemize}


\chapter{4. Cross-cutting target-to-method relationships}


The merged taxonomy should be used with the following logic:

\begin{itemize}
\item \textbf{Goals} explain the attacker's end state
\item \textbf{Targets} identify where in the AM ecosystem the attack lands
\item \textbf{Methods} identify how the attack is carried out
\end{itemize}

A single goal can be pursued through multiple targets and methods.
A single target can support multiple goals.
A single method can impact multiple targets.

Examples:

\begin{itemize}
\item \textbf{IP compromise} may target:
\begin{itemize}
\item part specification
\item CAD model integrity
\item manufacturing process specification
\item post-processing specification
\item tool-path integrity
\item machine-resident process knowledge
\end{itemize}
\end{itemize}

\begin{itemize}
\item \textbf{Sabotage} may target:
\begin{itemize}
\item final manufactured part
\item manufacturing process execution
\item post-processing execution
\item feedstock
\item AM equipment
\item production environment
\end{itemize}
\end{itemize}

\begin{itemize}
\item \textbf{Availability disruption} may target:
\begin{itemize}
\item AM equipment
\item logistics chain
\item third-party software/firmware services
\item power supply stability
\item production scheduling/throughput
\end{itemize}
\end{itemize}


\chapter{5. CRA Annex I Crosswalk}


This section links the full canonical AM threat taxonomy to the Cyber Resilience Act (CRA), Regulation (EU) 2024/2847, using the attached CRA text as the source. The crosswalk is intentionally requirements-first: each taxonomy node is linked primarily to \textbf{Annex I} and then to the relevant CRA articles that make those requirements legally operative.

The crosswalk covers every canonical node in sections \textbf{1. Attack goals}, \textbf{2. Attack targets}, and \textbf{3. Attack methods and mitigation mappings}. The cross-cutting relationships in section 4 are not mapped separately because they are combinations of those already mapped nodes.


\section{5.1 CRA scope and reference key}


The CRA mapping should always be read through the following scope gate:

\begin{itemize}
\item The CRA applies to products with digital elements made available on the market where the intended purpose or reasonably foreseeable use includes a direct or indirect logical or physical data connection to a device or network, as set out in Article 2(1).
\item Article 3(1) defines a product with digital elements as a software or hardware product and its remote data processing solutions, including software or hardware components placed on the market separately.
\item Article 6 links market availability to Annex I: products must meet Annex I Part I, and manufacturer processes must meet Annex I Part II.
\item Article 13 requires manufacturers to perform cybersecurity risk assessment, account for the results across planning, design, development, production, delivery, and maintenance, and document that assessment in the technical documentation.
\item Threats that concern purely physical materials, business authorization, export control, counterfeiting as a legal act, or general logistics are not directly solved by the CRA unless they are enabled by or affect a product with digital elements.
\end{itemize}

Common CRA anchors used below:

\begingroup
\scriptsize
\setlength{\tabcolsep}{4pt}
\rowcolors{2}{ReportTableStripe}{white}
\begin{longtable}{@{}>{\RaggedRight\arraybackslash}p{0.27\linewidth}>{\RaggedRight\arraybackslash}p{0.66\linewidth}@{}}
\toprule
\textbf{CRA reference} & \textbf{Relevance to this taxonomy} \\
\midrule
\endfirsthead
\toprule
\textbf{CRA reference} & \textbf{Relevance to this taxonomy} \\
\midrule
\endhead
Article 2(1) & Scope gate for products with digital elements with direct or indirect logical or physical data connection. \\
Article 2(6) & Excludes identical spare parts made available to replace identical components manufactured to the same specifications. \\
Article 3(1) & Defines product with digital elements, including software, hardware, remote data processing solutions, and separately placed software or hardware components. \\
Article 3(6) & Defines component as software or hardware intended for integration into an electronic information system. \\
Article 3(23)-(25) & Defines intended purpose, reasonably foreseeable use, and reasonably foreseeable misuse. \\
Article 3(37)-(42) & Defines cybersecurity risk, significant cybersecurity risk, SBOM, vulnerability, exploitable vulnerability, and actively exploited vulnerability. \\
Article 6 & Requires products with digital elements to meet Annex I Part I and manufacturer processes to meet Annex I Part II. \\
Article 13(1)-(4) & Requires secure design, development, production, cybersecurity risk assessment, and technical-documentation inclusion. \\
Article 13(5)-(6) & Requires due diligence for third-party components and reporting/remediation of vulnerabilities in integrated components. \\
Article 13(7)-(11) & Requires documentation of cybersecurity aspects, effective vulnerability handling during the support period, and security-update availability. \\
Article 13(12)-(14) & Links technical documentation, conformity assessment, EU declaration of conformity, CE marking, and ongoing conformity of series production. \\
Article 13(15)-(20) & Covers identification, manufacturer contact details, single point of contact, user information, support-period visibility, and declaration of conformity. \\
Article 13(21)-(22) & Requires corrective measures and cooperation with market surveillance authorities where conformity problems or cybersecurity risks exist. \\
Article 14 & Requires reporting of actively exploited vulnerabilities and severe incidents affecting product security. \\
Article 31 and Annex VII & Define technical-documentation content, including architecture, vulnerability-handling processes, production and monitoring processes, risk assessment, tests, standards, and SBOM where applicable. \\
Article 32 and Annex VIII & Define conformity-assessment procedures against Annex I. \\
Annex II & Defines required user information and instructions, including intended purpose, security environment, foreseeable cybersecurity risks, support, updates, secure use, and decommissioning. \\
\bottomrule
\end{longtable}
\endgroup

Annex I Part I product-property references:

\begingroup
\scriptsize
\setlength{\tabcolsep}{4pt}
\rowcolors{2}{ReportTableStripe}{white}
\begin{longtable}{@{}>{\RaggedRight\arraybackslash}p{0.27\linewidth}>{\RaggedRight\arraybackslash}p{0.66\linewidth}@{}}
\toprule
\textbf{Annex I reference} & \textbf{Requirement theme} \\
\midrule
\endfirsthead
\toprule
\textbf{Annex I reference} & \textbf{Requirement theme} \\
\midrule
\endhead
Annex I.I(1) & Products must be designed, developed, and produced to ensure appropriate cybersecurity based on risks. \\
Annex I.I(2)(a) & Products shall be made available without known exploitable vulnerabilities. \\
Annex I.I(2)(b) & Products shall have secure-by-default configuration, including reset to original state where relevant. \\
Annex I.I(2)(c) & Vulnerabilities must be addressable through security updates, including automatic security updates where applicable. \\
Annex I.I(2)(d) & Products must protect against unauthorised access using appropriate control mechanisms and report possible unauthorised access. \\
Annex I.I(2)(e) & Products must protect confidentiality of stored, transmitted, or otherwise processed data. \\
Annex I.I(2)(f) & Products must protect integrity of stored, transmitted, or otherwise processed data, commands, programs, and configuration against unauthorised manipulation or modification, and report corruptions. \\
Annex I.I(2)(g) & Products shall process only data that is adequate, relevant, and limited to the intended purpose. \\
Annex I.I(2)(h) & Products must protect availability of essential and basic functions, including resilience and mitigation against denial-of-service attacks. \\
Annex I.I(2)(i) & Products must minimise negative impact on the availability of services provided by other devices or networks. \\
Annex I.I(2)(j) & Products must limit attack surfaces, including external interfaces. \\
Annex I.I(2)(k) & Products must reduce incident impact using exploitation-mitigation mechanisms and techniques. \\
Annex I.I(2)(l) & Products must record and monitor relevant internal activity, including access to or modification of data, services, or functions. \\
Annex I.I(2)(m) & Products must allow users to securely remove data and settings and securely transfer data where applicable. \\
\bottomrule
\end{longtable}
\endgroup

Annex I Part II vulnerability-handling references:

\begingroup
\scriptsize
\setlength{\tabcolsep}{4pt}
\rowcolors{2}{ReportTableStripe}{white}
\begin{longtable}{@{}>{\RaggedRight\arraybackslash}p{0.27\linewidth}>{\RaggedRight\arraybackslash}p{0.66\linewidth}@{}}
\toprule
\textbf{Annex I reference} & \textbf{Requirement theme} \\
\midrule
\endfirsthead
\toprule
\textbf{Annex I reference} & \textbf{Requirement theme} \\
\midrule
\endhead
Annex I.II(1) & Identify and document vulnerabilities and components, including an SBOM covering at least top-level dependencies. \\
Annex I.II(2) & Address and remediate vulnerabilities without delay, including through security updates. \\
Annex I.II(3) & Apply effective and regular security tests and reviews. \\
Annex I.II(4) & Publicly disclose information about fixed vulnerabilities after a security update is available, subject to justified delay. \\
Annex I.II(5) & Put in place and enforce a coordinated vulnerability disclosure policy. \\
Annex I.II(6) & Facilitate vulnerability reporting, including by providing a contact address. \\
Annex I.II(7) & Securely distribute updates so vulnerabilities are fixed or mitigated in a timely manner. \\
Annex I.II(8) & Disseminate security updates without delay and free of charge, with advisory messages. \\
\bottomrule
\end{longtable}
\endgroup


\section{5.2 CRA mapping of attack goals}


\begin{landscape}
\begingroup
\scriptsize
\setlength{\tabcolsep}{4pt}
\rowcolors{2}{ReportTableStripe}{white}
\begin{longtable}{@{}>{\RaggedRight\arraybackslash}p{0.15\linewidth}>{\RaggedRight\arraybackslash}p{0.22\linewidth}>{\RaggedRight\arraybackslash}p{0.19\linewidth}>{\RaggedRight\arraybackslash}p{0.12\linewidth}>{\RaggedRight\arraybackslash}p{0.22\linewidth}@{}}
\toprule
\textbf{Taxonomy node} & \textbf{Main CRA requirement link} & \textbf{Legal anchor} & \textbf{Strength} & \textbf{CRA interpretation} \\
\midrule
\endfirsthead
\toprule
\textbf{Taxonomy node} & \textbf{Main CRA requirement link} & \textbf{Legal anchor} & \textbf{Strength} & \textbf{CRA interpretation} \\
\midrule
\endhead
1.1 Intellectual property compromise / technical data theft \texttt{[C]} & Annex I.I(2)(d), (e), (j), (l), (m) & Arts. 2(1), 3(1), 6, 13(1)-(4), 13(17)-(18), 31; Annex II; Annex VII & Direct for digital AM data; partial for IP law & The CRA supports confidentiality, access control, attack-surface limitation, monitoring, and secure data removal for AM data handled by a product with digital elements. It does not itself create an IP-rights enforcement regime. \\
1.2 Unauthorised production / counterfeit / restricted-part production \texttt{[I]} & Annex I.I(2)(b), (d), (f), (l) & Arts. 2(1), 3(1), 6, 13(2)-(4), 31; Annex VII & Partial/gap & The CRA can support access control, configuration integrity, command integrity, and monitoring where unauthorised production is enabled through AM software, firmware, equipment, or remote processing. It does not directly regulate licensing, counterfeiting, export control, grey-market production, or internal production authorization. \\
1.3 Data breach and data manipulation \texttt{[C, I]} & Annex I.I(2)(d), (e), (f), (j), (k), (l), (m); Annex I.II(1)-(8) where vulnerabilities are involved & Arts. 6, 13(1)-(8), 14, 31, 32; Annex VII & Direct & This is the strongest goal-level CRA fit: confidentiality, integrity, access control, attack-surface limitation, monitoring, vulnerability handling, secure updates, and reporting all map directly. \\
1.4 Sabotage / integrity degradation \texttt{[I, A]} & Annex I.I(2)(a), (d), (f), (h), (i), (j), (k), (l); Annex I.II(2), (3), (7), (8) & Arts. 6, 13(1)-(8), 13(21)-(22), 14, 31; Annex VII & Direct where cyber-enabled; partial for purely physical sabotage & The CRA is relevant when sabotage manipulates digital data, commands, configurations, software, firmware, connected equipment, updates, or product behavior. Purely physical sabotage is outside the direct CRA cybersecurity requirements. \\
1.5 Availability disruption / denial of service \texttt{[A]} & Annex I.I(2)(h), (i), (j), (k), (l); Annex I.II(2), (7), (8) & Arts. 6, 13(2)-(8), 14, 31; Annex VII & Direct for products with digital elements & The CRA directly addresses availability of essential and basic functions, DoS resilience, negative impact on other devices or networks, incident impact reduction, monitoring, and timely remediation. \\
\bottomrule
\end{longtable}
\endgroup
\end{landscape}


\section{5.3 CRA mapping of attack targets}


\begin{landscape}
\begingroup
\scriptsize
\setlength{\tabcolsep}{4pt}
\rowcolors{2}{ReportTableStripe}{white}
\begin{longtable}{@{}>{\RaggedRight\arraybackslash}p{0.15\linewidth}>{\RaggedRight\arraybackslash}p{0.22\linewidth}>{\RaggedRight\arraybackslash}p{0.19\linewidth}>{\RaggedRight\arraybackslash}p{0.12\linewidth}>{\RaggedRight\arraybackslash}p{0.22\linewidth}@{}}
\toprule
\textbf{Taxonomy node} & \textbf{Main CRA requirement link} & \textbf{Legal anchor} & \textbf{Strength} & \textbf{CRA interpretation} \\
\midrule
\endfirsthead
\toprule
\textbf{Taxonomy node} & \textbf{Main CRA requirement link} & \textbf{Legal anchor} & \textbf{Strength} & \textbf{CRA interpretation} \\
\midrule
\endhead
2.1 Design and specification assets \texttt{[C, I]} & Annex I.I(2)(d), (e), (f), (l), (m) & Arts. 6, 13(1)-(4), 31; Annex VII & Direct where digitally stored, processed, or transmitted & CRA requirements apply to confidentiality and integrity of design data handled by products with digital elements, including access control and monitoring of access or modification. \\
2.1.1 Part specification \texttt{[C, I]} & Annex I.I(2)(e), (f), (l); Annex VII point 3 & Arts. 13(2)-(4), 31; Annex VII & Conditional/direct & Digital part specifications used by AM software, firmware, or connected equipment map directly. Pure paper or manual specifications are outside direct CRA scope. \\
2.1.2 CAD model integrity \texttt{[C, I]} & Annex I.I(2)(d), (e), (f), (l), (m) & Arts. 6, 13(1)-(4), 31; Annex VII & Direct & CAD files are covered when stored, transmitted, or processed by AM products with digital elements. Integrity protection and monitoring are especially relevant. \\
2.1.3 STL / AMF / 3MF / intermediate model integrity \texttt{[C, I]} & Annex I.I(2)(d), (e), (f), (l) & Arts. 6, 13(1)-(4), 31; Annex VII & Direct & Intermediate build representations are data assets whose confidentiality and integrity map directly to Annex I.I(2)(e) and (f) when processed by AM digital products. \\
2.1.4 Tool-path / slicing output / machine-instruction integrity \texttt{[C, I]} & Annex I.I(2)(d), (f), (h), (l) & Arts. 6, 13(1)-(4), 31; Annex VII & Direct & Tool paths and machine instructions are commands or processed data, so Annex I.I(2)(f) is the central requirement. Availability and monitoring also matter where manipulated instructions affect essential functions. \\
2.2 Process knowledge and production recipe assets \texttt{[C, I, A]} & Annex I.I(2)(d), (e), (f), (h), (l); Annex VII points 2 and 3 & Arts. 13(2)-(4), 31; Annex VII & Direct where digital & Digital process recipes, settings, and production knowledge handled by connected AM systems map to confidentiality, integrity, availability, monitoring, and technical-documentation requirements. \\
2.2.1 Manufacturing process specification \texttt{[C, I]} & Annex I.I(2)(e), (f), (l); Annex VII point 3 & Arts. 13(2)-(4), 31; Annex VII & Direct where digitally represented & CRA relevance is strongest when process specifications are software settings, machine profiles, configuration files, or remote-processing inputs. \\
2.2.2 Manufacturing process execution \texttt{[I, A]} & Annex I.I(2)(d), (f), (h), (j), (k), (l) & Arts. 6, 13(1)-(4), 14, 31 & Direct for digitally controlled execution & Live execution maps to command/configuration integrity, access control, attack-surface reduction, resilience, exploitation mitigation, and monitoring. \\
2.2.3 Post-processing specification and execution \texttt{[C, I, A]} & Annex I.I(2)(e), (f), (h), (l); Annex VII points 2 and 3 & Arts. 13(2)-(4), 31; Annex VII & Conditional & CRA coverage is direct if post-processing is digitally specified, monitored, or controlled. Manual post-processing manipulation is only indirectly relevant through risk assessment if it affects a product with digital elements. \\
2.2.4 Sensor information and quality-control data \texttt{[C, I, A]} & Annex I.I(2)(e), (f), (h), (l); Annex VII points 2, 3, and 6 & Arts. 13(2)-(4), 31; Annex VII & Direct & Sensor and quality-control data are processed data. Integrity, confidentiality, availability, monitoring, production validation, and test reporting are the core CRA links. \\
2.2.5 Indirect manufacturing dependencies \texttt{[C, I, A]} & Annex I.I(2)(h), (i), (j), (k); Annex I.II(1)-(3), (7), (8) & Arts. 2(1), 13(2)-(8), 31; Annex VII & Conditional/direct & Article 2(1) explicitly includes indirect logical or physical connections. Dependencies such as remote services, components, update channels, and connected systems are CRA-relevant when necessary to product function or cybersecurity. \\
2.3 Supply-chain assets and dependencies \texttt{[C, I, A]} & Annex I.II(1)-(8); Annex I.I(2)(a), (c), (d), (f), (h), (l) & Arts. 13(5)-(8), 14, 31; Annex VII & Direct for digital components; partial for physical supply chain & The CRA is strong on third-party software/hardware components, SBOM, vulnerabilities, secure updates, and component due diligence. It is weaker for raw material and logistics issues. \\
2.3.1 Feedstock / source materials \texttt{[I, A]} & Annex I.I(1), (2)(f), (h), (l) only if digitally controlled or monitored & Arts. 13(2)-(4), 31; Annex II item 5; Annex VII & Partial/gap & The CRA does not directly regulate material quality or contamination. It becomes relevant if feedstock selection, validation, sensing, or release is controlled by a product with digital elements. \\
2.3.2 Replacement parts and upgrades \texttt{[I, A]} & Annex I.I(2)(a), (c), (f), (h); Annex I.II(1), (2), (7), (8) & Arts. 2(6), 3(6), 13(5)-(8), 31 & Conditional & Digital upgrades, firmware, software components, and security updates map directly. Identical spare parts made to replace identical components under the same specifications may be excluded by Article 2(6). \\
2.3.3 AM equipment delivered through the supply chain \texttt{[C, I, A]} & Annex I.I(1), (2)(a), (b), (c), (d), (f), (h), (j), (k), (l); Annex I.II(1)-(8) & Arts. 2(1), 3(1), 6, 13, 31, 32 & Direct if connected AM equipment is placed on the market & Connected AM equipment is a strong CRA target if it is a product with digital elements. Secure defaults, no known exploitable vulnerabilities, access control, integrity, availability, monitoring, updates, and conformity assessment all apply. \\
2.3.4 Third-party software, firmware, and service dependencies \texttt{[C, I, A]} & Annex I.II(1)-(8); Annex I.I(2)(a), (c), (d), (f), (j), (k), (l) & Arts. 3(1), 3(6), 13(5)-(8), 14, 31; Annex VII & Direct & This is one of the strongest target-level links. The CRA expressly addresses third-party component due diligence, component vulnerabilities, SBOM, testing, disclosure, secure update distribution, and vulnerability reporting. \\
2.3.5 Logistics and distribution chain \texttt{[C, I, A]} & Annex I.I(1), (2)(f), (h), (i), (l); Annex I.II(7), (8) & Arts. 13(2), 13(8), 14, 31 & Partial/gap & The CRA covers secure delivery and maintenance of products with digital elements only through product risk assessment, update distribution, availability, and incident handling. It does not generally regulate logistics operations. \\
2.3.6 Manufactured objects in transit or return flows \texttt{[I, A]} & Annex I.I(2)(d), (f), (h), (m) if the object is or contains a product with digital elements & Arts. 2(1), 3(1), 13(2)-(4), 31 & Conditional/mostly outside & A manufactured object is CRA-relevant if it is itself a product with digital elements or contains one. A non-digital AM part in transit is mostly outside the CRA. \\
2.4 Physical production system assets \texttt{[C, I, A]} & Annex I.I(2)(d), (f), (h), (i), (j), (k), (l); Annex I.II(1)-(8) & Arts. 2(1), 6, 13, 14, 31 & Conditional/direct & The CRA strongly covers connected production-system assets that are products with digital elements. Pure physical site conditions are only indirectly relevant. \\
2.4.1 Final manufactured part \texttt{[C, I]} & Annex I.I(1), (2)(d), (e), (f), (h), (m) if the part is a product with digital elements & Arts. 2(1), 3(1), 6, 13, 31 & Conditional & A final AM part with software, hardware, connectivity, or remote data processing may be covered. A purely mechanical AM part is mostly outside the CRA. \\
2.4.2 AM equipment \texttt{[C, I, A]} & Annex I.I(1), (2)(a), (b), (c), (d), (e), (f), (h), (j), (k), (l); Annex I.II(1)-(8) & Arts. 2(1), 3(1), 6, 13, 14, 31, 32 & Direct if connected & Connected AM equipment is directly linked to most CRA product-property and vulnerability-handling requirements. \\
2.4.3 AM production environment \texttt{[I, A]} & Annex I.I(1), (2)(h), (i), (j), (k), (l); Annex II items 4, 5, and 8 & Arts. 13(2)-(4), 13(18), 31 & Partial & The CRA does not regulate the entire production environment as a site-security regime. It does require risk assessment that considers conditions of use, operational environment, and assets to be protected. \\
2.4.4 Downstream environment affected by deployed part \texttt{[I, A]} & Annex I.I(1), (2)(h), (i), (k); Annex II items 4, 5, and 8 & Arts. 13(2)-(4), 13(18), 31 & Conditional/partial & The CRA is relevant when product behavior or compromise can affect other devices, networks, or essential functions. It does not comprehensively regulate downstream physical harm outside product cybersecurity. \\
2.5 Rights, governance, and restricted-object targets \texttt{[C, I]} & Annex I.I(2)(d), (e), (f), (l) where enforced through digital controls & Arts. 13(2)-(4), 31; Annex VII & Partial/gap & CRA controls can support governance through access control, integrity, confidentiality, and monitoring. The CRA does not directly enforce IP rights, contract limits, export controls, or internal policy restrictions. \\
2.5.1 Intellectual property rights and legally protected knowledge \texttt{[C]} & Annex I.I(2)(d), (e), (l), (m) & Arts. 6, 13(1)-(4), 31 & Direct for confidentiality; gap for IP enforcement & The CRA can protect digital confidentiality and access to protected knowledge, but it is not an intellectual-property enforcement instrument. \\
2.5.2 Unauthorised overproduction \texttt{[I, A]} & Annex I.I(2)(b), (d), (f), (l) & Arts. 13(2)-(4), 31; Annex VII & Partial/gap & CRA controls may help prevent or detect unauthorised commands, configurations, or access that enable overproduction. It does not regulate production quotas or contractual authorization as such. \\
2.5.3 Export-controlled parts \texttt{[C, I]} & Annex I.I(2)(d), (e), (f), (l) where digital restrictions or protected data are involved & Arts. 13(2)-(4), 31 & Mostly outside CRA & The CRA can protect digital files and access paths, but export-control status is outside the CRA's essential cybersecurity requirements. \\
2.5.4 Internally restricted parts \texttt{[C, I]} & Annex I.I(2)(b), (d), (e), (f), (l) & Arts. 13(2)-(4), 31 & Partial & CRA requirements can support internal restriction mechanisms if implemented through digital access control, secure configuration, integrity protection, and monitoring. Internal policy content itself is outside CRA scope. \\
\bottomrule
\end{longtable}
\endgroup
\end{landscape}


\section{5.4 CRA mapping of attack methods}


\begin{landscape}
\begingroup
\scriptsize
\setlength{\tabcolsep}{4pt}
\rowcolors{2}{ReportTableStripe}{white}
\begin{longtable}{@{}>{\RaggedRight\arraybackslash}p{0.15\linewidth}>{\RaggedRight\arraybackslash}p{0.22\linewidth}>{\RaggedRight\arraybackslash}p{0.19\linewidth}>{\RaggedRight\arraybackslash}p{0.12\linewidth}>{\RaggedRight\arraybackslash}p{0.22\linewidth}@{}}
\toprule
\textbf{Taxonomy node} & \textbf{Main CRA requirement link} & \textbf{Legal anchor} & \textbf{Strength} & \textbf{CRA interpretation} \\
\midrule
\endfirsthead
\toprule
\textbf{Taxonomy node} & \textbf{Main CRA requirement link} & \textbf{Legal anchor} & \textbf{Strength} & \textbf{CRA interpretation} \\
\midrule
\endhead
3.1 Acquisition and exfiltration methods \texttt{[C, I]} & Annex I.I(2)(d), (e), (j), (l), (m); Annex I.II(3) & Arts. 6, 13(1)-(4), 13(17)-(18), 31; Annex II; Annex VII & Direct for digital acquisition/exfiltration & CRA requirements directly address unauthorised access, confidentiality, attack-surface limitation, monitoring, secure data removal, and security testing. \\
3.1.1 Theft / intellectual property theft \texttt{[C]} & Annex I.I(2)(d), (e), (j), (l), (m); Annex I.II(3), (5), (6) & Arts. 6, 13(1)-(8), 14 where actively exploited vulnerabilities or severe incidents are involved, 31 & Direct for digital theft; partial for IP enforcement & Theft of CAD files, process data, credentials, or machine-resident data maps directly to access control, confidentiality, monitoring, vulnerability handling, and reporting. Legal ownership and IP remedies sit outside the CRA. \\
3.1.2 Reverse engineering \texttt{[C, I]} & Annex I.I(2)(e), (g), (j), (m) & Arts. 13(2)-(4), 13(18), 31; Annex II; Annex VII & Partial/gap & The CRA can reduce exposed data, interfaces, and residual information, but it does not prohibit reverse engineering as such. It is strongest where reverse engineering exploits exposed interfaces, stored data, insecure decommissioning, or unnecessary data processing. \\
3.2 Protection-circumvention and unauthorised reproduction methods \texttt{[C, I]} & Annex I.I(2)(b), (d), (f), (j), (k), (l) & Arts. 6, 13(1)-(4), 31; Annex VII & Direct for circumvention of digital protections; partial for reproduction & Bypassing authentication, configuration, or integrity controls is a direct CRA concern. Unauthorised reproduction as a legal or commercial act is only partly covered. \\
3.2.1 Bypassing digital and physical protection / protection removal \texttt{[C, I]} & Annex I.I(2)(b), (d), (f), (j), (k), (l) & Arts. 6, 13(1)-(4), 31 & Direct for digital controls; conditional for physical controls & Digital protection removal maps to secure default configuration, access control, integrity, attack-surface limitation, exploitation mitigation, and monitoring. Physical protection removal is CRA-relevant only where tied to product cybersecurity. \\
3.2.2 Counterfeiting / unlicensed reproduction \texttt{[I]} & Annex I.I(2)(d), (f), (l) where enabled through digital access or command manipulation & Arts. 13(2)-(4), 31 & Partial/gap & The CRA can constrain unauthorised access and manipulation of digital production paths. It does not directly regulate counterfeiting, licensing, grey-market sales, or reproduction rights. \\
3.3 Design and specification tampering methods \texttt{[I]} & Annex I.I(2)(d), (f), (j), (l); Annex VII points 2, 3, and 6 & Arts. 6, 13(1)-(4), 31 & Direct & Annex I.I(2)(f) is the core link because the threat manipulates data, commands, programs, or configuration. Monitoring and technical documentation support traceability. \\
3.3.1 Part specification tampering \texttt{[I]} & Annex I.I(2)(f), (l); Annex VII point 3 & Arts. 13(2)-(4), 31 & Direct where digital & Digital part-specification tampering directly maps to integrity protection and monitoring. Manual specification tampering is outside direct CRA scope unless it affects a product with digital elements. \\
3.3.2 CAD/STL/intermediate-representation tampering \texttt{[I]} & Annex I.I(2)(d), (e), (f), (l) & Arts. 6, 13(1)-(4), 31 & Direct & CAD, STL, AMF, 3MF, and intermediate representations are processed data. Integrity protection against unauthorised modification is the primary CRA link. \\
3.3.3 Tool-path and execution-file tampering \texttt{[I]} & Annex I.I(2)(d), (f), (h), (l) & Arts. 6, 13(1)-(4), 14, 31 & Direct & Tool paths and execution files are commands or machine instructions. The CRA directly requires integrity protection and reporting of corruptions, with availability implications where tampering disrupts essential functions. \\
3.4 Process tampering methods \texttt{[I, A]} & Annex I.I(2)(d), (f), (h), (j), (k), (l); Annex VII points 2 and 3 & Arts. 6, 13(1)-(4), 14, 31 & Direct for digitally controlled processes & Process tampering maps to access control, command/configuration integrity, availability, attack-surface limitation, incident-impact reduction, monitoring, and risk documentation. \\
3.4.1 Compromising AM process parameters \texttt{[I, A]} & Annex I.I(2)(d), (f), (h), (l); Annex VII point 3 & Arts. 13(2)-(4), 31 & Direct where parameters are digital & Digital process parameters are configuration or processed data. Integrity and monitoring are the core CRA requirements. \\
3.4.2 Process-control manipulation \texttt{[I, A]} & Annex I.I(2)(d), (f), (h), (i), (j), (k), (l) & Arts. 6, 13(1)-(4), 14, 31 & Direct & Manipulation of control behavior maps directly to access control, command integrity, availability of essential functions, reduction of impact on other devices or networks, and monitoring. \\
3.4.3 Sensor-information falsification \texttt{[I]} & Annex I.I(2)(e), (f), (h), (l); Annex VII points 2, 3, and 6 & Arts. 13(2)-(4), 31 & Direct & Sensor readings and quality data are processed data. Integrity, confidentiality where relevant, monitoring, production validation, and test reporting are the main CRA links. \\
3.4.4 Timing disruption \texttt{[I, A]} & Annex I.I(2)(f), (h), (i), (l) & Arts. 13(2)-(4), 14, 31 & Conditional/direct & Timing attacks map directly where timing is controlled by software, firmware, scheduling, machine instructions, or remote processing. Pure physical delays are outside direct CRA scope. \\
3.4.5 Power manipulation / power surge \texttt{[I, A]} & Annex I.I(2)(h), (i), (k), (l) where power behavior is digitally mediated or affects product availability & Arts. 13(2)-(4), 14, 31 & Partial/conditional & The CRA can address resilience and monitoring of products with digital elements against availability impacts. It does not directly regulate external electrical infrastructure or physical power sabotage. \\
3.5 Post-processing tampering methods \texttt{[I]} & Annex I.I(2)(d), (f), (h), (l); Annex VII points 2 and 3 & Arts. 13(2)-(4), 31 & Conditional & CRA relevance depends on whether post-processing is specified, controlled, monitored, or validated by a product with digital elements. Manual tampering is mostly outside direct CRA scope. \\
3.5.1 Post-processing procedure tampering \texttt{[I]} & Annex I.I(2)(f), (l); Annex VII point 3 & Arts. 13(2)-(4), 31 & Conditional & Digital procedure records, recipes, work instructions, or equipment configurations map to integrity and monitoring. Purely manual deviations are only indirectly covered through risk assessment if product cybersecurity is implicated. \\
3.5.2 Heat-treatment manipulation \texttt{[I]} & Annex I.I(2)(d), (f), (h), (l) if digitally controlled or monitored & Arts. 13(2)-(4), 31 & Conditional/partial & Digitally controlled heat treatment maps to command/configuration integrity and monitoring. Physical heat-treatment manipulation outside digital control is not directly regulated by the CRA. \\
3.5.3 Finish-machining manipulation \texttt{[I]} & Annex I.I(2)(d), (f), (h), (l) if digitally controlled or monitored & Arts. 13(2)-(4), 31 & Conditional/partial & CNC-like or digitally controlled finish machining can map to CRA integrity and availability requirements. Manual finishing manipulation is mostly outside direct CRA scope. \\
3.6 Supply-chain compromise methods \texttt{[C, I, A]} & Annex I.I(2)(a), (c), (d), (f), (h), (j), (k), (l); Annex I.II(1)-(8) & Arts. 13(5)-(8), 14, 31; Annex VII & Direct for digital/component supply chain; partial for material/logistics supply chain & The CRA is strongest on software, firmware, hardware components, vulnerabilities, SBOM, secure updates, tests, reviews, and component due diligence. \\
3.6.1 Exploitation of supply-chain weaknesses \texttt{[C, I, A]} & Annex I.II(1)-(8); Annex I.I(2)(a), (c), (d), (f), (j), (k), (l) & Arts. 3(6), 13(5)-(8), 14, 31 & Direct for digital components; partial for non-digital suppliers & Exploited weaknesses in components, update channels, dependencies, or remote services map directly. Supplier governance, logistics disruption, and physical material compromise are not fully covered. \\
3.6.2 Feedstock compromise \texttt{[I, A]} & Annex I.I(1), (2)(f), (h), (l) only if feedstock checks or release are digitally controlled & Arts. 13(2)-(4), 31; Annex II item 5 & Partial/gap & The CRA does not directly regulate feedstock composition or contamination. It may apply to digital systems that monitor, authorize, or record material use. \\
3.6.3 Replacement-parts compromise \texttt{[I, A]} & Annex I.I(2)(a), (c), (f), (h); Annex I.II(1), (2), (7), (8) & Arts. 2(6), 3(6), 13(5)-(8), 31 & Conditional & Compromised software, firmware, or digital components map directly. Identical replacement spare parts may be excluded by Article 2(6). Non-digital parts are only indirectly covered. \\
3.6.4 AM equipment compromise before deployment \texttt{[C, I, A]} & Annex I.I(1), (2)(a), (b), (c), (d), (e), (f), (h), (j), (k), (l); Annex I.II(1)-(8) & Arts. 2(1), 3(1), 6, 13, 14, 31, 32 & Direct if equipment is a product with digital elements & Pre-deployment compromise of connected AM equipment is a core CRA concern because it affects secure design, production, known vulnerabilities, secure defaults, updates, access control, integrity, availability, and conformity. \\
3.6.5 Manufactured-object compromise in logistics \texttt{[I, A]} & Annex I.I(2)(d), (f), (h), (m) where the object is a product with digital elements & Arts. 2(1), 3(1), 13(2)-(4), 31 & Conditional/mostly outside & If the manufactured object is a product with digital elements, CRA product-security requirements apply. If it is a non-digital part physically altered in transit, the CRA is mostly not the relevant instrument. \\
3.6.6 Disrupting supply-chain logistics \texttt{[C, I, A]} & Annex I.I(2)(h), (i), (l); Annex I.II(7), (8) where update delivery or remote services are affected & Arts. 13(2), 13(8), 14, 31 & Partial/gap & The CRA can address availability of product functions, secure update distribution, and incidents affecting product security. It does not generally regulate transport, warehousing, or logistics continuity. \\
3.7 Equipment-compromise methods \texttt{[C, I, A]} & Annex I.I(2)(a), (b), (c), (d), (e), (f), (h), (j), (k), (l); Annex I.II(1)-(8) & Arts. 2(1), 3(1), 6, 13, 14, 31, 32 & Direct for connected AM equipment & Connected AM equipment fits CRA product-security logic strongly. Purely physical equipment sabotage is less directly covered. \\
3.7.1 Cyber intrusion on AM equipment \texttt{[C, I, A]} & Annex I.I(2)(a), (b), (d), (e), (f), (h), (j), (k), (l); Annex I.II(1)-(8) & Arts. 6, 13, 14, 31, 32 & Direct & Cyber intrusion maps directly to known-vulnerability handling, secure defaults, access control, confidentiality, integrity, availability, attack-surface reduction, exploitation mitigation, monitoring, updates, and reporting. \\
3.7.2 Physical sabotage on AM equipment \texttt{[I, A]} & Annex I.I(2)(h), (k), (l) where sabotage affects or is detected by the product with digital elements & Arts. 13(2)-(4), 14, 31 & Partial/conditional & The CRA is not a physical-sabotage regulation. It matters where physical sabotage affects product cybersecurity, availability, monitoring, incident impact, or foreseeable use/misuse of the product. \\
3.7.3 Manipulating 3D equipment \texttt{[C, I, A]} & Annex I.I(2)(d), (e), (f), (h), (j), (k), (l) & Arts. 6, 13(1)-(4), 14, 31 & Direct where manipulation is digital & Digital manipulation of AM equipment maps strongly to access control, data/command/configuration integrity, availability, attack-surface limitation, exploitation mitigation, and monitoring. \\
3.7.4 Compromised AM equipment as a method-enabler \texttt{[C, I, A]} & Annex I.I(2)(a), (c), (d), (f), (h), (j), (k), (l); Annex I.II(1)-(8) & Arts. 6, 13(1)-(8), 13(21)-(22), 14, 31, 32 & Direct & Where compromised equipment enables later attacks, the CRA links to vulnerability remediation, security updates, access control, integrity, availability, monitoring, corrective measures, reporting, and conformity assessment. \\
3.8 Part-compromise methods \texttt{[I, A]} & Annex I.I(1), (2)(d), (f), (h), (l) where compromise is cyber-enabled or affects a product with digital elements & Arts. 2(1), 3(1), 13(2)-(4), 31 & Conditional/partial & The CRA is relevant if part compromise results from digital design/process manipulation or if the part itself is a product with digital elements. Pure physical defects or material tampering are mostly outside direct CRA scope. \\
3.8.1 Compromising 3D-printed parts \texttt{[I]} & Annex I.I(2)(f), (h), (l) where caused by digital manipulation; Annex I.I(1) where the part is a product with digital elements & Arts. 2(1), 3(1), 13(2)-(4), 31 & Conditional/partial & The CRA maps strongly to the upstream cyber manipulation that produces a compromised part. It does not directly regulate every physical defect in a non-digital AM part. \\
3.8.2 Contamination and embedded-component insertion \texttt{[I, A]} & Annex I.I(2)(a), (d), (f), (h); Annex I.II(1), (2), (3) where the inserted component is digital & Arts. 3(1), 3(6), 13(5)-(8), 31 & Conditional/gap & Embedded software or hardware components can trigger CRA component and vulnerability obligations. Material contamination without a digital component is mostly outside CRA scope. \\
3.9 Context-exploitation and destabilization methods \texttt{[I, A]} & Annex I.I(1), (2)(h), (i), (k), (l); Annex I.II(2), (7), (8) & Arts. 13(2)-(8), 14, 31; Annex II item 5 & Indirect/partial & The CRA can support resilience, risk assessment, secure updates, and incident handling during destabilizing contexts. It does not regulate the external crisis conditions themselves. \\
3.9.1 Exploitation of economic and market instability \texttt{[I, A]} & Annex I.I(1), (2)(d), (f), (h), (l); Annex I.II(1) where unstable sourcing affects components & Arts. 13(2)-(8), 31; Annex VII & Indirect/gap & Economic instability is not a CRA threat category. CRA relevance appears only where instability leads to component risk, unauthorised access, insecure configuration, integrity failure, or vulnerability-handling failures in products with digital elements. \\
3.9.2 Exploitation of natural or man-made disruption windows \texttt{[I, A]} & Annex I.I(1), (2)(h), (i), (k), (l); Annex I.II(2), (7), (8) & Arts. 13(2)-(8), 14, 31; Annex II item 5 & Indirect/partial & The CRA supports product-level resilience, secure updates, monitoring, and risk documentation under foreseeable conditions. It does not directly govern disaster response, labour disruption, war, or emergency production policy. \\
\bottomrule
\end{longtable}
\endgroup
\end{landscape}


\section{5.5 CRA coverage summary}


\begingroup
\scriptsize
\setlength{\tabcolsep}{4pt}
\rowcolors{2}{ReportTableStripe}{white}
\begin{longtable}{@{}>{\RaggedRight\arraybackslash}p{0.20\linewidth}>{\RaggedRight\arraybackslash}p{0.34\linewidth}>{\RaggedRight\arraybackslash}p{0.36\linewidth}@{}}
\toprule
\textbf{CRA coverage level} & \textbf{Threat-taxonomy areas} & \textbf{Reason} \\
\midrule
\endfirsthead
\toprule
\textbf{CRA coverage level} & \textbf{Threat-taxonomy areas} & \textbf{Reason} \\
\midrule
\endhead
Strong CRA coverage & Data breach/manipulation; CAD/STL/toolpath tampering; process-control manipulation; cyber intrusion on AM equipment; third-party software/firmware/service compromise; insecure updates; known exploitable vulnerabilities; DoS against product functions & These map directly to Annex I.I(2)(a), (c), (d), (e), (f), (h), (j), (k), (l), and Annex I.II(1)-(8). \\
Conditional CRA coverage & AM equipment, post-processing systems, manufactured parts, replacement parts, embedded components, remote services, sensor/QC data, indirect manufacturing dependencies & These are CRA-relevant when they are products with digital elements, components, remote data processing solutions, or connected systems within Article 2(1) and Article 3(1). \\
Partial CRA coverage & Feedstock compromise; logistics disruption; manufactured-object compromise in transit; AM production environment; downstream environment affected by a deployed part & The CRA may cover digital monitoring, availability, instructions, product impact, or risk assessment, but not the entire physical or operational threat. \\
Mostly outside direct CRA coverage & Counterfeiting as a legal/commercial act; unlicensed reproduction as an IP violation; overproduction as a contractual or business-control issue; export-controlled part production; internally restricted part policy; natural disasters; labour strikes; geopolitical instability; purely physical sabotage & These threats can become CRA-relevant only if they are enabled by or affect a product with digital elements. The CRA is not a general AM governance, export-control, IP, logistics, material-quality, or physical-security regulation. \\
\bottomrule
\end{longtable}
\endgroup

The useful research conclusion is therefore narrow but strong: the CRA can be linked rigorously to the \textbf{digital cyber-resilience layer} of AM threats, especially data, commands, software, firmware, connected equipment, vulnerabilities, updates, and component dependencies. It only partially covers the broader AM supply-chain threat surface.


\chapter{6. What was merged, and what was kept distinct}



\section{6.1 Items merged as duplicates or near-duplicates}


These were merged because they describe essentially the same attack idea at different abstraction levels.

\begin{itemize}
\item \textbf{Theft} + \textbf{Intellectual property theft}
\item \textbf{Protection removal} + \textbf{Bypassing digital and physical protection}
\item \textbf{3D part specification manipulation} + \textbf{Tempering with part specification}
\item \textbf{AM manufacturing process manipulation} + \textbf{Compromising AM process parameters}
\item \textbf{Post-processing manipulation} + \textbf{Tempering with post-processing procedures}
\item \textbf{Compromised AM equipment} + \textbf{Cyber intrusion / physical sabotage / manipulating 3D equipment}
\item \textbf{Compromised part} + \textbf{Compromising 3D-printed parts}
\item \textbf{Part specification} + \textbf{3D design and part specification} (while preserving Yampolskiy's internal substructure)
\item \textbf{Post-processing specification} + \textbf{Post-processing stage}
\item \textbf{Environment} + \textbf{AM production environment} (while preserving Yampolskiy's downstream-environment distinction)
\end{itemize}


\section{6.2 Items intentionally kept distinct}


These were \textbf{not} collapsed because the distinction is analytically useful.

\begin{itemize}
\item \textbf{Intellectual property compromise} vs \textbf{Unauthorised production}
\begin{itemize}
\item stealing knowledge is not the same as producing illegal parts
\end{itemize}
\item \textbf{Manufacturing process specification} vs \textbf{Manufacturing process execution}
\begin{itemize}
\item recipe knowledge is different from live process control
\end{itemize}
\item \textbf{CAD integrity} vs \textbf{STL/intermediate integrity} vs \textbf{Tool-path integrity}
\begin{itemize}
\item they are distinct control points in the AM digital chain
\end{itemize}
\item \textbf{Feedstock compromise} vs \textbf{Equipment compromise}
\begin{itemize}
\item both are supply-chain vectors but attack different physical dependencies
\end{itemize}
\item \textbf{Final manufactured part} vs \textbf{Downstream environment affected by the part}
\begin{itemize}
\item the part itself and the harm it later causes are not the same target surface
\end{itemize}
\item \textbf{Unauthorised overproduction} vs \textbf{Counterfeit/restricted-part production}
\begin{itemize}
\item overproduction is an operational/business-control issue that can occur even without a new design being introduced
\end{itemize}
\end{itemize}


\chapter{7. Elements preserved as secondary annotations rather than top-level nodes}


Some source concepts are important, but they fit better as \textbf{annotations} than as top-level merged dimensions.


\section{7.1 Yampolskiy's first-order causal effects}


Yampolskiy includes \textbf{1st order causal effects} under the sabotage target side. In this merged taxonomy, these are treated as \textbf{outcomes/consequences}, not primary targets.

Typical examples:

\begin{itemize}
\item 1st order causal effects affecting the manufacturing site
\begin{itemize}
\item manufacturing-site availability
\end{itemize}
\item 1st order causal effects affecting equipment
\begin{itemize}
\item equipment availability
\item equipment life time
\end{itemize}
\item 1st order causal effects affecting manufactured-part outcomes
\begin{itemize}
\item average material use
\item average time per part
\item average part acceptance
\item breakage conditions
\end{itemize}
\item equipment wear or damage
\item failed quality acceptance
\item damage to the system using the sabotaged part
\item downstream contamination or physical harm
\end{itemize}


\section{7.2 Kumar's exploited opportunities}


Kumar distinguishes between \textbf{attacker goals} and \textbf{exploited opportunities}. In this merged taxonomy, exploited opportunities are preserved as \textbf{contextual exploit conditions}, not primary goals.

Examples:

\begin{itemize}
\item attacker goals
\item cyber goals
\item physical goals
\item natural disruptions
\item man-made disruptions / man-made events
\item emergency or instability windows
\item economic/market instability
\item natural disasters
\item power outages
\item material contamination
\item temperature and humidity fluctuations
\item labour strikes
\item geopolitical instability
\end{itemize}

This keeps the core taxonomy cleaner while preserving the operational insight.


\chapter{8. Recommended short-form version of the merged taxonomy}


For practical use, the taxonomy can be remembered as:


\section{Goals}


\begin{itemize}
\item IP compromise / technical data theft \texttt{[C]}
\item unauthorised production \texttt{[I]}
\item data breach/manipulation \texttt{[C, I]}
\item sabotage \texttt{[I, A]}
\item availability disruption \texttt{[A]}
\end{itemize}


\section{Targets}


\begin{itemize}
\item design/specification assets \texttt{[C, I]}
\item process and recipe assets \texttt{[C, I, A]}
\item supply-chain assets and dependencies \texttt{[C, I, A]}
\item physical production system assets \texttt{[C, I, A]}
\item rights/restricted-object targets \texttt{[C, I]}
\end{itemize}


\section{Methods}


\begin{itemize}
\item acquisition/exfiltration \texttt{[C, I]}
\item protection bypass / unauthorised reproduction \texttt{[C, I]}
\item design/specification tampering \texttt{[I]}
\item process tampering \texttt{[I, A]}
\item post-processing tampering \texttt{[I]}
\item supply-chain compromise \texttt{[C, I, A]}
\item equipment compromise \texttt{[C, I, A]}
\item part compromise \texttt{[I, A]}
\item contextual disruption exploitation \texttt{[I, A]}
\end{itemize}


\section{Mitigation mapping}


Each concrete attack method carries a derived MITRE ATT\&CK mitigation block:

\begin{itemize}
\item \textbf{Primary mitigations:} the ATT\&CK mitigation categories most directly relevant to constraining the method
\item \textbf{Supporting mitigations:} hardening, detection, recovery, governance, or resilience controls that reduce prerequisites or impact
\item \textbf{AM-specific application:} how the mitigation categories apply to AM assets, workflows, equipment, or supply-chain stages
\item \textbf{Residual coverage note:} where ATT\&CK is incomplete or only indirectly useful, especially for physical, material, logistics, legal, or emergency-context risks
\end{itemize}

The mitigation mapping is a CRA-oriented traceability layer, not a claim of product compliance. It links AM attack methods to ATT\&CK mitigation families and then to CRA control areas such as \texttt{Annex I.I(2)(d)} unauthorised-access protection, \texttt{Annex I.I(2)(e)} confidentiality, \texttt{Annex I.I(2)(f)} integrity of data, commands, programs, and configuration, \texttt{Annex I.I(2)(j)} attack-surface limitation, \texttt{Annex I.I(2)(k)} exploitation mitigation, \texttt{Annex I.I(2)(l)} monitoring, and \texttt{Annex I.II(1)-(8)} vulnerability handling.

For the AM digital thread, this means the mapping is strongest for slicer software, printer firmware, machine-control interfaces, build-job transfer, cloud design platforms, update channels, design repositories, sensor streams, and quality records. It is weaker for purely physical, material, contractual, export-control, or logistics issues, which still require AM-specific assurance such as process qualification, geometry comparison, material testing, inspection, supplier audits, and chain-of-custody controls.

The short version is: \textbf{goals explain why, targets explain where, methods explain how, the MITRE mapping explains what defensive control families are relevant, and the CRA crosswalk explains which regulatory requirement those control families support.}


\section{CRA mapping}


The CRA crosswalk should be read as: \textbf{taxonomy node -> Annex I cybersecurity requirement -> legal article anchor -> coverage strength}.

\begin{itemize}
\item Annex I is the primary mapping target.
\item Articles 6 and 13 are the most common legal anchors.
\item Article 14 is relevant when the threat involves an actively exploited vulnerability or severe incident affecting product security.
\item Article 31 and Annex VII are relevant when the threat needs to be reflected in technical documentation, risk assessment, tests, architecture descriptions, vulnerability-handling processes, or SBOM-related evidence.
\item Threats involving materials, logistics, counterfeiting, export control, overproduction, or physical sabotage are only partly covered unless they are enabled by or affect a product with digital elements.
\end{itemize}


\chapter{9. Suggested canonical tree notation}


The tree remains goal-target-method oriented. The MITRE ATT\&CK content is attached at the \textbf{concrete method leaf} level rather than as a fourth top-level taxonomy branch, because mitigations respond to how an attack is carried out. The CRA crosswalk is attached across the \textbf{goal}, \textbf{target}, and \textbf{method} nodes because the CRA is used as a regulatory-requirements traceability layer rather than as a mitigation taxonomy.

\begin{Verbatim}[fontsize=\small,breaklines=true,breakanywhere=true]
Merged AM Threat Taxonomy
+-- 1. Attack goals
|   +-- 1.1 IP compromise / technical data theft [C]
|   +-- 1.2 Unauthorised production / counterfeit / restricted-part production [I]
|   +-- 1.3 Data breach and data manipulation [C, I]
|   +-- 1.4 Sabotage / integrity degradation [I, A]
|   +-- 1.5 Availability disruption / denial of service [A]
|
+-- 2. Attack targets
|   +-- 2.1 Design and specification assets [C, I]
|   |   +-- Part specification [C, I]
|   |   +-- CAD model integrity [C, I]
|   |   +-- STL / intermediate model integrity [C, I]
|   |   +-- Tool-path / machine-instruction integrity [C, I]
|   +-- 2.2 Process knowledge and production recipe assets [C, I, A]
|   |   +-- Manufacturing process specification [C, I]
|   |   +-- Manufacturing process execution [I, A]
|   |   +-- Post-processing specification and execution [C, I, A]
|   |   +-- Sensor / quality-control data [C, I, A]
|   |   +-- Indirect manufacturing dependencies [C, I, A]
|   +-- 2.3 Supply-chain assets and dependencies [C, I, A]
|   |   +-- Feedstock / source materials [I, A]
|   |   +-- Replacement parts and upgrades [I, A]
|   |   +-- AM equipment delivered through supply chain [C, I, A]
|   |   +-- Third-party software / firmware / service dependencies [C, I, A]
|   |   +-- Logistics and distribution chain [C, I, A]
|   |   +-- Manufactured objects in transit or return flows [I, A]
|   +-- 2.4 Physical production system assets [C, I, A]
|   |   +-- Final manufactured part [C, I]
|   |   +-- AM equipment [C, I, A]
|   |   +-- AM production environment [I, A]
|   |   +-- Downstream environment affected by deployed part [I, A]
|   +-- 2.5 Rights, governance, and restricted-object targets [C, I]
|       +-- Intellectual property rights and protected knowledge [C]
|       +-- Unauthorised overproduction [I, A]
|       +-- Export-controlled parts [C, I]
|       +-- Internally restricted parts [C, I]
|
+-- 3. Attack methods
    +-- 3.1 Acquisition and exfiltration methods [C, I]
    |   +-- Theft / IP theft [C]
    |   +-- Reverse engineering [C, I]
    +-- 3.2 Protection circumvention and unauthorised reproduction [C, I]
    |   +-- Bypassing protection / protection removal [C, I]
    |   +-- Counterfeiting / unlicensed reproduction [I]
    +-- 3.3 Design and specification tampering [I]
    |   +-- Part specification tampering [I]
    |   +-- CAD/STL tampering [I]
    |   +-- Tool-path tampering [I]
    +-- 3.4 Process tampering [I, A]
    |   +-- Process parameter compromise [I, A]
    |   +-- Process-control manipulation [I, A]
    |   +-- Sensor-information falsification [I]
    |   +-- Timing disruption [I, A]
    |   +-- Power manipulation / surge [I, A]
    +-- 3.5 Post-processing tampering [I]
    |   +-- Post-processing procedure tampering [I]
    |   +-- Heat-treatment manipulation [I]
    |   +-- Finish-machining manipulation [I]
    +-- 3.6 Supply-chain compromise [C, I, A]
    |   +-- Exploiting supply-chain weaknesses [C, I, A]
    |   +-- Feedstock compromise [I, A]
    |   +-- Replacement-parts compromise [I, A]
    |   +-- Pre-deployment equipment compromise [C, I, A]
    |   +-- Manufactured-object compromise in logistics [I, A]
    |   +-- Logistics disruption [C, I, A]
    +-- 3.7 Equipment compromise [C, I, A]
    |   +-- Cyber intrusion on AM equipment [C, I, A]
    |   +-- Physical sabotage on AM equipment [I, A]
    |   +-- Manipulating 3D equipment [C, I, A]
    |   +-- Compromised AM equipment as method-enabler [C, I, A]
    +-- 3.8 Part compromise [I, A]
    |   +-- Compromising 3D-printed parts [I]
    |   +-- Contamination / embedded-component insertion [I, A]
    +-- 3.9 Context exploitation and destabilization [I, A]
        +-- Economic / market instability exploitation [I, A]
        +-- Natural or man-made disruption exploitation [I, A]
\end{Verbatim}

Each concrete attack-method leaf expands with the same mitigation structure:

\begin{Verbatim}[fontsize=\small,breaklines=true,breakanywhere=true]
Concrete attack method
+-- Includes
+-- MITRE ATT&CK mitigation mapping (derived)
    +-- Primary mitigations
    +-- Supporting mitigations
    +-- AM-specific application
    +-- Residual coverage note
\end{Verbatim}

Example:

\begin{Verbatim}[fontsize=\small,breaklines=true,breakanywhere=true]
3.1.1 Theft / IP theft [C]
+-- Includes
|   +-- theft of design files
|   +-- theft of process files
|   +-- insider / malware / network-intrusion vectors
+-- MITRE ATT&CK mitigation mapping (derived)
    +-- Primary mitigations
    |   +-- Data Loss Prevention (Enterprise M1057; ICS M0803)
    |   |   +-- short MITRE-based explanation
    |   +-- Encrypt Sensitive Information (Enterprise M1041; ICS M0941)
    |   |   +-- short MITRE-based explanation
    |   +-- Access / authorization / authentication controls
    |       +-- short MITRE-based explanation
    +-- Supporting mitigations
    |   +-- Network Segmentation (Enterprise M1030; ICS M0930)
    |   |   +-- short MITRE-based explanation
    |   +-- Audit (Enterprise M1047; ICS M0947)
    |   |   +-- short MITRE-based explanation
    |   +-- Update / vulnerability-management controls
    |       +-- short MITRE-based explanation
    +-- AM-specific application
    +-- Residual coverage note
\end{Verbatim}


\chapter{10. References used for the merge, mitigation mapping, and CRA crosswalk}


\begin{itemize}
\item Yampolskiy, M., King, W. E., Gatlin, J., Belikovetsky, S., Brown, A., Skjellum, A., \& Elovici, Y. (2018). \emph{Security of additive manufacturing: Attack taxonomy and survey}. Additive Manufacturing, 21, 431-457.
\item Kumar, M., Epiphaniou, G., \& Maple, C. (2025). \emph{Security of cyber-physical Additive Manufacturing supply chain: Survey, attack taxonomy and solutions}. Computers \& Security, 157, 104557.
\item Ruohonen, J., Kang, E.-Y., \& Ramadan, Q. (2025). \emph{An Alignment Between the CRA's Essential Requirements and the ATT\&CK's Mitigations}. In 2025 IEEE 33rd International Requirements Engineering Conference Workshops (REW) (pp. 209-214). IEEE. \url{https://doi.org/10.1109/REW66121.2025.00033}.
\item MITRE. (2026). \emph{MITRE ATT\&CK Enterprise Mitigations}. ATT\&CK v19. Accessed 2026-05-01. \url{https://attack.mitre.org/mitigations/enterprise/}
\item MITRE. (2026). \emph{MITRE ATT\&CK Industrial Control Systems (ICS) Mitigations}. ATT\&CK v19. Accessed 2026-05-01. \url{https://attack.mitre.org/mitigations/ics/}
\item MITRE. (2026). \emph{MITRE ATT\&CK Mobile Mitigations}. ATT\&CK v19. Accessed 2026-05-01. \url{https://attack.mitre.org/mitigations/mobile/}
\item European Parliament and Council. (2024). \emph{Regulation (EU) 2024/2847 of the European Parliament and of the Council of 23 October 2024 on horizontal cybersecurity requirements for products with digital elements and amending Regulations (EU) No 168/2013 and (EU) No 2019/1020 and Directive (EU) 2020/1828 (Cyber Resilience Act)}. Official Journal of the European Union, L series, 20.11.2024. Local source used: \texttt{Official CRA document/EUR-Lex Forordning (EU) 2024:2847 (Cyber Resilience Act).pdf}.
\end{itemize}

Source-use notes:

\begin{itemize}
\item The AM taxonomy structure itself is still derived from Yampolskiy and Kumar.
\item The MITRE ATT\&CK entries are used only for the method-level mitigation crosswalk.
\item Ruohonen et al. is used as supporting literature for the CRA/ATT\&CK alignment rationale, not as the source for AM taxonomy nodes or MITRE mitigation identifiers.
\item The CRA entries are used only for the regulatory-requirements crosswalk and are based on the attached Regulation (EU) 2024/2847 text.
\item Enterprise and Industrial Control Systems (ICS) mitigations are used as the primary catalogs because AM combines conventional IT, engineering workstations, production equipment, industrial networks, and cyber-physical control surfaces.
\item Mobile mitigations are referenced for completeness but are not mapped into every method. They should be applied only where the concrete AM product or workflow includes a mobile app, mobile device, or mobile-controlled production interface.
\item The mitigation mappings are \textbf{derived}, not MITRE-authored AM mappings. They identify relevant mitigation families; they do not prove that a control is sufficient for a specific product or risk scenario.
\end{itemize}


\section{Final note}


This merged taxonomy should be treated as a \textbf{synthesized analytical model}, not as a verbatim reproduction of either source paper.

It preserves:

\begin{itemize}
\item \textbf{Kumar's supply-chain-wide scope and goal-target-method structure}
\item \textbf{Yampolskiy's AM-specific detail on part specification, process manipulation, sabotage, equipment compromise, and environment}
\end{itemize}

The result is a more complete AM threat taxonomy for research, threat modeling, and compliance-oriented security analysis. It contains AM threat goals, targets, methods, and method-level mitigation mappings, not ready-made product-specific scenarios. Product-specific risk assessment items are created later by applying those taxonomy entries to the selected product, version, remote data processing solution scope, environment, and assets.





\appendix
\chapter{Appendix: Source Document Outline}


\begin{itemize}
\item Purpose
\item Source-model summary
\begin{itemize}
\item Yampolskiy et al. (2018)
\item Kumar et al. (2025)
\end{itemize}
\item Merge decision
\item CIA Triad Labeling Convention
\item MITRE ATT\&CK Mitigation Mapping Convention
\item CRA Mapping Convention
\item Canonical merged taxonomy
\item 1. Attack goals
\item 2. Attack targets
\item 3. Attack methods and mitigation mappings
\item 4. Cross-cutting target-to-method relationships
\item 5. CRA Annex I Crosswalk
\begin{itemize}
\item 5.1 CRA scope and reference key
\item 5.2 CRA mapping of attack goals
\item 5.3 CRA mapping of attack targets
\item 5.4 CRA mapping of attack methods
\item 5.5 CRA coverage summary
\end{itemize}
\item 6. What was merged, and what was kept distinct
\begin{itemize}
\item 6.1 Items merged as duplicates or near-duplicates
\item 6.2 Items intentionally kept distinct
\end{itemize}
\item 7. Elements preserved as secondary annotations rather than top-level nodes
\begin{itemize}
\item 7.1 Yampolskiy's first-order causal effects
\item 7.2 Kumar's exploited opportunities
\end{itemize}
\item 8. Recommended short-form version of the merged taxonomy
\begin{itemize}
\item Goals
\item Targets
\item Methods
\item Mitigation mapping
\item CRA mapping
\end{itemize}
\item 9. Suggested canonical tree notation
\item 10. References used for the merge, mitigation mapping, and CRA crosswalk
\item Final note
\end{itemize}

\end{document}
